% Les alcènes et les alcynes — Chimie 1ère D — Cours signé OnBuch+
% Programme officiel MINESEC, Première C, D, E, TI. Compiler avec : tectonic N02-les-alcenes-et-les-alcynes.tex
\newcommand{\DOCMATIERE}{Chimie}
\newcommand{\DOCNIVEAU}{1ère D}
\newcommand{\DOCMODULE}{Module 1 : Chimie organique}
\newcommand{\DOCLECON}{02}
\newcommand{\DOCTITRE}{Les alcènes et les alcynes}
% OnBuch+ — préambule « cours signés » ENRICHI (compilé avec tectonic / XeLaTeX)
% Système visuel « Pop » d'OnBuch+ : fond crème (jamais blanc), bordures encre
% épaisses, ombres « dures » décalées, titres Archivo Black en capitales.
% Version enrichie pour les cours de chimie : chimie (mhchem, chemfig),
% graphes (pgfplots), boîtes pédagogiques supplémentaires (définition,
% propriété, méthode, attention, le savais-tu, exercice résolu, exercice,
% TP, bilan), page de garde refaite et signature OnBuch+ ancrée en bas.
%
% Macros attendues dans main.tex AVANT \input{preamble} :
%   \DOCMATIERE  \DOCNIVEAU  \DOCMODULE  \DOCLECON (numéro)  \DOCTITRE
\documentclass[11pt]{article}

\usepackage{fontspec}
\usepackage[french]{babel}
\usepackage[a4paper,margin=2cm,top=3.4cm,bottom=2.8cm,headheight=1.4cm,footskip=1.3cm]{geometry}
\usepackage{amsmath,amssymb,amsfonts}
\usepackage[table]{xcolor} % [table] : fournit aussi \rowcolors (alternance de lignes)
\usepackage{pagecolor}
\usepackage{tikz}
\usetikzlibrary{arrows.meta,positioning,calc,shapes.geometric,shapes.misc,shapes.arrows,decorations.pathmorphing,decorations.markings,patterns,shadows,mindmap,trees,fit,backgrounds,matrix,intersections}
\usepackage{pgfplots}
\pgfplotsset{compat=1.18}
\usepackage[version=4]{mhchem}
\usepackage{chemfig}
\usepackage{enumitem}
\usepackage{fancyhdr}
% [most] charge toutes les bibliothèques tcolorbox (listings, minted, raster,
% documentation...) et gonfle inutilement la mémoire TeX sur un document long
% (~300 boîtes) : on ne charge que ce qui est réellement utilisé.
\usepackage[skins,breakable]{tcolorbox}
\usepackage{graphicx}
\usepackage{colortbl}
\usepackage{booktabs}
\usepackage{tabularx}
\usepackage{array}
\usepackage{multicol}
\usepackage{siunitx}
% parse-numbers=false : accepte aussi \SI{2,5 \times 10^{-3}}{...}
\sisetup{locale=FR,output-decimal-marker={,},group-minimum-digits=4,parse-numbers=false}
\DeclareSIUnit\ohm{\ensuremath{\symup{\Omega}}} % U+2126 absent de la police
\usepackage{qrcode}
% \degree : commande fréquemment utilisée par les modèles pour un angle ou une
% température en dehors de \SI{}{} (non fournie par les paquets chargés).
\providecommand{\degree}{^{\circ}}
\usepackage{lastpage}
\usepackage{hyperref}
\hypersetup{hidelinks}

% ============================================================
% Palette « Pop » OnBuch+ — mêmes hex que la classe P (plus_ui.dart)
% ============================================================
\definecolor{poporange}{HTML}{F59321}
\definecolor{popdark}{HTML}{E07A0C}
\definecolor{popcream}{HTML}{FFF6E8}
\definecolor{popink}{HTML}{1C1714}
\definecolor{poppurple}{HTML}{7A5AE0}
\definecolor{popgreen}{HTML}{2E9E6A}
\definecolor{poppink}{HTML}{E0457B}
\definecolor{popblue}{HTML}{2F7DE1}
\definecolor{popgold}{HTML}{C98A12}
\definecolor{popmuted}{HTML}{8A7F76}
\definecolor{popline}{HTML}{EADFD2}
\definecolor{popgray}{HTML}{8A7F76}
\colorlet{popgrey}{popgray}
% Alias tolérants (le modèle nomme parfois les couleurs différemment)
\colorlet{popwhite}{white}
\colorlet{popblack}{popink}
\colorlet{popred}{poppink}
\colorlet{popyellow}{popgold}
% Teintes claires pour fonds de tableaux / zones
\colorlet{popblueL}{popblue!10!white}
\colorlet{poporangeL}{poporange!14!white}
\colorlet{popgreenL}{popgreen!12!white}
\colorlet{poppurpleL}{poppurple!12!white}
\colorlet{poppinkL}{poppink!12!white}
\colorlet{popgoldL}{popgold!14!white}

\pagecolor{popcream}

% ============================================================
% Typographie
% ============================================================
\defaultfontfeatures{Ligatures=TeX}
\setmainfont{PlusJakartaSans-Medium.ttf}[Path=../../../fonts/,BoldFont=PlusJakartaSans-Bold.ttf]
% Maths en sans-serif (Fira Math) avec chiffres et lettres droites de Plus
% Jakarta, pour que les formules se fondent dans le texte.
\usepackage[math-style=ISO,bold-style=ISO]{unicode-math}
\setmathfont{FiraMath-Regular.otf}[Path=../../../fonts/]
\setmathfont{PlusJakartaSans-Medium.ttf}[Path=../../../fonts/,range={up/num,up/{Latin,latin},\mathup}]
\usepackage{icomma} % virgule décimale sans espace en mode math
% Caractères absents de la police de texte (sinon carré vide) : rendus par
% leur équivalent mathématique, en texte comme en formule.
\usepackage{newunicodechar}
\newunicodechar{α}{\ensuremath{\alpha}}
\newunicodechar{Α}{\ensuremath{\alpha}}
\newunicodechar{β}{\ensuremath{\beta}}
\newunicodechar{δ}{\ensuremath{\delta}}
\newunicodechar{✓}{\popcheck{popgreen}}
\newfontfamily\popdisplay{ArchivoBlack-Regular.ttf}[Path=../../../fonts/]
\newfontfamily\poplogo{SpaceGrotesk-Bold.ttf}[Path=../../../fonts/]
\newfontfamily\popsemi{PlusJakartaSans-SemiBold.ttf}[Path=../../../fonts/]
\newfontfamily\popxbold{PlusJakartaSans-ExtraBold.ttf}[Path=../../../fonts/]
\newfontfamily\popxboldls{PlusJakartaSans-ExtraBold.ttf}[Path=../../../fonts/,LetterSpace=3.0]

\setlist[enumerate]{leftmargin=1.7em,itemsep=5pt,topsep=4pt}
\setlist[itemize]{leftmargin=1.7em,itemsep=4pt,topsep=4pt,label={\color{poporange}\textbullet}}
\setlength{\parindent}{0pt}
\setlength{\parskip}{5pt}
\renewcommand{\arraystretch}{1.25}

% chemfig : liaisons un peu plus épaisses, encre Pop
\setchemfig{atom sep=2em,bond style={line width=0.9pt,popink},cram width=3pt}

% pgfplots : style Pop par défaut
\pgfplotsset{
  every axis/.append style={axis line style={popink,line width=1pt},tick style={popink},
    grid style={popline,line width=0.5pt},label style={font=\small},tick label style={font=\footnotesize}},
  popcurve/.style={poporange,line width=1.8pt,no markers,smooth},
}

% ============================================================
% Pill / squiggle
% ============================================================
\newcommand{\poppill}[3]{%
  \tikz[baseline=-0.55ex]\node[draw=popink,line width=1.2pt,rounded corners=2.6pt,%
    fill=#2,inner xsep=6.5pt,inner ysep=3pt]{%
    \popxboldls\fontsize{7}{7}\selectfont\color{#3}\MakeUppercase{#1}};%
}
\newcommand{\popsquiggle}[1]{%
  \tikz[baseline=-0.4ex]{
    \draw[#1,line width=2.6pt,line cap=round]
      plot [smooth] coordinates {(0,0.09) (0.34,0.01) (0.68,0.21) (1.02,0.01) (1.36,0.10)};
  }%
}
% Mot-clé surligné (marqueur orange sous le texte)
\newcommand{\cle}[1]{{\popxbold\color{popdark}#1}}

% ============================================================
% En-tête / pied
% ============================================================
\pagestyle{fancy}
\fancyhf{}
\renewcommand{\headrulewidth}{0pt}
\renewcommand{\footrulewidth}{0pt}
\lhead{\raisebox{-0.15ex}{\poplogo\fontsize{13}{13}\selectfont\color{popink}OnBuch\color{poporange}+}}
\rhead{\poppill{\DOCMATIERE\ \textbullet\ \DOCNIVEAU\ \textbullet\ Leçon \DOCLECON}{white}{popink}}
\lfoot{\popxbold\fontsize{7.6}{7.6}\selectfont\color{popmuted}\MakeUppercase{Cours signé OnBuch+}\ \textbullet\ {\popsemi\DOCTITRE}}
\rfoot{\tikz[baseline=-0.6ex]\node[rounded corners=8.5pt,fill=popink,minimum height=17pt,minimum width=17pt,inner xsep=5pt,inner sep=0]{\popxbold\fontsize{8}{8}\selectfont\color{popcream}\thepage\,/\,\pageref*{LastPage}};}

% ============================================================
% Titres
% ============================================================
\newcommand{\chaptitre}[2]{%
  \begin{tcolorbox}[enhanced,colback=poporange,colframe=popink,boxrule=2.6pt,arc=11pt,
      drop shadow=popink,left=15pt,right=15pt,top=12pt,bottom=11pt,before skip=3pt,after skip=18pt]
    \poppill{#2}{popink}{popcream}\par\vspace{9pt}
    {\raggedright\hyphenpenalty=10000\exhyphenpenalty=10000\popdisplay\fontsize{22}{24}\selectfont\color{popcream}\MakeUppercase{#1}\par}
  \end{tcolorbox}
}
% Section numérotée : pastille orange avec le numéro + titre Archivo
\newcounter{popsec}
\newcommand{\coursec}[1]{%
  \stepcounter{popsec}\setcounter{popsub}{0}%
  \par\vspace{16pt}\noindent
  \begin{minipage}{\linewidth}
  \tikz[baseline=-0.7ex]\node[draw=popink,line width=1.6pt,fill=poporange,rounded corners=4pt,minimum size=22pt,inner sep=0]{\popdisplay\fontsize{12}{12}\selectfont\color{popcream}\thepopsec};%
  \hspace{8pt}{\popdisplay\fontsize{15}{17}\selectfont\color{popink}\MakeUppercase{#1}}\par
  \vspace{4pt}\noindent{\color{poporange}\rule{36pt}{3.6pt}}
  \end{minipage}\par\nopagebreak\vspace{8pt}
}
\newcounter{popsub}
\newcommand{\courssub}[1]{%
  \stepcounter{popsub}%
  \par\vspace{9pt}\noindent{\popxbold\fontsize{11.5}{13}\selectfont\color{popink}{\color{poporange}\thepopsec.\thepopsub}\hspace{6pt}#1}\par\nopagebreak\vspace{4pt}
}

% ============================================================
% Boîtes pédagogiques — même recette : fond blanc, bordure épaisse colorée,
% ombre dure encre, badge plein. Titre optionnel [#1].
% ============================================================
\tcbset{popbase/.style={breakable,enhanced,colback=white,boxrule=2.3pt,arc=9pt,
  shadow={3pt}{-3pt}{0pt}{popink},left=14pt,right=14pt,top=11pt,bottom=11pt,before skip=11pt,after skip=15pt,
  fontupper=\popsemi\fontsize{10.6}{14.5}\selectfont\color{popink},
  fontlower=\popsemi\fontsize{10.6}{14.5}\selectfont\color{popink}}}
% Coche dessinée (le glyphe ✓ n'existe pas dans les polices du document)
\newcommand{\popcheck}[1]{\tikz[baseline=-0.5ex]\draw[#1,line width=1.6pt,line cap=round,line join=round](-0.13,0.0)--(-0.04,-0.09)--(0.13,0.1);}
\newcommand{\popboxhead}[3]{\poppill{#1}{#2}{white}\ifx\relax#3\relax\else\hspace{6pt}{\popxbold\fontsize{10.6}{12}\selectfont\color{popink}#3}\fi\par\vspace{7pt}}

\newtcolorbox{aretenir}[1][]{popbase,colframe=popgreen,before upper={\popboxhead{À retenir}{popgreen}{#1}}}
\newtcolorbox{exemplebox}[1][]{popbase,colframe=poporange,before upper={\popboxhead{Exemple}{poporange}{#1}}}
\newtcolorbox{definition}[1][]{popbase,colframe=popblue,before upper={\popboxhead{Définition}{popblue}{#1}}}
\newtcolorbox{propriete}[1][]{popbase,colframe=poppurple,before upper={\popboxhead{Propriété}{poppurple}{#1}}}
\newtcolorbox{methode}[1][]{popbase,colframe=popgold,before upper={\popboxhead{Méthode}{popgold}{#1}}}
\newtcolorbox{attention}[1][]{popbase,colframe=poppink,before upper={\popboxhead{Attention}{poppink}{#1}}}
\newtcolorbox{experience}[1][]{popbase,colframe=popblue,colback=popblueL,before upper={\popboxhead{Expérience}{popblue}{#1}}}
% « Le savais-tu ? » : carte violette pleine (contexte, culture, Cameroun)
\newtcolorbox{savaistu}[1][]{popbase,colback=poppurple,colframe=popink,
  fontupper=\popsemi\fontsize{10.4}{14.2}\selectfont\color{white},
  before upper={\poppill{Le savais-tu\,?}{white}{poppurple}\ifx\relax#1\relax\else\hspace{6pt}{\popxbold\color{white}#1}\fi\par\vspace{7pt}}}
% Exercice résolu : énoncé en haut, \tcblower puis la solution sur fond crème
\newcounter{popexo}\newcounter{popexr}
\newtcolorbox{exoresolu}[1][]{popbase,colframe=poporange,colbacklower=popcream,
  segmentation style={popink,line width=1.2pt,dashed},
  before upper={\stepcounter{popexr}\popboxhead{Exercice résolu \thepopexr}{poporange}{#1}},
  before lower={\poppill{Solution}{popink}{popcream}\par\vspace{6pt}}}
% Exercice (non résolu en place) — la correction est dans la partie Corrigés
\newtcolorbox{exercice}[1][]{popbase,colframe=popink,
  before upper={\stepcounter{popexo}\popboxhead{Exercice \thepopexo}{popink}{#1}}}
% Corrigé
\newtcolorbox{corrige}[1][]{popbase,colframe=popgreen,colback=popgreenL,
  before upper={\popboxhead{Corrigé}{popgreen}{#1}}}
% Objectifs / prérequis / bilan
\newtcolorbox{objectifs}{popbase,colframe=popink,colback=poporangeL,
  before upper={\popboxhead{Objectifs de la leçon}{popink}{}}}
\newtcolorbox{prerequis}{popbase,colframe=popink,colback=popblueL,
  before upper={\popboxhead{Prérequis}{popblue}{}}}
\newtcolorbox{bilan}{popbase,colframe=popink,colback=white,boxrule=2.8pt,
  before upper={\popboxhead{Fiche bilan}{poporange}{L'essentiel à connaître pour l'examen}}}
% Figure encadrée avec légende
\newtcolorbox{popfigure}[1][]{enhanced,breakable=false,colback=white,colframe=popline,boxrule=1.4pt,arc=8pt,
  left=10pt,right=10pt,top=10pt,bottom=8pt,before skip=10pt,after skip=12pt,halign=center,
  lower separated=false,halign lower=center,
  fontlower=\popsemi\fontsize{9}{11.5}\selectfont\color{popmuted},
  #1}
\newcommand{\legende}[1]{\par\vspace{6pt}{\popsemi\fontsize{9}{11.5}\selectfont\color{popmuted}#1}}

% ============================================================
% Signature OnBuch+
% ============================================================
\newcommand{\poplogoinline}[1]{{\poplogo\fontsize{#1}{#1}\selectfont\color{popink}OnBuch\color{poporange}+}}
\newcommand{\signatureonbuch}{%
  \begin{tcolorbox}[enhanced,colback=popink,colframe=popink,boxrule=2.2pt,arc=10pt,
      drop shadow=poporange,left=14pt,right=14pt,top=10pt,bottom=10pt,before skip=0pt,after skip=0pt]
    \tikz[baseline=-0.7ex]\node[circle,fill=popgreen,draw=popcream,line width=1.3pt,minimum size=22pt,inner sep=0]{\popcheck{white}};%
    \hspace{9pt}\begin{minipage}[c]{0.62\linewidth}
      {\popxbold\fontsize{12}{13}\selectfont\color{popcream}Signé\ \textbullet\ {\poplogo OnBuch\color{poporange}+}}\par\vspace{2pt}
      {\raggedright\popsemi\fontsize{8}{10.5}\selectfont\color{popline}\MakeUppercase{Équipe pédagogique OnBuch+}\\\MakeUppercase{Conforme au programme officiel MINESEC}\par}
    \end{minipage}\hfill
    {\poplogo\fontsize{22}{22}\selectfont\color{popcream}OnBuch\color{poporange}+}
  \end{tcolorbox}%
}

% ============================================================
% Page de garde
% #1 titre · #2 matière · #3 niveau · #4 module · #5 n° leçon · #6 durée ·
% #7 description / objectifs (texte ou itemize)
% ============================================================
\newcommand{\pagegarde}[7]{%
  \thispagestyle{empty}
  \newgeometry{a4paper,margin=1.7cm,top=1.6cm,bottom=1.4cm}%
  \begin{tikzpicture}[remember picture,overlay]
    \fill[poporange!18] ([xshift=-3.2cm,yshift=-3.6cm]current page.north east) circle (4.6cm);
    \fill[poppurple!14] ([xshift=2.4cm,yshift=7.6cm]current page.south west) circle (3.8cm);
  \end{tikzpicture}%
  {\poplogo\fontsize{30}{30}\selectfont\color{popink}OnBuch\color{poporange}+}\hfill
  \raisebox{0.6ex}{\poppill{Cours signé \textbullet\ Premium}{popink}{popcream}}\par
  \vspace{1pt}\popsquiggle{poppurple}\par
  \vspace{8pt}
  \begin{tcolorbox}[enhanced,colback=poporange,colframe=popink,boxrule=2.8pt,arc=13pt,
      drop shadow=popink,left=18pt,right=18pt,top=12pt,bottom=13pt,after skip=10pt]
    \poppill{#2 \textbullet\ #3}{popink}{popcream}\hspace{5pt}\poppill{#4}{white}{popink}\par\vspace{8pt}
    {\popdisplay\fontsize{14}{14}\selectfont\color{popink}LEÇON #5}\par\vspace{4pt}
    {\raggedright\hyphenpenalty=10000\exhyphenpenalty=10000\popdisplay\fontsize{25}{27}\selectfont\color{popcream}\MakeUppercase{#1}\par}
    \vspace{8pt}
    {\popxbold\fontsize{9.5}{9.5}\selectfont\color{popink}\MakeUppercase{Durée conseillée : #6}}
  \end{tcolorbox}
  \begin{tcolorbox}[enhanced,colback=white,colframe=popink,boxrule=2.2pt,arc=10pt,
      drop shadow=popink,left=16pt,right=16pt,top=10pt,bottom=10pt,after skip=8pt]
    \poppill{Dans cette leçon}{poppurple}{white}\par\vspace{5pt}
    \popsemi\fontsize{9.3}{12.6}\selectfont\color{popink}#7
  \end{tcolorbox}
  \noindent\begin{minipage}[t]{0.487\linewidth}
    \begin{tcolorbox}[enhanced,colback=popgreen,colframe=popink,boxrule=2.2pt,arc=10pt,
        drop shadow=popink,left=11pt,right=11pt,top=10pt,bottom=10pt,height=3.1cm,valign=center]
      \tcbox[colback=white,colframe=popink,boxrule=1.3pt,arc=5pt,boxsep=0pt,nobeforeafter,
        left=3pt,right=3pt,top=3pt,bottom=3pt,valign=center]{\qrcode[height=1.9cm]{https://play.google.com/store/apps/details?id=com.luvvix.onbuch&hl=fr}}%
      \hspace{8pt}\begin{minipage}[c]{0.5\linewidth}
        {\popdisplay\fontsize{11}{12.5}\selectfont\color{white}\MakeUppercase{Télécharge OnBuch}}\par\vspace{4pt}
        {\raggedright\popsemi\fontsize{8.4}{11}\selectfont\color{white}Gratuit \textbullet\ Google Play\\Tuteur IA, annales, concours, résultats\par}
      \end{minipage}
    \end{tcolorbox}
  \end{minipage}\hfill
  \begin{minipage}[t]{0.487\linewidth}
    \begin{tcolorbox}[enhanced,colback=poppurple,colframe=popink,boxrule=2.2pt,arc=10pt,
        drop shadow=popink,left=11pt,right=11pt,top=10pt,bottom=10pt,height=3.1cm,valign=center]
      \tikz[baseline=-0.7ex]\node[circle,fill=white,draw=popink,line width=1.3pt,minimum size=1.6cm,inner sep=0]{\popdisplay\fontsize{24}{24}\selectfont\color{poppurple}+};%
      \hspace{8pt}\begin{minipage}[c]{0.6\linewidth}
        {\popdisplay\fontsize{11}{12.5}\selectfont\color{white}\MakeUppercase{Passe à OnBuch+}}\par\vspace{4pt}
        {\raggedright\popsemi\fontsize{8.4}{11}\selectfont\color{white}Tous les cours signés, Tuteur IA illimité, fiches PDF — onglet OnBuch+ de l'app\par}
      \end{minipage}
    \end{tcolorbox}
  \end{minipage}
  \vspace{8pt}
  \popcontenu
  \vfill
  \signatureonbuch
  \restoregeometry
  \newpage
}

% Rangée « Ce cours contient » (page de garde)
\newcommand{\poptile}[3]{% #1 couleur #2 gros texte #3 légende
  \begin{tcolorbox}[enhanced,colback=#1,colframe=popink,boxrule=2pt,arc=9pt,shadow={2.5pt}{-2.5pt}{0pt}{popink},
      left=6pt,right=6pt,top=9pt,bottom=9pt,halign=center,valign=center,height=1.75cm,width=\linewidth,nobeforeafter]
    {\popdisplay\fontsize{12.5}{13}\selectfont\color{white}\MakeUppercase{#2}}\par\vspace{3pt}
    {\centering\popsemi\fontsize{7.8}{9.5}\selectfont\color{white}#3\par}
  \end{tcolorbox}}
\newcommand{\popcontenu}{%
  \par\noindent{\popxboldls\fontsize{7.5}{7.5}\selectfont\color{popmuted}CE COURS SIGNÉ CONTIENT}\par\vspace{7pt}
  \noindent\begin{minipage}[t]{0.235\linewidth}\poptile{popblue}{Cours}{complet, illustré, pas à pas}\end{minipage}\hfill
  \begin{minipage}[t]{0.235\linewidth}\poptile{poporange}{Méthodes}{et exercices résolus}\end{minipage}\hfill
  \begin{minipage}[t]{0.235\linewidth}\poptile{poppink}{Exercices}{type examen + corrigés}\end{minipage}\hfill
  \begin{minipage}[t]{0.235\linewidth}\poptile{popgreen}{Fiche bilan}{l'essentiel à retenir}\end{minipage}\par}

% Sommaire
\newtcolorbox{sommaire}{enhanced,colback=white,colframe=popink,boxrule=2.2pt,arc=10pt,
  drop shadow=popink,left=18pt,right=18pt,top=13pt,bottom=13pt,before skip=6pt,after skip=20pt,
  before upper={\poppill{Sommaire}{poppurple}{white}\par\vspace{9pt}\popsemi\fontsize{10.6}{14.5}\selectfont\color{popink}}}

% Signature de fin de cours — poussée tout en bas de la dernière page
\newcommand{\findecours}{%
  \par\vfill\nopagebreak
  \begin{center}\popsquiggle{poporange}\end{center}\vspace{4pt}
  \signatureonbuch
}
% Glyphes absents de Fira Math : redéfinis à partir de glyphes disponibles
\AtBeginDocument{%
  \def\setminus{\mathbin{\backslash}}\let\smallsetminus\setminus
  \def\vdots{\mathord{\vbox{\baselineskip3.2pt\lineskiplimit0pt\kern4pt\hbox{.}\hbox{.}\hbox{.}}}}%
  \def\ddots{\mathinner{\mkern1mu\raise6.5pt\hbox{.}\mkern2mu\raise3.6pt\hbox{.}\mkern2mu\raise0.7pt\hbox{.}\mkern1mu}}%
  \def\triangle{\mathord{\scalebox{0.9}{$\symup{\Delta}$}}}%
  \def\nabla{\mathord{\raisebox{\depth}{\scalebox{1}[-1]{$\symup{\Delta}$}}}}%
  \def\checkmark{\popcheck{popdark}}%
  \def\star{\mathbin{\ast}}\def\bigstar{\mathord{\ast}}%
  \def\diamond{\mathbin{\scalebox{0.6}{$\Diamond$}}}%
  \def\bigcup{\mathop{\mathchoice{\raisebox{-0.3em}{\scalebox{1.7}{$\cup$}}}{\scalebox{1.25}{$\cup$}}{\cup}{\cup}}\displaylimits}%
  \def\bigcap{\mathop{\mathchoice{\raisebox{-0.3em}{\scalebox{1.7}{$\cap$}}}{\scalebox{1.25}{$\cap$}}{\cap}{\cap}}\displaylimits}%
  \def\lesseqqgtr{\mathrel{\lessgtr}}\def\bigoplus{\mathop{\oplus}\displaylimits}%
}
\providecommand{\ssi}{si et seulement si} % abréviation fréquente

%% FIGFIX-BEGIN v3 — correctifs globaux des figures (docs/onbuch-generation/tools/figfix)
\usepackage{adjustbox}\usepackage{etoolbox}
% toute figure TikZ/pgfplots plus large que la ligne est réduite à la largeur disponible
\BeforeBeginEnvironment{tikzpicture}{\begin{adjustbox}{max width=\linewidth}}
\AfterEndEnvironment{tikzpicture}{\end{adjustbox}}
% légendes pgfplots lisibles par-dessus les courbes
\pgfplotsset{every axis/.append style={legend style={fill=white,fill opacity=0.92,text opacity=1,draw=black!15,inner sep=3pt}}}
% étiquettes d'arêtes (syntaxe edge["..."]) avec fond blanc
\tikzset{every edge quotes/.append style={fill=white,inner sep=1.2pt}}
% bulles de cartes mentales : texte à la ligne dans la bulle, bulle un peu plus grande
\tikzset{every concept/.append style={text width=2.0cm,align=center,minimum size=2.3cm,inner sep=2pt}}
%% FIGFIX-END
\begin{document}

\pagegarde{Les alcènes et les alcynes}{Chimie}{1ère D}{Module 1 : Chimie organique}{02}{6 h}{Cette leçon présente les hydrocarbures insaturés : les alcènes (double liaison C=C) et les alcynes (triple liaison C≡C). Elle couvre leur structure spatiale, leur nomenclature, l'isomérie Z/E, leurs réactions caractéristiques (addition, polymérisation, combustion) et la préparation de l'acétylène au laboratoire, gaz utilisé dans le soudage au Cameroun.
\vspace{4pt}
\begin{multicols}{2}\small
\begin{itemize}
\item À la fin de cette leçon, je sais décrire la structure spatiale (géométrie, angles, longueurs de liaisons) des alcènes et des alcynes à partir des exemples de l'éthylène et de l'éthyne.
\item À la fin de cette leçon, je sais représenter géométriquement les alcènes et les alcynes à l'aide de modèles moléculaires ou de schémas.
\item À la fin de cette leçon, je sais écrire les formules générales, développées et semi-développées des alcènes (CnH2n) et des alcynes (CnH2n-2).
\item À la fin de cette leçon, je sais nommer un alcène ou un alcyne en appliquant les règles de nomenclature officielles.
\item À la fin de cette leçon, je sais distinguer et représenter les isomères Z et E d'un alcène.
\item À la fin de cette leçon, je sais écrire et équilibrer les équations des réactions d'addition (hydrogénation, hydratation, HCl, Cl2, Br2) en appliquant la règle de Markovnikov.
\end{itemize}\end{multicols}}

\begin{sommaire}
\begin{multicols}{2}
\begin{enumerate}
\item Structure et géométrie des alcènes et des alcynes
\item Formules générales et écriture des formules
\item Nomenclature des alcènes et des alcynes
\item Isomérie Z et E (stéréoisomérie)
\item Propriétés chimiques : combustion et réactions d'addition
\item Polymérisation et préparation de l'acétylène
\item Travaux pratiques
\item Méthodes et exercices résolus
\item Exercices
\item Corrigés des exercices
\item Fiche bilan
\end{enumerate}
\end{multicols}
\end{sommaire}

\begin{objectifs}
\begin{itemize}
\item À la fin de cette leçon, je sais décrire la structure spatiale (géométrie, angles, longueurs de liaisons) des alcènes et des alcynes à partir des exemples de l'éthylène et de l'éthyne.
\item À la fin de cette leçon, je sais représenter géométriquement les alcènes et les alcynes à l'aide de modèles moléculaires ou de schémas.
\item À la fin de cette leçon, je sais écrire les formules générales, développées et semi-développées des alcènes (CnH2n) et des alcynes (CnH2n-2).
\item À la fin de cette leçon, je sais nommer un alcène ou un alcyne en appliquant les règles de nomenclature officielles.
\item À la fin de cette leçon, je sais distinguer et représenter les isomères Z et E d'un alcène.
\item À la fin de cette leçon, je sais écrire et équilibrer les équations des réactions d'addition (hydrogénation, hydratation, HCl, Cl2, Br2) en appliquant la règle de Markovnikov.
\item À la fin de cette leçon, je sais écrire l'équation d'une réaction de polymérisation d'un alcène et identifier le motif du polymère.
\item À la fin de cette leçon, je sais décrire et schématiser le dispositif expérimental de préparation de l'acétylène au laboratoire.
\item À la fin de cette leçon, je sais écrire l'équation de combustion complète d'un alcène ou d'un alcyne.
\end{itemize}
\end{objectifs}

\begin{prerequis}
\begin{itemize}
\item Les alcanes : formule générale CnH2n+2, nomenclature des chaînes carbonées (leçon 01)
\item La représentation de Lewis et la notion de liaison covalente (simple, double, triple)
\item L'hybridation du carbone tétragonal (sp3) vue dans la leçon sur les alcanes
\item L'équilibrage des équations-bilans chimiques (Seconde)
\item La notion d'isomérie de constitution vue avec les alcanes
\end{itemize}
\end{prerequis}

\newpage

\coursec{Structure et géométrie des alcènes et des alcynes}

Dans les ateliers de soudure de Douala et de Yaoundé, les soudeurs utilisent une flamme bleue très chaude alimentée par un gaz appelé \cle{acétylène}, capable de fondre l'acier. Dans les marchés, certains revendeurs font mûrir rapidement les bananes et les mangues grâce au \cle{carbure de calcium}, qui dégage ce même gaz au contact de l'eau. D'où vient ce gaz ? Pourquoi brûle-t-il avec une flamme si chaude ? Et comment la même famille de composés permet-elle de fabriquer les sachets plastiques et les tuyaux en PVC omniprésents au Cameroun ?

Toutes ces questions ont une même racine : la structure des molécules d'\emph{hydrocarbures insaturés}. Comprendre comment les atomes de carbone s'assemblent dans ces molécules — en double ou en triple liaison — c'est comprendre pourquoi elles sont si réactives et si utiles. C'est précisément l'objet de cette première partie : décrire la \cle{structure} et la \cle{géométrie} des alcènes et des alcynes.

\courssub{Les hydrocarbures insaturés}

\textbf{Des molécules qui n'ont pas « assez » d'hydrogène}\par

Dans la leçon précédente, tu as rencontré les \emph{alcanes} : des hydrocarbures dont tous les atomes de carbone sont reliés par des liaisons \emph{simples}. On dit qu'ils sont \cle{saturés} : chaque carbone porte le maximum d'atomes d'hydrogène possible. Leur formule générale est $C_nH_{2n+2}$.

Or il existe des hydrocarbures dans lesquels deux atomes de carbone sont reliés par une liaison \emph{multiple} : une \cle{double liaison} ou une \cle{triple liaison}. Ces molécules possèdent moins d'atomes d'hydrogène que l'alcane correspondant : on dit qu'elles sont \cle{insaturées}.

\begin{definition}[Alcènes et alcynes]
Un \cle{alcène} est un hydrocarbure \textbf{acylclique} insaturé possédant une \textbf{double liaison} carbone-carbone $\ce{C=C}$. Sa formule générale est $C_nH_{2n}$ (avec $n \geq 2$).

Un \cle{alcyne} est un hydrocarbure \textbf{acylclique} insaturé possédant une \textbf{triple liaison} carbone-carbone $\ce{C\bond{3}C}$. Sa formule générale est $C_nH_{2n-2}$ (avec $n \geq 2$).
\end{definition}

L'intuition est simple : former une liaison supplémentaire entre deux carbones « consomme » deux bras de liaison qui auraient pu porter des hydrogènes. Chaque liaison multiple fait donc perdre deux atomes d'hydrogène par rapport à l'alcane :
\begin{itemize}
\item éthane \ce{C2H6} (saturé) $\longrightarrow$ éthylène \ce{C2H4} (une double liaison, $-2$ H) ;
\item éthane \ce{C2H6} $\longrightarrow$ éthyne \ce{C2H2} (une triple liaison, $-4$ H).
\end{itemize}

Les deux plus petits représentants de ces familles, l'\cle{éthylène} \ce{C2H4} et l'\cle{éthyne} \ce{C2H2} (appelé couramment \emph{acétylène}), vont nous servir de modèles pour toute la leçon.

\courssub{L'éthylène \ce{C2H4} : une molécule plane}

\textbf{Observation expérimentale}\par

Les mesures physiques (diffraction des rayons X, spectroscopie) montrent que dans la molécule d'éthylène, les \textbf{six atomes sont dans un même plan}. Les angles de liaison autour de chaque carbone valent environ \SI{120}{\degree}, et la distance entre les deux carbones est d'environ \SI{134}{pm}, nettement plus courte que la liaison simple $\ce{C-C}$ de l'éthane (\SI{154}{pm}).

\begin{popfigure}
\centering
\begin{tikzpicture}
\node at (0,0) {\chemfig{H-[:-150]C(-[:-90]H)=[::0]C(-[::60]H)-[::-60]H}};
% angles
\draw[poporange,thick,->] (-1.05,0.18) arc[start angle=150,end angle=270,radius=0.42];
\node[poporange,font=\small] at (-1.75,-0.55) {$\approx \SI{120}{\degree}$};
\draw[poporange,thick,->] (1.05,-0.18) arc[start angle=-30,end angle=90,radius=0.42];
\node[poporange,font=\small] at (1.8,0.6) {$\approx \SI{120}{\degree}$};
% longueur C=C
\draw[popblue,thick,<->] (-0.42,0.55) -- (0.42,0.55);
\node[popblue,font=\small] at (0,0.85) {$d_{\ce{C=C}} \approx \SI{134}{pm}$};
% plan
\draw[popgreen!60,dashed,thick] (-3.2,-1.1) -- (-2.2,-1.6) -- (3.2,-1.6) -- (2.2,-1.1) -- cycle;
\node[popgreen!70!black,font=\small,align=center] at (0,-1.85) {les 6 atomes sont dans un même plan};
\end{tikzpicture}
\legende{Structure de l'éthylène \ce{C2H4} : molécule plane, angles voisins de \SI{120}{\degree}, double liaison $\ce{C=C}$ de longueur \SI{134}{pm}.}
\end{popfigure}

\begin{definition}[Carbone trigonal]
Un \cle{carbone trigonal} est un atome de carbone lié à \textbf{trois} atomes (ou groupes d'atomes). Il est dans l'état d'\textbf{hybridation $sp^2$} : ses trois directions de liaisons sont situées dans un même plan et pointent vers les sommets d'un triangle équilatéral, avec des angles de \SI{120}{\degree}. On dit que sa géométrie est \emph{plane triangulaire} (ou trigonale plane).
\end{definition}

Dans l'éthylène, chaque carbone est lié à trois atomes : deux hydrogènes et un carbone. Chaque carbone est donc \cle{trigonal}, et comme les deux triangles partagent le sommet commun $\ce{C=C}$, toute la molécule est \textbf{plane}.

\begin{exemplebox}[Vérification sur l'éthylène]
Pour le carbone de gauche de \ce{CH2=CH2} :
\begin{itemize}
\item nombre d'atomes liés : H, H, C $\Rightarrow$ \textbf{3 groupes} ;
\item géométrie prévue : plane triangulaire, angles de \SI{120}{\degree} ;
\item hybridation : $sp^2$ (une orbitale $s$ + deux orbitales $p$ hybridées).
\end{itemize}
La mesure expérimentale de l'angle $\widehat{\ce{HCH}}$ donne environ \SI{117}{\degree}, très proche de la valeur idéale : le modèle est confirmé.
\end{exemplebox}

\courssub{L'éthyne \ce{C2H2} : une molécule linéaire}

\textbf{Le gaz des soudeurs}\par

L'éthyne, plus connu sous le nom d'\cle{acétylène}, est le plus petit alcyne. Les mesures montrent que ses \textbf{quatre atomes sont alignés} : la molécule est \emph{linéaire}. L'angle $\widehat{\ce{HCC}}$ vaut exactement \SI{180}{\degree}, et la triple liaison $\ce{C\bond{3}C}$ ne mesure que \SI{120}{pm} environ : c'est la liaison carbone-carbone la plus courte et la plus forte.

\begin{popfigure}
\centering
\begin{tikzpicture}
\node at (0,0) {\chemfig{H-C~C-H}};
\draw[poporange,thick,->] (-0.9,0.45) arc[start angle=180,end angle=0,radius=0.9];
\node[poporange,font=\small] at (0,1.55) {$\widehat{\ce{HCC}} = \SI{180}{\degree}$ : molécule linéaire};
\draw[popblue,thick,<->] (-0.5,-0.55) -- (0.5,-0.55);
\node[popblue,font=\small] at (0,-0.9) {$d_{\ce{C\bond{3}C}} \approx \SI{120}{pm}$};
\end{tikzpicture}
\legende{Structure de l'éthyne (acétylène) \ce{C2H2} : les quatre atomes sont alignés, triple liaison $\ce{C\bond{3}C}$ de longueur \SI{120}{pm}.}
\end{popfigure}

\begin{definition}[Carbone digonal]
Un \cle{carbone digonal} est un atome de carbone lié à \textbf{deux} atomes (ou groupes d'atomes). Il est dans l'état d'\textbf{hybridation $sp$} : ses deux directions de liaisons sont \emph{alignées}, formant un angle de \SI{180}{\degree}. Sa géométrie est dite \emph{linéaire} (ou digonale).
\end{definition}

Dans l'éthyne $\ce{H-C\bond{3}C-H}$, chaque carbone n'est lié qu'à deux atomes : un hydrogène et un carbone. Chaque carbone est donc \cle{digonal}, d'où la linéarité parfaite de la molécule.

\begin{attention}[Ne pas confondre longueur et force de liaison]
Plus une liaison est multiple, plus elle est \textbf{courte} et plus elle est \textbf{forte} : $\ce{C-C}$ (\SI{154}{pm}) $>$ $\ce{C=C}$ (\SI{134}{pm}) $>$ $\ce{C\bond{3}C}$ (\SI{120}{pm}) en longueur, mais c'est l'inverse en énergie de liaison. Attention cependant : une triple liaison forte ne signifie pas une molécule peu réactive ! L'acétylène est au contraire très réactif, car une partie de ses liaisons (les liaisons $\pi$) est fragile et s'ouvre facilement lors des réactions d'addition.
\end{attention}

\courssub{Comparaison des trois types de carbone}

Tu connais désormais les trois « géométries » possibles du carbone dans les molécules organiques. Le point de départ du raisonnement est toujours le même : \emph{compter le nombre de groupes liés au carbone}.

\begin{methode}[Déterminer la géométrie autour d'un atome de carbone]
\begin{enumerate}
\item Repérer l'atome de carbone étudié dans la formule développée.
\item \textbf{Compter le nombre d'atomes} (ou groupes d'atomes) directement liés à ce carbone, sans tenir compte de la multiplicité des liaisons (une double ou triple liaison ne compte que pour \emph{une} direction).
\item Conclure :
\begin{itemize}
\item \textbf{4 groupes} $\Rightarrow$ carbone \cle{tétragonal} ($sp^3$), angles de \SI{109,5}{\degree} (soit $109\degree 28'$) ;
\item \textbf{3 groupes} $\Rightarrow$ carbone \cle{trigonal} ($sp^2$), angles de \SI{120}{\degree}, géométrie plane ;
\item \textbf{2 groupes} $\Rightarrow$ carbone \cle{digonal} ($sp$), angle de \SI{180}{\degree}, géométrie linéaire.
\end{itemize}
\end{enumerate}
\end{methode}

\begin{popfigure}
\centering
\begin{tikzpicture}
% sp3
\begin{scope}[shift={(-4.6,0)}]
\draw[thick,popdark] (0,0) -- (60:0.9);
\draw[thick,popdark] (0,0) -- (180:0.9);
\draw[thick,popdark] (0,0) -- (300:0.9);
\draw[thick,popdark,dashed] (0,0) -- (20:0.75);
\fill[popblue] (0,0) circle (0.16);
\fill[popmuted] (60:0.9) circle (0.10);
\fill[popmuted] (180:0.9) circle (0.10);
\fill[popmuted] (300:0.9) circle (0.10);
\fill[popmuted] (20:0.75) circle (0.10);
\node[font=\small\bfseries,popblue] at (0,-1.5) {tétragonal $sp^3$};
\node[font=\small] at (0,-1.95) {\SI{109,5}{\degree} — méthane \ce{CH4}};
\end{scope}
% sp2
\begin{scope}[shift={(0,0)}]
\draw[thick,popdark] (0,0) -- (90:0.9);
\draw[thick,popdark] (0,0) -- (210:0.9);
\draw[thick,popdark] (0,0) -- (330:0.9);
\fill[poporange] (0,0) circle (0.16);
\fill[popmuted] (90:0.9) circle (0.10);
\fill[popmuted] (210:0.9) circle (0.10);
\fill[popmuted] (330:0.9) circle (0.10);
\draw[popgreen!50,dashed] (90:1.05) -- (210:1.05) -- (330:1.05) -- cycle;
\node[font=\small\bfseries,poporange] at (0,-1.5) {trigonal $sp^2$};
\node[font=\small] at (0,-1.95) {\SI{120}{\degree} — éthylène \ce{C2H4}};
\end{scope}
% sp
\begin{scope}[shift={(4.6,0)}]
\draw[thick,popdark] (-0.9,0) -- (0.9,0);
\fill[poppurple] (0,0) circle (0.16);
\fill[popmuted] (-0.9,0) circle (0.10);
\fill[popmuted] (0.9,0) circle (0.10);
\node[font=\small\bfseries,poppurple] at (0,-1.5) {digonal $sp$};
\node[font=\small] at (0,-1.95) {\SI{180}{\degree} — éthyne \ce{C2H2}};
\end{scope}
\end{tikzpicture}
\legende{Les trois géométries de l'atome de carbone : tétragonale ($sp^3$), trigonale plane ($sp^2$) et digonale linéaire ($sp$).}
\end{popfigure}

\begin{center}
\begin{tabularx}{\linewidth}{|l|X|X|X|}
\hline
\rowcolor{popblueL} \textbf{Type de carbone} & \textbf{Tétragonal} & \textbf{Trigonal} & \textbf{Digonal} \\
\hline
Groupes liés & 4 & 3 & 2 \\
\hline
Hybridation & $sp^3$ & $sp^2$ & $sp$ \\
\hline
Géométrie & tétraédrique & plane triangulaire & linéaire \\
\hline
Angle de liaison & \SI{109,5}{\degree} & \SI{120}{\degree} & \SI{180}{\degree} \\
\hline
Type de liaison C--C & simple & double & triple \\
\hline
Exemple & \ce{CH4} (méthane) & \ce{C2H4} (éthylène) & \ce{C2H2} (éthyne) \\
\hline
\end{tabularx}
\end{center}

\begin{exoresolu}[Identifier les carbones d'une molécule]
Énoncé. On considère la molécule de méthylacétylène (propyne) de formule semi-développée $\ce{CH3-C\bond{3}CH}$. Pour chaque atome de carbone, préciser le nombre de groupes liés, l'hybridation et la géométrie locale.
\tcblower
Solution détaillée. On numérote les carbones : $\ce{C^1H3 - C^2 \bond{3} C^3H}$.
\begin{itemize}
\item \textbf{Carbone $\ce{C^1}$} : lié à 3 H et 1 C, soit \textbf{4 groupes} $\Rightarrow$ carbone \cle{tétragonal}, hybridation $sp^3$, géométrie tétraédrique (angles $\approx \SI{109,5}{\degree}$).
\item \textbf{Carbone $\ce{C^2}$} : lié à $\ce{C^1}$ et $\ce{C^3}$, soit \textbf{2 groupes} (la triple liaison ne compte que pour une direction) $\Rightarrow$ carbone \cle{digonal}, hybridation $sp$, géométrie linéaire.
\item \textbf{Carbone $\ce{C^3}$} : lié à $\ce{C^2}$ et H, soit \textbf{2 groupes} $\Rightarrow$ carbone \cle{digonal}, hybridation $sp$, géométrie linéaire.
\end{itemize}
Conséquence : les atomes $\ce{C^1}$, $\ce{C^2}$, $\ce{C^3}$ et le H terminal sont \textbf{alignés} ; seuls les trois hydrogènes du groupe $\ce{CH3}$ sortent de cet axe.
\end{exoresolu}

\courssub{Conséquence structurale majeure : la rotation bloquée autour de \ce{C=C}}

\textbf{Une liaison simple tourne, une double liaison non}\par

Autour d'une liaison \emph{simple} $\ce{C-C}$ (comme dans l'éthane), les deux moitiés de la molécule peuvent tourner librement l'une par rapport à l'autre : on parle de \cle{libre rotation}. C'est pourquoi les alcanes peuvent adopter de nombreuses conformations.

Il n'en est rien pour la double liaison $\ce{C=C}$. Une double liaison est constituée de deux liaisons de natures différentes : une liaison $\sigma$ (dans le plan de la molécule) et une liaison $\pi$ (au-dessus et au-dessous du plan). Tourner autour de l'axe $\ce{C=C}$ imposerait de \emph{rompre} la liaison $\pi$, ce qui demande une énergie considérable. À température ordinaire, la rotation est donc \textbf{bloquée}.

\begin{attention}[Erreur fréquente à éviter absolument]
Beaucoup d'élèves croient que la molécule d'éthylène peut tourner autour de la liaison $\ce{C=C}$, comme l'éthane autour de sa liaison simple. C'est \textbf{faux} : la double liaison $\ce{C=C}$ \textbf{bloque la rotation}. Les quatre atomes d'hydrogène restent fixés dans le plan de la molécule. Cette rigidité a une conséquence capitale que tu étudieras dans la suite de la leçon : lorsque chaque carbone de la double liaison porte deux groupes \emph{différents}, il existe deux arrangements spatiaux distincts et non interconvertibles — c'est l'\cle{isomérie Z/E}.
\end{attention}

\courssub{Représentation spatiale et modèles moléculaires}

Pour visualiser ces géométries, les chimistes utilisent des \cle{modèles moléculaires} : le modèle \emph{boules-bâtonnets} (les atomes sont des boules, les liaisons des bâtonnets) et le modèle \emph{compact} (les atomes sont des sphères emboîtées). La figure suivante représente les modèles boules-bâtonnets de nos deux molécules de référence.

\begin{popfigure}
\centering
\begin{tikzpicture}[ball/.style={circle,shading=ball,minimum size=#1,inner sep=0pt}]
% ---- Ethylene ----
\begin{scope}[shift={(-3.6,0)}]
\draw[line width=2.2pt,popmuted] (-0.55,0) -- (0.55,0);
\draw[line width=2.2pt,popmuted] (-0.55,0) -- (-1.5,0.75);
\draw[line width=2.2pt,popmuted] (-0.55,0) -- (-1.5,-0.75);
\draw[line width=2.2pt,popmuted] (0.55,0) -- (1.5,0.75);
\draw[line width=2.2pt,popmuted] (0.55,0) -- (1.5,-0.75);
\node[ball=0.62cm,ball color=popdark] at (-0.55,0) {};
\node[ball=0.62cm,ball color=popdark] at (0.55,0) {};
\node[ball=0.4cm,ball color=popblue!60] at (-1.5,0.75) {};
\node[ball=0.4cm,ball color=popblue!60] at (-1.5,-0.75) {};
\node[ball=0.4cm,ball color=popblue!60] at (1.5,0.75) {};
\node[ball=0.4cm,ball color=popblue!60] at (1.5,-0.75) {};
\node[font=\small\bfseries] at (0,-1.5) {éthylène \ce{C2H4} : plane};
\end{scope}
% ---- Acetylene ----
\begin{scope}[shift={(3.6,0)}]
\draw[line width=2.2pt,popmuted] (-1.5,0) -- (1.5,0);
\node[ball=0.62cm,ball color=popdark] at (-0.55,0) {};
\node[ball=0.62cm,ball color=popdark] at (0.55,0) {};
\node[ball=0.4cm,ball color=popblue!60] at (-1.5,0) {};
\node[ball=0.4cm,ball color=popblue!60] at (1.5,0) {};
\node[font=\small\bfseries] at (0,-1.5) {éthyne \ce{C2H2} : linéaire};
\end{scope}
% legend
\node[ball=0.3cm,ball color=popdark] at (-5.6,-2.3) {};
\node[font=\small,right] at (-5.4,-2.3) {carbone};
\node[ball=0.24cm,ball color=popblue!60] at (-3.4,-2.3) {};
\node[font=\small,right] at (-3.2,-2.3) {hydrogène};
\end{tikzpicture}
\legende{Modèles boules-bâtonnets de l'éthylène (molécule plane, angles de \SI{120}{\degree}) et de l'éthyne (molécule linéaire, angle de \SI{180}{\degree}).}
\end{popfigure}

\begin{experience}[Construire les modèles moléculaires de \ce{C2H4} et \ce{C2H2}]
\textbf{Matériel} : boules noires percées de 4 trous (carbone), boules blanches (hydrogène), bâtonnets rigides et ressorts.

\textbf{Protocole} :
\begin{enumerate}
\item Pour l'éthylène : relier deux boules noires par \emph{deux} connecteurs (double liaison), puis fixer deux boules blanches sur chaque carbone de façon que les trois directions soient dans un même plan, à \SI{120}{\degree}.
\item Pour l'éthyne : relier deux boules noires par \emph{trois} connecteurs (triple liaison), puis fixer une boule blanche sur chaque carbone, dans le prolongement de l'axe $\ce{C\bond{3}C}$.
\item Essayer de faire tourner une moitié de chaque modèle par rapport à l'autre.
\end{enumerate}

\textbf{Observations} : le modèle de l'éthylène est rigide et plat ; impossible de le tordre sans déboîter les connecteurs. Le modèle de l'éthyne est parfaitement droit.

\textbf{Interprétation} : les modèles traduisent fidèlement la planéité de l'éthylène (carbones trigonaux $sp^2$), la linéarité de l'éthyne (carbones digonaux $sp$) et le blocage de la rotation autour de la double liaison.
\end{experience}

\begin{savaistu}[L'éthylène, l'hormone qui fait mûrir les bananes]
L'éthylène \ce{C2H4} n'est pas qu'un produit de l'industrie : c'est une \textbf{hormone végétale naturelle}. Les fruits la dégagent eux-mêmes pour déclencher et synchroniser leur mûrissement. C'est pourquoi une banane mûre placée dans un sac avec des mangues vertes accélère leur mûrissement ! Au Cameroun, certains revendeurs utilisent du \cle{carbure de calcium} \ce{CaC2} : au contact de l'humidité, il dégage de l'acétylène \ce{C2H2} selon la réaction $\ce{CaC2 + 2H2O -> Ca(OH)2 + C2H2}$, un gaz voisin de l'éthylène qui provoque un mûrissement rapide (mais parfois de qualité inférieure, et dont l'usage est encadré pour des raisons sanitaires). Quant au chalumeau oxyacétylénique des soudeurs de Douala, il exploite la combustion de ce même acétylène dans le dioxygène, dont la flamme atteint environ \SI{3000}{\celsius} :
\[ \ce{2C2H2 + 5O2 -> 4CO2 + 2H2O} \]
Cette combustion très exothermique s'explique par la grande quantité d'énergie stockée dans la molécule d'acétylène.
\end{savaistu}

\begin{aretenir}[L'essentiel de cette section]
\begin{itemize}
\item Les \cle{alcènes} ($C_nH_{2n}$) possèdent une double liaison $\ce{C=C}$ ; les \cle{alcynes} ($C_nH_{2n-2}$) possèdent une triple liaison $\ce{C\bond{3}C}$ : ce sont des hydrocarbures \emph{insaturés} (acylcliques).
\item L'éthylène \ce{C2H4} est une molécule \textbf{plane} : carbones \cle{trigonaux} ($sp^2$), angles de \SI{120}{\degree}, $d_{\ce{C=C}} \approx \SI{134}{pm}$.
\item L'éthyne \ce{C2H2} (acétylène) est une molécule \textbf{linéaire} : carbones \cle{digonaux} ($sp$), angle de \SI{180}{\degree}, $d_{\ce{C\bond{3}C}} \approx \SI{120}{pm}$.
\item Pour trouver la géométrie d'un carbone, on compte ses groupes liés : 4 $\Rightarrow$ tétragonal (\SI{109,5}{\degree}) ; 3 $\Rightarrow$ trigonal (\SI{120}{\degree}) ; 2 $\Rightarrow$ digonal (\SI{180}{\degree}).
\item La double liaison $\ce{C=C}$ \textbf{bloque la rotation} : c'est l'origine de l'isomérie Z/E, que nous étudierons bientôt.
\end{itemize}
\end{aretenir}

\coursec{Formules générales et écriture des formules}

Dans la section précédente, nous avons découvert que les alcènes possèdent une \cle{double liaison} \ce{C=C} entre deux carbones \cle{trigonaux}, tandis que les alcynes possèdent une \cle{triple liaison} \ce{C#C} entre deux carbones \cle{digonaux}. Nous savons donc maintenant \emph{à quoi ressemblent} ces molécules dans l'espace. Mais un chimiste, au laboratoire ou devant sa copie de Probatoire, doit aussi savoir les \emph{écrire} rapidement et sans erreur. C'est l'objet de cette section : établir les \cle{formules générales} des alcènes et des alcynes, puis apprendre à passer d'une formule brute à une formule développée, semi-développée ou topologique, et enfin à dénombrer les isomères de constitution d'une formule brute donnée.

\courssub{Formule générale des alcènes : $C_nH_{2n}$}

\textbf{Intuition.} Partons d'un alcane, par exemple l'éthane \ce{C2H6}, dont tous les carbones sont tétragonaux et reliés par des liaisons simples. Pour fabriquer l'éthylène \ce{C2H4}, il faut créer une double liaison \ce{C=C} : or, pour respecter la tétravalence du carbone, chacun des deux carbones doit \emph{lâcher} un atome d'hydrogène. La création d'une double liaison fait donc disparaître \textbf{deux atomes d'hydrogène} par rapport à l'alcane correspondant.

\begin{propriete}[Formule générale des alcènes]
Un alcène \cle{non cyclique} possédant \textbf{une seule double liaison} \ce{C=C} a pour formule générale :
\[ C_nH_{2n} \qquad \text{avec } n \geq 2 \]
où $n$ est le nombre d'atomes de carbone. La condition $n \geq 2$ s'explique : il faut au minimum deux carbones pour former une double liaison \ce{C=C}.
\end{propriete}

\begin{attention}[Ne pas confondre avec les alcanes]
L'erreur la plus fréquente au Probatoire est de confondre les trois formules générales :
\begin{itemize}
\item alcane : $C_nH_{2n+2}$ (que des liaisons simples) ;
\item alcène : $C_nH_{2n}$ (une double liaison : \textbf{moins 2 hydrogènes}) ;
\item alcyne : $C_nH_{2n-2}$ (une triple liaison : \textbf{moins 4 hydrogènes}).
\end{itemize}
Retiens cette règle simple : \textbf{chaque liaison multiple supplémentaire enlève 2 atomes d'hydrogène}. Une double liaison en enlève 2, une triple liaison en enlève 4 par rapport à l'alcane.
\end{attention}

\textbf{Vérification sur les exemples.} Appliquons la formule $C_nH_{2n}$ aux premiers alcènes :
\begin{itemize}
\item éthylène (éthène) : $n = 2$, donc $2n = 4$ hydrogènes $\Rightarrow$ \ce{C2H4} \checkmark ;
\item propène : $n = 3$, donc $2n = 6$ hydrogènes $\Rightarrow$ \ce{C3H6} \checkmark ;
\item butène : $n = 4$, donc $2n = 8$ hydrogènes $\Rightarrow$ \ce{C4H8} \checkmark.
\end{itemize}

\courssub{Formule générale des alcynes : $C_nH_{2n-2}$}

\textbf{Intuition.} Reprenons le même raisonnement. Pour transformer l'éthane \ce{C2H6} en éthyne (acétylène) \ce{C2H2}, il faut créer une triple liaison \ce{C#C}. Chaque carbone ne peut alors porter qu'\emph{un seul} hydrogène : la triple liaison a fait disparaître \textbf{quatre atomes d'hydrogène} par rapport à l'alcane.

\begin{propriete}[Formule générale des alcynes]
Un alcyne \cle{non cyclique} possédant \textbf{une seule triple liaison} \ce{C#C} a pour formule générale :
\[ C_nH_{2n-2} \qquad \text{avec } n \geq 2 \]
\end{propriete}

\textbf{Vérification sur les exemples.}
\begin{itemize}
\item éthyne (acétylène) : $n = 2$, donc $2n - 2 = 2$ hydrogènes $\Rightarrow$ \ce{C2H2} \checkmark ;
\item propyne : $n = 3$, donc $2n - 2 = 4$ hydrogènes $\Rightarrow$ \ce{C3H4} \checkmark ;
\item butyne : $n = 4$, donc $2n - 2 = 6$ hydrogènes $\Rightarrow$ \ce{C4H6} \checkmark.
\end{itemize}

\begin{methode}[Retrouver rapidement la formule brute]
Pour trouver la formule brute d'un hydrocarbure insaturé à $n$ atomes de carbone :
\begin{enumerate}
\item \textbf{Alcène} : double le nombre $n$ de carbones pour obtenir le nombre d'hydrogènes $\Rightarrow C_nH_{2n}$.
\item \textbf{Alcyne} : double $n$, puis retire 2 $\Rightarrow C_nH_{2n-2}$.
\item \textbf{Vérification} : contrôle toujours la tétravalence de chaque carbone sur ta formule développée (chaque carbone doit former exactement 4 liaisons).
\end{enumerate}
\end{methode}

\begin{exemplebox}[Exemple chiffré : retrouver $n$ à partir de la masse molaire]
Un alcène a une masse molaire $M = \SI{56}{g.mol^{-1}}$. Quelle est sa formule brute ?

La masse molaire d'un alcène $C_nH_{2n}$ vaut :
\[ M = 12\,n + 1 \times 2n = 14\,n \]
On résout :
\[ 14\,n = 56 \quad \Rightarrow \quad n = \frac{56}{14} = 4 \]
La formule brute est donc \ce{C4H8} (un butène). \textbf{Vérification} : $M = 4 \times 12 + 8 \times 1 = 48 + 8 = \SI{56}{g.mol^{-1}}$ \checkmark
\end{exemplebox}

\begin{savaistu}[Un grand classique du Probatoire]
Au Probatoire C et D, la question « déterminer la formule brute d'un hydrocarbure à partir de sa masse molaire ou de sa densité de vapeur » revient presque chaque année. Rappelle-toi le lien entre densité de vapeur et masse molaire : $M = 29\,d$ (avec $d$ la densité par rapport à l'air). Ainsi, un alcyne de densité $d \approx 1{,}86$ a pour masse molaire $M \approx 29 \times 1{,}86 \approx \SI{54}{g.mol^{-1}}$, d'où $12n + 2n - 2 = 54$, soit $n = 4$ : c'est \ce{C4H6}. Entraîne-toi sur ce type de calcul : il rapporte des points faciles !
\end{savaistu}

\courssub{De la formule brute à la formule développée}

La \cle{formule brute} (ex. \ce{C4H8}) indique seulement la \emph{composition} de la molécule : elle ne dit rien sur la manière dont les atomes sont reliés. Pour représenter la structure, on utilise la \cle{formule développée}.

\begin{definition}[Formule développée]
La \cle{formule développée} d'une molécule est une représentation plane dans laquelle \textbf{tous les atomes et toutes les liaisons} (simples, doubles ou triples) sont explicitement dessinés.
\end{definition}

\textbf{Méthode de construction} (exemple du but-1-ène \ce{C4H8}) :
\begin{enumerate}
\item Écrire le squelette carboné : 4 carbones en chaîne \ce{C-C-C-C}.
\item Placer la double liaison à l'extrémité (position 1) : \ce{C=C-C-C}.
\item Compléter chaque carbone avec des hydrogènes jusqu'à ce qu'il forme \textbf{exactement 4 liaisons} : le premier carbone (2 liaisons vers le carbone voisin) reçoit 2 H ; le deuxième (3 liaisons déjà engagées) reçoit 1 H ; les carbones 3 et 4 reçoivent respectivement 2 H et 3 H.
\end{enumerate}

On obtient ainsi la formule développée du but-1-ène, représentée dans le tableau ci-dessous.

\begin{popfigure}
\begin{center}
\begin{tabularx}{\linewidth}{|l|X|}
\hline
\rowcolor{popblueL} \textbf{Type de formule} & \textbf{But-1-ène (\ce{C4H8})} \\
\hline
Formule brute & \ce{C4H8} \\
\hline
Formule développée & \chemfig{H-C(-[2]H)=C(-[2]H)-C(-[2]H)(-[6]H)-C(-[2]H)(-[6]H)-H} \\
\hline
Formule semi-développée & \chemfig{CH_2=CH-CH_2-CH_3} \\
\hline
Formule topologique & \chemfig{=[:30]-[:-30]} \\
\hline
\end{tabularx}
\end{center}
\legende{Les quatre écritures du but-1-ène : de la formule brute à la formule topologique.}
\end{popfigure}

\courssub{Formule semi-développée : l'écriture du chimiste pressé}

Dessiner toutes les liaisons devient vite fastidieux pour les grosses molécules. On utilise alors la \cle{formule semi-développée}, qui est l'écriture la plus utilisée dans les exercices.

\begin{definition}[Formule semi-développée]
La \cle{formule semi-développée} est obtenue en \textbf{regroupant sur chaque atome de carbone les atomes d'hydrogène} qui lui sont liés. On n'écrit plus les liaisons \ce{C-H}, mais on conserve les liaisons entre carbones et, obligatoirement, les liaisons multiples.
\end{definition}

\textbf{Exemples à connaître par cœur :}
\begin{itemize}
\item éthylène : \chemfig{CH_2=CH_2} ;
\item propène : \chemfig{CH_2=CH-CH_3} ;
\item but-2-ène : \chemfig{CH_3-CH=CH-CH_3} ;
\item éthyne (acétylène) : \chemfig{CH#CH} ;
\item propyne : \chemfig{CH#C-CH_3}.
\end{itemize}

\begin{attention}[Pièges fréquents des formules semi-développées]
\begin{itemize}
\item \textbf{Jamais} de formule semi-développée sans la liaison multiple : écrire \ce{CH2CH2} pour l'éthylène est faux, car on croirait lire un fragment d'alcane. Écris toujours \chemfig{CH_2=CH_2}.
\item Vérifie la tétravalence : dans \chemfig{CH_2=CH-CH_3}, le carbone central porte bien 1 H car il forme déjà 3 liaisons (1 double + 1 simple). Un carbone à 5 liaisons est une erreur qui coûte cher !
\item Pour les alcynes, la triple liaison doit être visible : \chemfig{CH#C-CH_3} et non \ce{CHC-CH3}.
\end{itemize}
\end{attention}

\textbf{Formules topologiques (mention).} Pour les chaînes carbonées longues, les chimistes utilisent la \cle{formule topologique} : les carbones et les hydrogènes qui leur sont liés ne sont plus écrits ; chaque sommet et chaque extrémité de segment représente un carbone, et la double liaison est figurée par un double trait. Par exemple, le but-1-ène se note simplement \chemfig{=[:30]-[:-30]}. Tu la rencontreras surtout en Terminale et dans les documents scientifiques ; au Probatoire, la formule semi-développée reste l'écriture attendue.

\courssub{Isomères de constitution d'une formule brute donnée}

Une même formule brute peut correspondre à \textbf{plusieurs molécules différentes} : ce sont des \cle{isomères de constitution}. Ils diffèrent par :
\begin{itemize}
\item la \textbf{position} de la liaison multiple dans la chaîne (isomérie de position) ;
\item la \textbf{ramification} de la chaîne carbonée (isomérie de chaîne).
\end{itemize}

\begin{methode}[Recherche systématique des isomères d'un alcène $C_nH_{2n}$]
\begin{enumerate}
\item Écrire la chaîne carbonée \textbf{linéaire} à $n$ carbones, puis déplacer la double liaison le long de la chaîne (attention aux positions équivalentes par symétrie !).
\item \textbf{Raccourcir la chaîne principale} d'un carbone et greffer le carbone retiré comme ramification (méthyle), puis replacer la double liaison à toutes les positions possibles.
\item Répéter si possible, puis vérifier qu'aucun isomère n'est compté deux fois et que chaque carbone reste tétravalent.
\end{enumerate}
\end{methode}

\textbf{Application : les isomères de constitution de \ce{C4H8}.}
\begin{enumerate}
\item Chaîne linéaire à 4 carbones, double liaison en position 1 : \textbf{but-1-ène} \chemfig{CH_2=CH-CH_2-CH_3}.
\item Chaîne linéaire, double liaison en position 2 : \textbf{but-2-ène} \chemfig{CH_3-CH=CH-CH_3} (la position 3 est équivalente à la position 1 par symétrie : on s'arrête là).
\item Chaîne à 3 carbones + ramification méthyle : \textbf{2-méthylprop-1-ène} \chemfig{CH_2=C(-[2]CH_3)-CH_3}.
\end{enumerate}

Il existe donc \textbf{trois isomères de constitution} pour \ce{C4H8}, représentés ci-dessous.

\begin{popfigure}
\begin{center}
\chemfig{CH_2=CH-CH_2-CH_3}
\hspace{1.2cm}
\chemfig{CH_3-CH=CH-CH_3}
\hspace{1.2cm}
\chemfig{CH_2=C(-[2]CH_3)-CH_3}
\end{center}
\begin{center}
\textbf{but-1-ène} \hspace{1.6cm} \textbf{but-2-ène} \hspace{1.6cm} \textbf{2-méthylprop-1-ène}
\end{center}
\legende{Les trois isomères de constitution de formule brute \ce{C4H8}.}
\end{popfigure}

\begin{attention}[Ne pas compter deux fois le même isomère]
Le but-1-ène et le « but-3-ène » sont la \emph{même} molécule : on numérote toujours la chaîne de façon à donner à la double liaison \textbf{le plus petit indice possible}. De même, « but-4-ène » n'existe pas. Avant d'ajouter un isomère à ta liste, retourne la molécule mentalement pour vérifier qu'elle n'y figure pas déjà.
\end{attention}

\courssub{Application aux alcynes : exercice résolu}

\begin{exoresolu}[Isomères de constitution des alcynes de formule \ce{C5H8}]
\textbf{Énoncé.} Un alcyne a pour formule brute \ce{C5H8}.
\begin{enumerate}
\item Vérifier que cette formule est compatible avec celle d'un alcyne.
\item Écrire les formules semi-développées de tous les isomères de constitution de cet alcyne.
\end{enumerate}
\tcblower
\textbf{Solution détaillée.}

\textbf{1.} Pour un alcyne, $C_nH_{2n-2}$. Avec $n = 5$ : $2n - 2 = 10 - 2 = 8$. La formule \ce{C5H8} est bien celle d'un alcyne à 5 carbones. \checkmark

\textbf{2.} Recherche systématique :

\emph{Étape 1 : chaîne linéaire à 5 carbones.}
\begin{itemize}
\item Triple liaison en position 1 : \textbf{pent-1-yne} \quad \chemfig{CH#C-CH_2-CH_2-CH_3}
\item Triple liaison en position 2 : \textbf{pent-2-yne} \quad \chemfig{CH_3-C#C-CH_2-CH_3}
\end{itemize}
(Les positions 3 et 4 sont équivalentes aux positions 2 et 1 par symétrie.)

\emph{Étape 2 : chaîne à 4 carbones + ramification méthyle.}
On part du squelette but-1-yne \ce{CH#C-CH2-CH3}. Les carbones de la triple liaison (C1 et C2) sont \emph{hybridés sp} : ils forment déjà deux liaisons $\sigma$ (une vers l'autre carbone, une vers H ou C), ils ne peuvent \textbf{pas} porter de ramification supplémentaire. La ramification ne est possible que sur C3.
\begin{itemize}
\item Méthyle sur C3 : \textbf{3-méthylbut-1-yne} \quad \chemfig{CH#C-CH(-[2]CH_3)-CH_3}
\end{itemize}

\emph{Étape 3 : chaîne à 3 carbones ?} Une chaîne propane avec deux ramifications méthyle sur le carbone central donnerait un carbone à 5 liaisons : impossible.

\textbf{Conclusion :} il existe \textbf{trois isomères de constitution} : le pent-1-yne, le pent-2-yne et le 3-méthylbut-1-yne.
\end{exoresolu}

\begin{exercice}[À toi de jouer]
\begin{enumerate}
\item Un alcène a une densité de vapeur $d = 1{,}45$. Détermine sa formule brute, puis écris les formules semi-développées de tous ses isomères de constitution.
\item Écris la formule développée du 2-méthylbut-2-ène, puis sa formule semi-développée. Vérifie que sa formule brute est bien $C_nH_{2n}$.
\end{enumerate}
\end{exercice}

\begin{corrige}[Exercice]
\textbf{1.} $M = 29\,d = 29 \times 1{,}45 \approx \SI{42}{g.mol^{-1}}$. Pour un alcène, $M = 14n$, d'où $n = 3$ : formule brute \ce{C3H6}. Un seul isomère de constitution : le propène \chemfig{CH_2=CH-CH_3} (déplacer la double liaison donne la même molécule par symétrie).

\textbf{2.} Le 2-méthylbut-2-ène a pour chaîne principale le but-2-ène (4 carbones, double liaison en 2) avec un méthyle sur le C2.
\begin{itemize}
\item \textbf{Formule développée :} \chemfig{H_3C-C(-[2]CH_3)=C(-[2]H)-CH_3} (chaque carbone forme 4 liaisons).
\item \textbf{Formule semi-développée :} \chemfig{CH_3-C(-[2]CH_3)=CH-CH_3}.
\item \textbf{Formule brute :} 5 atomes C, 10 atomes H $\Rightarrow$ \ce{C5H10}. Vérification : $n=5$, $2n=10$ \checkmark
\end{itemize}
\end{corrige}

\begin{aretenir}[L'essentiel de la section]
\begin{itemize}
\item Alcène non cyclique à une double liaison : $C_nH_{2n}$ ($n \geq 2$) ; alcyne non cyclique à une triple liaison : $C_nH_{2n-2}$ ($n \geq 2$).
\item Chaque liaison multiple enlève 2 H par rapport à l'alcane $C_nH_{2n+2}$.
\item Formule développée : toutes les liaisons dessinées ; formule semi-développée : hydrogènes regroupés par carbone, liaisons multiples toujours visibles ; formule topologique : sommets = carbones.
\item Recherche d'isomères : chaîne linéaire d'abord (on déplace la liaison multiple), puis ramifications ; toujours vérifier la tétravalence et éviter les doublons.
\item Formule brute à partir de la masse molaire : alcène $M = 14n$ ; alcyne $M = 14n - 2$ ; et $M = 29d$ avec la densité de vapeur.
\end{itemize}
\end{aretenir}

Maintenant que tu sais écrire et compter ces molécules, il reste à apprendre à les \emph{nommer} rigoureusement : c'est l'objet de la section suivante, consacrée à la nomenclature des alcènes et des alcynes.

\coursec{Nomenclature des alcènes et des alcynes}

Dans la section précédente, tu as appris à écrire les formules développées et semi-développées des alcènes ($C_nH_{2n}$) et des alcynes ($C_nH_{2n-2}$). Mais une formule comme \ce{CH3-CH=CH-CH2-CH3} ne suffit pas à communiquer : deux chimistes, l'un à Yaoundé et l'autre à Tokyo, doivent pouvoir évoquer la même molécule avec le même nom. C'est toute l'utilité de la \cle{nomenclature} : un langage international, codifié par l'IUPAC (Union internationale de chimie pure et appliquée), qui permet de passer d'une formule à un nom, et inversement, sans aucune ambiguïté.

Bonne nouvelle : tu connais déjà l'essentiel des règles, car elles prolongent celles des alcanes. Il suffit d'y ajouter une idée nouvelle, simple mais capitale : \textbf{la liaison multiple commande tout}. C'est elle qui impose le choix de la chaîne principale, le sens de la numérotation et le suffixe du nom. Voyons cela pas à pas.

\courssub{Les quatre règles fondamentales}

\textbf{Règle 1 : la chaîne principale doit contenir la liaison multiple}\par

\begin{definition}[Chaîne principale d'un alcène ou d'un alcyne]
La \cle{chaîne principale} d'un hydrocarbure insaturé est la chaîne carbonée la plus longue \textbf{qui contient la liaison multiple} (double ou triple). Elle détermine la racine du nom : éth- (2 C), prop- (3 C), but- (4 C), pent- (5 C), hex- (6 C), etc.
\end{definition}

L'intuition est la suivante : la liaison multiple est la « personnalité » de la molécule, son groupe caractéristique. Une chaîne qui la laisserait de côté ne pourrait pas la décrire. Il peut donc arriver que la chaîne la plus longue de la molécule ne soit \emph{pas} la chaîne principale, si elle ne contient pas la liaison multiple !

\begin{attention}[Piège classique]
Pour \chemfig{CH_2=C(-[2]CH_3)-CH_2-CH_3}, la plus longue chaîne de la molécule compte 4 carbones et contient bien la double liaison : la racine est donc « but ». Mais si l'on imaginait une molécule où la plus longue chaîne possible (5 carbones) passerait à côté de la double liaison, il faudrait lui préférer une chaîne de 4 carbones contenant \ce{C=C}. \textbf{Critère prioritaire : contenir la liaison multiple ; critère secondaire : être la plus longue.}
\end{attention}

\textbf{Règle 2 : la numérotation donne à la liaison multiple l'indice le plus petit}\par

\begin{definition}[Numérotation de la chaîne principale]
On numérote les carbones de la chaîne principale de façon à attribuer à la \cle{liaison multiple} l'\textbf{indice de position le plus petit possible}. Pour une double ou triple liaison entre les carbones $i$ et $i+1$, on retient l'indice $i$ (celui du premier carbone de la liaison).
\end{definition}

C'est la différence majeure avec les alcanes : chez les alcanes, on numérotait pour donner les indices les plus petits aux ramifications. Ici, la liaison multiple est \textbf{toujours prioritaire} sur les groupes alkyles, même si cela donne de plus gros numéros aux substituants.

\textbf{Règle 3 : le suffixe et l'indice de position}\par

\begin{definition}[Suffixes des hydrocarbures insaturés]
Le nom se termine par le suffixe :
\begin{itemize}
\item \cle{-ène} pour les alcènes (liaison double) ;
\item \cle{-yne} pour les alcynes (liaison triple).
\end{itemize}
Ce suffixe est précédé de l'indice de position de la liaison multiple, encadré de tirets : \textbf{but-1-ène}, \textbf{pent-2-yne}. L'ancienne écriture (but-1-ène noté « 1-butène ») est encore tolérée mais l'écriture moderne place l'indice juste devant le suffixe.
\end{definition}

\textbf{Règle 4 : les ramifications alkyles}\par

\begin{definition}[Nommage des substituants]
Les ramifications alkyles (méthyle \ce{CH3}-, éthyle \ce{C2H5}-, propyle...) sont nommées \textbf{par ordre alphabétique}, chacune précédée de son indice de position. Si un même groupe apparaît plusieurs fois, on utilise les préfixes di-, tri-, tétra- et on répète les indices : 2,3-diméthyl... Les préfixes di-, tri- ne comptent pas dans l'ordre alphabétique.
\end{definition}

\begin{methode}[Algorithme de nommage en 4 étapes]
Pour nommer un alcène ou un alcyne, applique méthodiquement :
\begin{enumerate}
\item \textbf{Identifier la chaîne principale} : la plus longue chaîne contenant la liaison multiple. Elle donne la racine (éth-, prop-, but-, pent-...).
\item \textbf{Numéroter} dans le sens qui donne à la liaison multiple l'indice le plus petit. La liaison multiple est prioritaire sur tout substituant.
\item \textbf{Écrire le suffixe} : « -ène » (double) ou « -yne » (triple), précédé de l'indice de position : racine-indice-suffixe (ex. pent-2-ène).
\item \textbf{Ajouter les substituants} en préfixe, par ordre alphabétique, avec leurs indices : ex. 4-méthylpent-2-ène.
\end{enumerate}
\end{methode}

\courssub{Premiers exemples progressifs}

\begin{exemplebox}[Exemple 1 : \ce{CH2=CH-CH2-CH3}]
\textbf{Étape 1} : la chaîne contenant la double liaison compte 4 carbones $\Rightarrow$ racine « but ».\\
\textbf{Étape 2} : en numérotant depuis la gauche, la double liaison porte l'indice 1 ; depuis la droite, elle porterait l'indice 3. On choisit 1.\\
\textbf{Étape 3} : suffixe « -ène ».\\
\textbf{Nom} : \textbf{but-1-ène}.
\end{exemplebox}

\begin{attention}[Erreur fréquente : « but-3-ène »]
Écrire « but-3-ène » pour \ce{CH2=CH-CH2-CH3} est une faute : la numérotation doit donner l'indice \textbf{le plus petit} à la liaison multiple. Or but-3-ène et but-1-ène décrivent la \emph{même} molécule (il suffit de la retourner) ; la nomenclature n'admet qu'un seul nom : \textbf{but-1-ène}. De même, « pent-4-yne » n'existe pas : c'est le pent-1-yne.
\end{attention}

\begin{exemplebox}[Exemple 2 : \ce{CH3-C##C-CH3}]
\textbf{Étape 1} : 4 carbones $\Rightarrow$ « but ». \textbf{Étape 2} : la triple liaison est entre C2 et C3 quel que soit le sens $\Rightarrow$ indice 2. \textbf{Étape 3} : suffixe « -yne ».\\
\textbf{Nom} : \textbf{but-2-yne}. Remarque : \ce{CH#C-CH2-CH3} serait le but-1-yne, isomère de position du but-2-yne : même formule brute \ce{C4H6}, mais triple liaison à des positions différentes.
\end{exemplebox}

\begin{exemplebox}[Exemple 3 : \chemfig{CH_2=C(-[2]CH_3)-CH_3}]
\textbf{Étape 1} : la plus longue chaîne contenant \ce{C=C} compte 3 carbones $\Rightarrow$ « prop ». \textbf{Étape 2} : la double liaison est nécessairement en position 1 (prop-1-ène ramifié). \textbf{Étape 4} : un groupe méthyle sur le carbone 2.\\
\textbf{Nom} : \textbf{2-méthylpropène} (on omet l'indice de la double liaison, car aucune ambiguïté n'est possible : prop-1-ène est la seule position envisageable pour un propène).
\end{exemplebox}

\courssub{Cas des composés ramifiés : priorité de la liaison multiple}

C'est ici que se joue le point le plus délicat de la leçon. Considérons la molécule \ce{CH3-CH=CH-CH(CH3)-CH3}.

\begin{exoresolu}[Nommer \ce{CH3-CH=CH-CH(CH3)-CH3}]
Énoncé : donner le nom IUPAC de la molécule de formule semi-développée \ce{CH3-CH=CH-CH(CH3)-CH3}.
\tcblower
\textbf{Solution détaillée.}

\textbf{Étape 1 — Chaîne principale.} La plus longue chaîne contenant la double liaison compte 5 carbones : racine « pent ».

\textbf{Étape 2 — Numérotation.} Deux sens sont possibles :
\begin{itemize}
\item de gauche à droite : la double liaison est en position \textbf{2}, le méthyle en position 4 ;
\item de droite à gauche : la double liaison serait en position 3, le méthyle en position 2.
\end{itemize}
La \textbf{liaison multiple est prioritaire} : on retient le sens qui lui donne l'indice le plus petit, c'est-à-dire 2, même si le méthyle reçoit alors l'indice 4.

\textbf{Étape 3 — Suffixe} : « -ène » avec l'indice 2 $\Rightarrow$ pent-2-ène.

\textbf{Étape 4 — Substituant} : un méthyle en position 4.

\textbf{Nom : 4-méthylpent-2-ène.} Le nom « 2-méthylpent-3-ène » est \textbf{faux}, bien qu'il donne un indice plus petit au substituant : il attribue à la double liaison l'indice 3 au lieu de 2.
\end{exoresolu}

\begin{attention}[La liaison multiple est TOUJOURS prioritaire]
L'erreur la plus fréquente au Probatoire consiste à numéroter la chaîne pour donner le plus petit indice au \textbf{groupe alkyle} (comme on le faisait pour les alcanes). C'est faux pour les hydrocarbures insaturés : la hiérarchie des priorités est
\[
\text{liaison multiple} \;>\; \text{ramifications alkyles}.
\]
On ne compare les indices des substituants que si les deux sens de numérotation donnent le \emph{même} indice à la liaison multiple.
\end{attention}

\begin{popfigure}
\begin{center}
\begin{tikzpicture}[scale=0.9, every node/.style={font=\small}]
% Numérotation correcte
\node[font=\bfseries\color{popgreen}] at (3.2,3.1) {Numérotation correcte : 4-méthylpent-2-ène};
\node at (3.2,2.2) {\chemfig{CH_3-[0]CH=[0]CH-[0]CH(-[2]CH_3)-[0]CH_3}};
\foreach \x/\n in {0.6/1, 1.75/2, 2.9/3, 4.05/4, 5.2/5} {
  \node[popgreen, font=\bfseries] at (\x,1.15) {\n};
}
\draw[popgreen, thick, ->] (2.35,1.55) -- (2.35,1.85);
\node[popgreen, font=\footnotesize] at (2.35,0.7) {double liaison : indice 2};
% Numérotation incorrecte
\node[font=\bfseries\color{popink}] at (10.8,3.1) {Numérotation incorrecte : « 2-méthylpent-3-ène »};
\node at (10.8,2.2) {\chemfig{CH_3-[0]CH=[0]CH-[0]CH(-[2]CH_3)-[0]CH_3}};
\foreach \x/\n in {8.2/5, 9.35/4, 10.5/3, 11.65/2, 12.8/1} {
  \node[popink, font=\bfseries] at (\x,1.15) {\n};
}
\draw[popink, thick, ->] (9.95,1.55) -- (9.95,1.85);
\node[popink, font=\footnotesize] at (9.95,0.7) {double liaison : indice 3 $\Rightarrow$ rejeté};
\end{tikzpicture}
\end{center}
\legende{Les deux numérotations possibles du 4-méthylpent-2-ène. On retient celle qui donne l'indice 2 à la double liaison (priorité de la liaison multiple), même si le groupe méthyle reçoit alors l'indice 4.}
\end{popfigure}

\courssub{Tableau récapitulatif et noms usuels}

\begin{center}
\begin{tabularx}{\linewidth}{|l|X|X|}
\hline
\rowcolor{popblueL} \textbf{Famille} & \textbf{Suffixe} & \textbf{Exemple} \\
\hline
Alcène (liaison double \ce{C=C}) & -ène & \ce{CH3-CH=CH-CH3} : but-2-ène \\
\hline
Alcyne (liaison triple \ce{C#C}) & -yne & \ce{CH#C-CH2-CH2-CH3} : pent-1-yne \\
\hline
\end{tabularx}
\end{center}

\begin{center}
\begin{tabularx}{\linewidth}{|l|X|X|}
\hline
\rowcolor{popblueL} \textbf{Formule} & \textbf{Nom IUPAC} & \textbf{Nom usuel} \\
\hline
\ce{CH2=CH2} & éthène & \cle{éthylène} \\
\hline
\ce{CH#CH} & éthyne & \cle{acétylène} \\
\hline
\ce{CH2=CH-CH3} & propène & \cle{propylène} \\
\hline
\end{tabularx}
\end{center}

\begin{savaistu}[Des noms usuels encore très vivants]
Dans l'industrie et le langage courant, les noms usuels résistent : on parle du \textbf{polyéthylène} (sachets, bidons) et du \textbf{polypropylène} plutôt que de polyéthène et polypropène. Quant à l'\textbf{acétylène}, c'est sous ce nom que les soudeurs camerounais le commandent : les chalumeaux oxyacétyléniques des ateliers de Douala et de Yaoundé utilisent la combustion de l'éthyne, dont la flamme atteint près de \SI{3000}{\celsius}, de quoi fondre l'acier. Au Probatoire, les deux noms sont acceptés, mais le nom IUPAC (éthène, éthyne, propène) est préféré dans une copie.
\end{savaistu}

\begin{aretenir}[L'essentiel de la nomenclature]
\begin{itemize}
\item Chaîne principale : la plus longue \textbf{contenant la liaison multiple}.
\item Numérotation : indice \textbf{le plus petit pour la liaison multiple}, prioritaire sur les substituants.
\item Suffixe : \textbf{-ène} (alcène), \textbf{-yne} (alcyne), précédé de l'indice : but-1-ène, pent-2-yne.
\item Substituants alkyles cités en préfixe, par \textbf{ordre alphabétique}, avec leurs indices.
\item Noms usuels à connaître : éthylène = éthène ; acétylène = éthyne ; propylène = propène.
\end{itemize}
\end{aretenir}

\begin{exercice}[À toi de jouer]
Nomme les composés suivants selon l'IUPAC :
\begin{enumerate}
\item \ce{CH3-CH2-CH=CH-CH3}
\item \ce{CH#C-CH2-CH2-CH3}
\item \ce{CH3-CH(CH3)-CH=CH2}
\item \ce{CH3-CH2-C#C-CH(CH3)-CH3}
\end{enumerate}
\end{exercice}

\begin{corrige}[Exercice n°1]
\begin{enumerate}
\item Chaîne de 5 C, double liaison en position 2 (et non 3) : \textbf{pent-2-ène}.
\item Chaîne de 5 C, triple liaison en position 1 : \textbf{pent-1-yne}.
\item Chaîne de 4 C contenant \ce{C=C}, double liaison en position 1 (prioritaire), méthyle en position 3 : \textbf{3-méthylbut-1-ène}.
\item Chaîne de 6 C, triple liaison en position 3 dans les deux sens ; on départage alors par le substituant : méthyle en position 2 (et non 5) : \textbf{2-méthylhex-3-yne}.
\end{enumerate}
\end{corrige}

Tu sais désormais nommer n'importe quel alcène ou alcyne. Mais certains alcènes cachent une subtilité supplémentaire : le but-2-ène, par exemple, existe sous deux formes spatiales distinctes. C'est l'objet de la prochaine section : l'\textbf{isomérie Z et E}.

\coursec{Isomérie Z et E (stéréoisomérie)}

Dans la section précédente, tu as appris à nommer les alcènes et les alcynes selon les règles de l'UICPA : choix de la chaîne principale, numérotation donnant le plus petit indice à la double liaison, suffixe \emph{-ène} ou \emph{-yne}. Tu as donc désormais un nom précis pour chaque molécule. Mais une question délicate subsiste : le nom \og but-2-ène \fg{} désigne-t-il une seule molécule ? Prends deux modèles moléculaires de \ce{CH3-CH=CH-CH3}. Tu peux placer les deux groupes méthyle \ce{CH3} du même côté de la double liaison, ou bien de part et d'autre. Et tu auras beau tourner, tordre, manipuler tes modèles : impossible de faire passer l'une de ces formes à l'autre sans casser la double liaison. Ces deux molécules ont la même formule brute, la même formule semi-développée, le même nom UICPA... et pourtant ce sont deux composés différents, aux propriétés différentes. Bienvenue dans le monde de la \cle{stéréoisomérie} Z/E, une notion classique du Probatoire que nous allons maîtriser pas à pas.

\courssub{Pourquoi la double liaison fige-t-elle la géométrie ?}

\textbf{Intuition : la règle collée}\par
Imagine deux règles plates posées l'une sur l'autre et collées ensemble : tu ne peux plus les faire pivoter l'une par rapport à l'autre sans décoller (c'est-à-dire casser) le collage. La double liaison \ce{C=C} fonctionne exactement ainsi.

Rappelons la structure vue dans la première section : dans un alcène, chaque carbone de la double liaison est \cle{trigonal} (hybridation $sp^2$). La double liaison est constituée de deux liaisons de nature différente :

\begin{itemize}
\item une liaison $\sigma$, formée par recouvrement frontal des orbitales hybrides $sp^2$ : elle autoriserait, à elle seule, la rotation ;
\item une liaison $\pi$, formée par recouvrement \emph{latéral} de deux orbitales $p$ parallèles, situées au-dessus et au-dessous du plan moléculaire.
\end{itemize}

Or, pour tourner autour de l'axe \ce{C=C}, il faudrait désaligner ces deux orbitales $p$, c'est-à-dire \textbf{rompre la liaison $\pi$}. Cela coûte une énergie considérable (environ \SI{270}{kJ.mol^{-1}}), que l'agitation thermique ordinaire ne peut pas fournir.

\begin{propriete}[Conséquence fondamentale]
La \cle{rotation} autour d'une double liaison \ce{C=C} est \textbf{bloquée} à température ordinaire. Les positions relatives des groupes fixés sur les deux carbones de la double liaison sont donc \textbf{figées} : la molécule est plane autour de la double liaison et conserve sa géométrie.
\end{propriete}

C'est exactement l'inverse de ce qui se passe dans les alcanes : autour d'une simple liaison \ce{C-C} (liaison $\sigma$ seule), la rotation est libre, et les différentes positions des groupes ne constituent pas des isomères distincts. Retiens bien ce contraste, il est souvent demandé au Probatoire.

\courssub{La condition d'existence de l'isomérie Z/E}

La rotation étant bloquée, deux dispositions spatiales différentes peuvent exister... mais pas toujours ! Tout dépend des groupes portés par les deux carbones de la double liaison.

Notons $a$ et $b$ les deux groupes portés par le premier carbone de la double liaison, et $c$ et $d$ les deux groupes portés par le second carbone :

\[ \chemfig{a-C(-[2]b)=[0]C(-[2]c)-[0]d} \]

\begin{propriete}[Condition d'existence]
Un alcène présente l'\cle{isomérie Z/E} si et seulement si \textbf{chaque} carbone de la double liaison porte \textbf{deux groupes différents} :
\[ a \neq b \quad \text{ET} \quad c \neq d \]
Si l'un des deux carbones porte deux groupes identiques (par exemple deux atomes d'hydrogène), l'isomérie Z/E \textbf{n'existe pas}.
\end{propriete}

Pourquoi ? Si le carbone de gauche porte deux groupes identiques ($a = b$), échanger ces deux groupes ne change rien à la molécule : les deux dispositions imaginées sont en réalité \emph{superposables} par une simple rotation de la molécule entière dans l'espace. Il n'y a alors qu'\textbf{une seule} molécule, pas deux isomères.

\begin{methode}[Vérifier si un alcène présente l'isomérie Z/E]
\begin{enumerate}
\item Repérer la double liaison \ce{C=C} dans la formule semi-développée.
\item Écrire explicitement les \textbf{deux groupes} portés par le \textbf{premier} carbone de la double liaison.
\item Écrire explicitement les \textbf{deux groupes} portés par le \textbf{second} carbone de la double liaison.
\item Comparer : si, sur \textbf{chaque} carbone, les deux groupes sont \textbf{différents}, alors l'isomérie Z/E existe : il y a deux stéréoisomères, l'un Z, l'autre E.
\item Si au moins un carbone porte deux groupes identiques, conclure : pas d'isomérie Z/E.
\end{enumerate}
\end{methode}

\courssub{Les descripteurs Z et E}

Lorsque la condition est remplie, il existe exactement deux stéréoisomères. Pour les distinguer, l'UICPA utilise deux lettres venues de l'allemand :

\begin{definition}[Isomères Z et E]
Deux isomères Z/E ont la \textbf{même formule brute} et la \textbf{même formule développée} (même enchaînement d'atomes), mais diffèrent par la \textbf{disposition spatiale} des groupes autour de la double liaison \ce{C=C} :
\begin{itemize}
\item l'isomère \cle{Z} (de l'allemand \emph{zusammen}, \og ensemble \fg) : les deux groupes identiques (ou prioritaires) sont situés \textbf{du même côté} de la double liaison ;
\item l'isomère \cle{E} (de l'allemand \emph{entgegen}, \og opposé \fg) : les deux groupes identiques (ou prioritaires) sont situés \textbf{de part et d'autre} de la double liaison.
\end{itemize}
Ce sont des \cle{stéréoisomères} : des isomères qui ne diffèrent que par l'arrangement des atomes dans l'espace.
\end{definition}

\begin{aretenir}[Moyen mnémotechnique — complément Probatoire]
Pour ne jamais confondre le jour de l'examen, retiens la phrase :
\begin{center}
\textbf{\og Z = Ze même côté ; E = En face \fg}
\end{center}
Z comme \emph{zusammen} (ensemble, réunis) ; E comme \emph{entgegen} (opposés, en face l'un de l'autre).
\end{aretenir}

\textbf{Remarque (pour aller un peu plus loin).}\par
Dans tous les exercices de Première, chaque carbone de la double liaison porte au moins un atome d'hydrogène, ou bien les deux carbones portent un groupe identique (\ce{CH3} par exemple) : la comparaison est alors immédiate. Dans les cas plus compliqués où les quatre groupes sont tous différents, on compare sur chaque carbone le groupe \emph{prioritaire} (celui dont l'atome directement lié a le plus grand numéro atomique) : Z si les deux groupes prioritaires sont du même côté, E sinon. Cette règle de priorité n'est pas exigible au Probatoire, mais elle te permet de comprendre la mention \og ou prioritaires \fg{} de la définition.

\courssub{Exemple type : le but-2-ène}

Le but-2-ène, de formule semi-développée \ce{CH3-CH=CH-CH3}, est l'exemple canonique à connaître par cœur. Appliquons la méthode :

\begin{itemize}
\item Premier carbone de la double liaison : il porte \ce{CH3} et \ce{H} $\rightarrow$ deux groupes \textbf{différents} ;
\item Second carbone de la double liaison : il porte \ce{CH3} et \ce{H} $\rightarrow$ deux groupes \textbf{différents}.
\end{itemize}

La condition est remplie : le but-2-ène existe sous \textbf{deux} formes stéréoisomères.

\begin{popfigure}
\begin{center}
\setchemfig{atom sep=2.2em}
\chemfig{CH_3-[0]C(-[2]H)=[0]C(-[2]H)-[0]CH_3}
\hspace{3cm}
\chemfig{CH_3-[0]C(-[2]H)=[0]C(-[-2]H)-[0]CH_3}
\end{center}
\vspace{0.3em}
\begin{center}
\textbf{(Z)-but-2-ène} : les deux \ce{CH3} sont du même côté \hspace{1.2cm} \textbf{(E)-but-2-ène} : les deux \ce{CH3} sont de part et d'autre
\end{center}
\legende{Les deux stéréoisomères du but-2-ène \ce{CH3-CH=CH-CH3}. La molécule est plane autour de la double liaison ; les groupes méthyle \ce{CH3} sont du même côté dans l'isomère Z, opposés dans l'isomère E.}
\end{popfigure}

Dans l'isomère \textbf{Z}, les deux groupes \ce{CH3} (et donc aussi les deux atomes \ce{H}) sont du même côté de la double liaison. Dans l'isomère \textbf{E}, ils sont de part et d'autre. Aucune rotation, aucune manipulation dans l'espace ne permet de superposer ces deux molécules : ce sont deux composés chimiques distincts, que l'on peut séparer et conserver séparément dans deux flacons.

Le nom complet s'écrit en ajoutant le descripteur entre parenthèses devant le nom : \textbf{(Z)-but-2-ène} et \textbf{(E)-but-2-ène}.

\begin{attention}[Erreur fréquente n°1 : le but-1-ène]
Beaucoup d'élèves attribuent à tort l'isomérie Z/E au \textbf{but-1-ène} \ce{CH2=CH-CH2-CH3}. Appliquons la méthode : le premier carbone de la double liaison est un groupe \ce{=CH2} qui porte \textbf{deux atomes d'hydrogène identiques}. La condition $a \neq b$ n'est pas respectée : échanger les deux hydrogènes ne change rien. Le but-1-ène ne présente donc \textbf{pas} d'isomérie Z/E. Au Probatoire, dès qu'une double liaison est en bout de chaîne (groupe \ce{=CH2}), la réponse est immédiate : pas de Z/E.
\end{attention}

\begin{attention}[Erreur fréquente n°2 : le 2-méthylpropène]
Le 2-méthylpropène \ce{(CH3)2C=CH2} est un autre piège classique. Ici, le premier carbone de la double liaison porte \textbf{deux groupes méthyle identiques} \ce{CH3}, et l'autre carbone porte deux hydrogènes identiques. Doublement, la condition n'est pas remplie : pas d'isomérie Z/E. Moralité : vérifie \textbf{les deux} carbones ; un seul carbone défaillant suffit à faire échouer la condition.
\end{attention}

\begin{popfigure}
\begin{center}
\begin{tikzpicture}[
  font=\small,
  decision/.style={diamond, draw=popblue, fill=popblueL, text width=4.6cm, align=center, inner sep=1pt, aspect=2.2},
  result/.style={rectangle, rounded corners, draw=popgreen, fill=popgreenL, text width=4.2cm, align=center, inner sep=4pt},
  nope/.style={rectangle, rounded corners, draw=poporange, fill=poporangeL, text width=4.2cm, align=center, inner sep=4pt},
  start/.style={rectangle, rounded corners, draw=popdark, fill=popcream, text width=5.5cm, align=center, inner sep=4pt},
  arr/.style={-{Stealth[length=2.5mm]}, thick, popdark}
]
\node[start] (s) {Alcène de formule connue : repérer la double liaison \ce{C=C}};
\node[decision, below=1.1cm of s] (d1) {Le premier carbone de \ce{C=C} porte-t-il deux groupes différents ?};
\node[nope, right=2.6cm of d1] (n1) {Pas d'isomérie Z/E\\ (ex. \ce{=CH2} du but-1-ène)};
\node[decision, below=1.3cm of d1] (d2) {Le second carbone de \ce{C=C} porte-t-il deux groupes différents ?};
\node[nope, right=2.6cm of d2] (n2) {Pas d'isomérie Z/E\\ (ex. \ce{(CH3)2C=} du 2-méthylpropène)};
\node[result, below=1.3cm of d2] (ok) {L'isomérie Z/E existe :\\ un isomère \textbf{Z} (même côté)\\ et un isomère \textbf{E} (en face)};
\draw[arr] (s) -- (d1);
\draw[arr] (d1) -- node[above, font=\footnotesize\itshape]{non} (n1);
\draw[arr] (d1) -- node[right, font=\footnotesize\itshape]{oui} (d2);
\draw[arr] (d2) -- node[above, font=\footnotesize\itshape]{non} (n2);
\draw[arr] (d2) -- node[right, font=\footnotesize\itshape]{oui} (ok);
\end{tikzpicture}
\end{center}
\legende{Arbre de décision : un alcène donné présente-t-il l'isomérie Z/E ? Il faut que \emph{chaque} carbone de la double liaison porte deux groupes différents.}
\end{popfigure}

\courssub{Z et E sont des stéréoisomères : des molécules vraiment différentes}

Insistons sur un point fondamental souvent mal compris : les isomères Z et E ne sont pas de simples \og dessins \fg{} différents d'une même molécule. Ce sont des \textbf{espèces chimiques distinctes}, que l'on peut isoler, peser, analyser séparément.

\begin{definition}[Stéréoisomères]
Des \cle{stéréoisomères} sont des composés qui ont la même formule brute et le même enchaînement d'atomes (même formule développée), mais qui diffèrent par l'\textbf{arrangement des atomes dans l'espace} (configuration différente) et qui ne sont \textbf{pas superposables}.
\end{definition}

Comme toutes les molécules différentes, les isomères Z et E possèdent des \textbf{propriétés physiques différentes} : températures de fusion et d'ébullition, densité, indice de réfraction... Leurs propriétés chimiques sont en revanche très voisines, puisqu'ils possèdent les mêmes groupes fonctionnels.

\begin{savaistu}[La géométrie change les propriétés physiques]
Le \textbf{(E)-but-2-ène} bout à $\theta_{eb} = \SI{0,9}{\celsius}$, tandis que le \textbf{(Z)-but-2-ène} bout à $\theta_{eb} = \SI{3,7}{\celsius}$ : près de \SI{3}{\celsius} d'écart pour deux molécules de même formule brute \ce{C4H8} ! L'explication : la molécule E, plus \og symétrique \fg, possède un moment dipolaire global quasi nul (les effets des deux groupes \ce{CH3} se compensent), alors que la molécule Z est légèrement polaire. Les interactions entre molécules Z sont donc un peu plus fortes, et il faut davantage d'énergie pour les vaporiser. Ce phénomène est partout autour de nous : dans la rétine, la vision déclenche la conversion d'un isomère Z en isomère E (le rétinal), et l'industrie alimentaire s'inquiète de la géométrie Z ou E des graisses insaturées (les fameuses graisses \emph{trans}).
\end{savaistu}

\courssub{Application : exercice résolu type Probatoire}

\begin{exoresolu}[Stéréoisomérie du pent-2-ène]
On considère le pent-2-ène, alcène de formule brute \ce{C5H10}.
\begin{enumerate}
\item Écrire sa formule semi-développée.
\item Montrer qu'il présente l'isomérie Z/E.
\item Représenter les deux stéréoisomères et les nommer.
\item Le pent-1-ène, isomère de position du pent-2-ène, présente-t-il l'isomérie Z/E ? Justifier.
\end{enumerate}
\tcblower
\textbf{Solution détaillée}

\textbf{1.} Le pent-2-ène possède une chaîne de 5 carbones avec la double liaison entre les carbones 2 et 3 :
\[ \ce{CH3-CH=CH-CH2-CH3} \]

\textbf{2.} Examinons les groupes portés par chaque carbone de la double liaison :
\begin{itemize}
\item carbone 2 : il porte \ce{CH3} et \ce{H} $\rightarrow$ deux groupes différents ;
\item carbone 3 : il porte \ce{CH2-CH3} (éthyle) et \ce{H} $\rightarrow$ deux groupes différents.
\end{itemize}
Chaque carbone de la double liaison porte deux groupes différents : \textbf{l'isomérie Z/E existe}.

\textbf{3.} Les deux stéréoisomères (les atomes d'hydrogène de la double liaison servent de référence, car chaque carbone n'en porte qu'un seul) :
\begin{center}
\setchemfig{atom sep=2.2em}
\chemfig{CH_3-[0]C(-[2]H)=[0]C(-[2]H)-[0]CH_2-CH_3}
\hspace{2.2cm}
\chemfig{CH_3-[0]C(-[2]H)=[0]C(-[-2]H)-[0]CH_2-CH_3}
\end{center}
\begin{itemize}
\item À gauche : les deux hydrogènes (et les deux groupes alkyles) sont du même côté : \textbf{(Z)-pent-2-ène} ;
\item À droite : ils sont de part et d'autre : \textbf{(E)-pent-2-ène}.
\end{itemize}

\textbf{4.} Le pent-1-ène a pour formule \ce{CH2=CH-CH2-CH2-CH3}. Le carbone 1 de la double liaison est un groupe \ce{=CH2} : il porte \textbf{deux atomes d'hydrogène identiques}. La condition d'existence n'est pas remplie : le pent-1-ène \textbf{ne présente pas} d'isomérie Z/E.
\end{exoresolu}

\begin{exercice}[À toi de jouer]
Pour chacun des alcènes suivants, dire s'il présente l'isomérie Z/E ; si oui, représenter et nommer les deux stéréoisomères :
\begin{enumerate}
\item \ce{CH3-CH=CH-CH2-CH2-CH3} (hex-2-ène) ;
\item \ce{CH3-C(CH3)=CH-CH3} (2-méthylbut-2-ène) ;
\item \ce{CH3-CH2-CH=CH-CH2-CH3} (hex-3-ène) ;
\item \ce{CH2=CH-CH2-CH3} (but-1-ène).
\end{enumerate}
\end{exercice}

\begin{corrige}[Exercice]
\textbf{1.} Hex-2-ène : le carbone 2 porte \ce{CH3} et \ce{H} ; le carbone 3 porte \ce{CH2CH2CH3} et \ce{H}. Deux groupes différents sur chaque carbone : \textbf{l'isomérie Z/E existe}. On distingue le \textbf{(Z)-hex-2-ène} (les deux \ce{H} du même côté) et le \textbf{(E)-hex-2-ène} (les deux \ce{H} de part et d'autre).

\textbf{2.} 2-Méthylbut-2-ène : le carbone 2 porte \textbf{deux groupes} \ce{CH3} \textbf{identiques} : \textbf{pas d'isomérie Z/E}.

\textbf{3.} Hex-3-ène : le carbone 3 porte \ce{CH2CH3} et \ce{H} ; le carbone 4 porte \ce{CH2CH3} et \ce{H}. Deux groupes différents sur chaque carbone : \textbf{l'isomérie Z/E existe}. On distingue le \textbf{(Z)-hex-3-ène} (les deux groupes éthyle du même côté) et le \textbf{(E)-hex-3-ène} (les deux groupes éthyle de part et d'autre).

\textbf{4.} But-1-ène : le carbone 1 porte deux hydrogènes identiques (\ce{=CH2}) : \textbf{pas d'isomérie Z/E}.
\end{corrige}

\courssub{Bilan de la section}

\begin{aretenir}[L'essentiel sur l'isomérie Z/E]
\begin{itemize}
\item La double liaison \ce{C=C} \textbf{bloque la rotation} (liaison $\pi$) : la géométrie autour d'elle est figée et plane.
\item \textbf{Condition d'existence} : chaque carbone de la double liaison doit porter \textbf{deux groupes différents}.
\item Isomère \textbf{Z} (\emph{zusammen}) : groupes identiques du \textbf{même côté} ; isomère \textbf{E} (\emph{entgegen}) : groupes identiques \textbf{en face}. Mémo : \og Z = Ze même côté ; E = En face \fg.
\item Z et E sont des \textbf{stéréoisomères} : même formule brute, même formule développée, configurations spatiales différentes, non superposables, propriétés physiques différentes (ex. ébullition : \SI{3,7}{\celsius} pour le Z contre \SI{0,9}{\celsius} pour le E du but-2-ène).
\item Pièges : toute double liaison en bout de chaîne (\ce{=CH2}, comme le but-1-ène) ou tout carbone portant deux groupes identiques (2-méthylpropène) exclut l'isomérie Z/E.
\end{itemize}
\end{aretenir}

Maintenant que tu sais distinguer et nommer ces stéréoisomères, il est temps d'étudier comment les alcènes et les alcynes \emph{réagissent} : la section suivante te fera découvrir leurs propriétés chimiques, de la combustion aux fameuses réactions d'addition qui font toute la richesse de la chimie des insaturations.

\coursec{Propriétés chimiques : combustion et réactions d'addition}

Dans la section précédente, nous avons découvert que la double liaison \ce{C=C} impose une géométrie plane et rigide, source de l'isomérie $Z$/$E$. Mais cette double liaison n'est pas seulement une contrainte géométrique : c'est aussi une \cle{zone de fragilité} de la molécule. Riche en électrons (liaison $\pi$), elle attire les réactifs électrophiles et peut s'ouvrir pour accueillir de nouveaux atomes. C'est toute la chimie des alcènes et des alcynes qui va se jouer ici : brûler, ou se laisser « additionner ». Tu comprendras dans cette section pourquoi l'eau de brome se décolore au contact d'un gaz insaturé, pourquoi le soudeur du quartier utilise de l'acétylène, et comment on fabrique industriellement l'éthanol.

\courssub{La combustion des alcènes et des alcynes}

\textbf{Intuition.} Comme tous les hydrocarbures, alcènes et alcynes brûlent dans le dioxygène. Mais observe la flamme d'une lampe à acétylène mal réglée ou celle d'un feu de plastique : elle est jaune, très éclairante, et dégage une fumée noire. Pourquoi ? Parce que ces molécules sont très \cle{riches en carbone} (rapport C/H élevé) : il n'y a pas assez d'oxygène au contact de la flamme pour tout brûler, et des particules de carbone incandescentes (la \emph{suie}) se forment.

\begin{propriete}[Combustion complète des alcènes]
En présence d'un excès de dioxygène, la combustion d'un alcène de formule générale $C_nH_{2n}$ est \cle{complète} : elle produit uniquement du dioxyde de carbone et de l'eau, avec dégagement de chaleur (réaction \emph{exothermique}) :
\[ \ce{C_{n}H_{2n} + \frac{3n}{2} O2 -> n CO2 + n H2O} \]
Si le dioxygène est en défaut, la combustion est \emph{incomplète} : la flamme devient fuligineuse (dépôt de carbone \ce{C}) et peut dégager du monoxyde de carbone \ce{CO}, gaz mortel.
\end{propriete}

\begin{propriete}[Combustion complète des alcynes]
Pour un alcyne de formule générale $C_nH_{2n-2}$, l'équation-bilan de la combustion complète s'écrit :
\[ \ce{C_{n}H_{2n-2} + \frac{3n-1}{2} O2 -> n CO2 + (n-1) H2O} \]
La flamme est encore plus fuligineuse que pour les alcènes (teneur en carbone plus grande).
\end{propriete}

\begin{exemplebox}[Équilibrer la combustion du propène \ce{C3H6}]
Cherchons les coefficients de l'équation : $\ce{C3H6} + a\,\ce{O2} \rightarrow b\,\ce{CO2} + c\,\ce{H2O}$.
\begin{itemize}
\item Conservation du carbone : $b = 3$.
\item Conservation de l'hydrogène : $2c = 6$, donc $c = 3$.
\item Conservation de l'oxygène : $2a = 2b + c = 6 + 3 = 9$, donc $a = \dfrac{9}{2}$.
\end{itemize}
D'où l'équation :
\[ \ce{C3H6 + 9/2 O2 -> 3 CO2 + 3 H2O} \]
Pour n'avoir que des coefficients entiers, on multiplie tout par 2 :
\[ \ce{2 C3H6 + 9 O2 -> 6 CO2 + 6 H2O} \]
Vérification : 6 C, 12 H et 18 O de chaque côté. L'équation est bien équilibrée.
\end{exemplebox}

\begin{attention}[Le piège du coefficient fractionnaire]
Écrire $\frac{3n}{2}$ ou $\frac{9}{2}$ devant \ce{O2} n'est pas une erreur : c'est une écriture correcte pour une équation-bilan rapportée à \textbf{une} mole d'hydrocarbure. En revanche, si l'énoncé exige des coefficients entiers, multiplie tous les coefficients par 2. Erreur fréquente : oublier de multiplier aussi le coefficient de l'hydrocarbure !
\end{attention}

\textbf{Cas particulier : la combustion de l'acétylène.} L'éthyne \ce{C2H2} brûle selon :
\[ \ce{2 C2H2 + 5 O2 -> 4 CO2 + 2 H2O} \]
Cette réaction, très exothermique, est le cœur du \cle{chalumeau oxyacétylénique} : en brûlant l'acétylène dans un courant de dioxygène pur (et non dans l'air), on obtient une flamme bleue, très compacte, dont la température atteint environ \SI{3000}{\celsius}.

\begin{savaistu}[La flamme oxyacétylénique au Cameroun]
À environ \SI{3000}{\celsius}, la flamme oxyacétylénique fond l'acier (qui fond vers \SI{1500}{\celsius}) sans difficulté. C'est elle que manipulent les soudeurs de Douala, Yaoundé ou Bafoussam pour assembler les portails métalliques, réparer les châssis de camions ou découper la tôle. L'acétylène est stocké dans des bouteilles dissous dans de l'acétone, car sous pression pure, il peut exploser : un bel exemple de la chimie au service de l'artisanat !
\end{savaistu}

\courssub{Le principe général des réactions d'addition}

\textbf{Intuition.} Imagine la double liaison \ce{C=C} comme une fermeture éclair à double rangée. Une molécule réactive $X-Y$ peut « casser » une des deux rangées (la liaison $\pi$) : la double liaison s'ouvre, devient simple, et chacun des deux carbones capte un fragment du réactif ($X$ d'un côté, $Y$ de l'autre). Rien ne se perd, rien ne s'élimine : tout le réactif est \emph{ajouté} à la molécule. D'où le nom d'\cle{addition}.

\begin{definition}[Réaction d'addition]
Une \cle{réaction d'addition} est une réaction au cours de laquelle une molécule se fixe sur les deux atomes de carbone d'une liaison multiple : la double liaison \ce{C=C} devient une liaison simple \ce{C-C}, ou la triple liaison \ce{C#C} devient double puis simple. L'équation générale pour un alcène s'écrit :
\[ \ce{C=C + X-Y -> X-C-C-Y} \]
L'addition est la \textbf{réaction caractéristique} des composés insaturés (alcènes et alcynes), par opposition à la \emph{substitution}, réaction caractéristique des alcanes.
\end{definition}

\begin{popfigure}
\begin{center}
\begin{tikzpicture}[scale=1.0]
\node at (0,0) {\chemfig{CH_2=CH_2}};
\node at (2.2,0) {$+$};
\node at (3.6,0) {\chemfig{Br-Br}};
\draw[-{Stealth[length=3mm]},very thick,poporange] (4.6,0) -- (6.2,0);
\node at (5.4,0.55) {\small addition};
\node at (7.8,0) {\chemfig{Br-CH_2-CH_2-Br}};
\draw[-{Stealth[length=2mm]},thick,popblue] (0.55,-0.75) to[bend right=20] (1.15,-0.75);
\node[popblue] at (0.85,-1.25) {\small la liaison $\pi$ s'ouvre};
\end{tikzpicture}
\end{center}
\legende{Mécanisme simplifié de l'addition du dibrome sur l'éthylène : la liaison $\pi$ s'ouvre et un atome de brome se fixe sur chaque carbone.}
\end{popfigure}

\courssub{L'hydrogénation : de l'alcène à l'alcane}

\begin{propriete}[Hydrogénation catalytique]
En présence d'un catalyseur métallique (nickel de Raney, platine ou palladium) et à chaud, le dihydrogène s'additionne sur la double liaison. L'alcène se transforme en \cle{alcane} :
\[ \ce{C_{n}H_{2n} + H2 ->[\text{Ni, Pt ou Pd, chauffage}] C_{n}H_{2n+2}} \]
Exemple avec l'éthylène :
\[ \ce{CH2=CH2 + H2 ->[Ni, \Delta] CH3-CH3} \]
Le catalyseur n'apparaît pas dans le bilan : il accélère la réaction sans être consommé.
\end{propriete}

Cette réaction a une application industrielle considérable : l'\cle{hydrogénation des huiles végétales} (huile de palme, huile d'arachide) transforme les chaînes insaturées en chaînes saturées, ce qui durcit l'huile en graisse solide — c'est le principe de fabrication des margarines.

\begin{attention}[Ne pas confondre les deux sens du test]
L'hydrogénation transforme un alcène (incolore) en alcane (incolore) : elle n'offre \textbf{aucun changement visible}. Ce n'est donc pas un test d'identification. Le test d'identification des insaturations, c'est la décoloration de l'eau de brome (voir plus bas). Ne dis jamais « on identifie un alcène par hydrogénation » !
\end{attention}

\courssub{L'addition des dihalogènes : le test à l'eau de brome}

\begin{propriete}[Addition des dihalogènes]
Le dichlore et le dibrome s'additionnent spontanément, à température ordinaire, sur la double liaison d'un alcène, sans catalyseur ni chauffage. On obtient un \cle{dihalogénoalcane} vicinal (les deux halogènes sur deux carbones voisins) :
\[ \ce{CH2=CH2 + Br2 -> CH2Br-CH2Br} \]
Le produit est le \textbf{1,2-dibromoéthane}, liquide incolore. De même : $\ce{CH2=CH2 + Cl2 -> CH2Cl-CH2Cl}$ (1,2-dichloroéthane).
\end{propriete}

Or le dibrome dissous dans l'eau (\emph{eau de brome}) ou dans un solvant organique est \textbf{orange}, alors que le 1,2-dibromoéthane formé est \textbf{incolore}. Conséquence spectaculaire : si on barbote un alcène dans l'eau de brome, la couleur orange disparaît. C'est le test officiel d'identification des insaturations.

\begin{methode}[Test d'identification d'une insaturation (eau de brome)]
\begin{enumerate}
\item Verser environ \SI{2}{mL} d'eau de brome (solution \textbf{orange}) dans deux tubes à essai.
\item Introduire dans le premier tube le composé à tester, dans le second un alcane (témoin).
\item Agiter et observer.
\item \textbf{Conclusion} : si l'eau de brome se \cle{décolore}, le composé est insaturé (alcène ou alcyne) ; si la couleur orange persiste, c'est un alcane (composé saturé).
\end{enumerate}
Remarque : ce test ne permet \emph{pas} de distinguer un alcène d'un alcyne, car les deux décolorent l'eau de brome (l'alcyne fixe deux molécules de \ce{Br2}).
\end{methode}

\begin{popfigure}
\begin{center}
\begin{tikzpicture}[scale=0.9]
% Tube 1 : avant
\draw[thick] (-4.2,0) -- (-4.2,-2.6) arc[start angle=180,end angle=360,radius=0.45] -- (-3.3,0);
\fill[poporange!70] (-4.2,-1.1) -- (-4.2,-2.6) arc[start angle=180,end angle=360,radius=0.45] -- (-3.3,-1.1) -- cycle;
\draw[thick] (-4.2,-1.1) -- (-3.3,-1.1);
\node[below] at (-3.75,-3.15) {\small eau de brome};
\node[below,poporange] at (-3.75,-3.6) {\small (orange)};
% flèche
\draw[-{Stealth[length=3mm]},very thick,popblue] (-2.6,-1.5) -- (-1.2,-1.5);
\node[above,popblue] at (-1.9,-1.4) {\small $+$ alcène};
% Tube 2 : après
\draw[thick] (-0.4,0) -- (-0.4,-2.6) arc[start angle=180,end angle=360,radius=0.45] -- (0.5,0);
\fill[popcream] (-0.4,-1.1) -- (-0.4,-2.6) arc[start angle=180,end angle=360,radius=0.45] -- (0.5,-1.1) -- cycle;
\draw[thick] (-0.4,-1.1) -- (0.5,-1.1);
\node[below] at (0.05,-3.15) {\small solution};
\node[below,popmuted] at (0.05,-3.6) {\small (incolore)};
% bulles gaz
\foreach \y in {0.2,0.5,0.8} {\draw[popblue!60] (-2.0,\y-0.2) circle (0.06);}
\node[popblue] at (-2.0,1.0) {\small gaz insaturé};
\end{tikzpicture}
\end{center}
\legende{Test à l'eau de brome : la solution orange se décolore au contact d'un alcène ou d'un alcyne.}
\end{popfigure}

\begin{experience}[Mise en évidence au laboratoire]
\textbf{Protocole} : on fait barboter de l'éthylène (obtenu par déshydratation de l'éthanol) dans un tube contenant de l'eau de brome. \textbf{Observation} : la solution orange se décolore progressivement et il apparaît un liquide huileux incolore, plus dense que l'eau. \textbf{Interprétation} : le dibrome s'est additionné sur la double liaison ; le liquide huileux est le 1,2-dibromoéthane. Ce test est aussi utilisé pour suivre la maturation des fruits, qui dégagent naturellement de l'éthylène.
\end{experience}

\courssub{L'addition des hydracides et la règle de Markovnikov}

\textbf{Cas d'un alcène symétrique.} Avec l'éthylène, les deux carbones de la double liaison sont identiques : peu importe de quel côté se fixe l'hydrogène, on obtient un seul produit, le chloroéthane :
\[ \ce{CH2=CH2 + HCl -> CH3-CH2Cl} \]

\textbf{Cas d'un alcène asymétrique.} Avec le propène \ce{CH2=CH-CH3}, les deux carbones de la double liaison ne jouent pas le même rôle : l'un porte deux atomes d'hydrogène (carbone \emph{le plus hydrogéné}), l'autre n'en porte qu'un. Deux produits sont alors envisageables : le 2-chloropropane et le 1-chloropropane. L'expérience montre qu'on obtient \textbf{presque exclusivement} le 2-chloropropane. Le chimiste russe Vladimir Markovnikov a résumé cette observation en 1869.

\begin{definition}[Règle de Markovnikov]
Lors de l'addition d'un hydracide \ce{HX} (ou de l'eau) sur un alcène \cle{asymétrique}, l'atome d'hydrogène se fixe préférentiellement sur l'atome de carbone de la double liaison \textbf{le plus hydrogéné} (celui qui porte déjà le plus d'atomes d'hydrogène). L'atome d'halogène (ou le groupe \ce{OH}) se fixe donc sur le carbone le \emph{moins} hydrogéné, c'est-à-dire le plus substitué.
\end{definition}

\begin{savaistu}[Pourquoi la règle de Markovnikov ? (Niveau approfondi)]
La réaction se passe en deux étapes via un intermédiaire : le \emph{carbocation}. L'hydrogène (\ce{H+}) attaque la double liaison. Si le \ce{H+} se fixe sur le carbone le plus hydrogéné, le carbocation se forme sur le carbone le plus substitué (secondaire), qui est plus stable que le carbocation primaire qui se formerait sinon. La stabilité des carbocations dicte la régiosélectivité : tertiaire $>$ secondaire $>$ primaire.
\end{savaistu}

\begin{popfigure}
\begin{center}
\begin{tikzpicture}[scale=0.95]
\node at (-3.4,0) {\chemfig{CH_2=CH-CH_3}};
\node at (-1.7,0) {$+$};
\node at (-0.7,0) {\chemfig{H-Cl}};
% flèche majoritaire
\draw[-{Stealth[length=3mm]},ultra thick,popgreen] (0.4,0.35) to[bend left=15] (2.6,0.9);
\node[popgreen] at (1.5,1.35) {\small\textbf{majoritaire}};
\node at (4.6,0.9) {\chemfig{CH_3-CH(-[2]Cl)-CH_3}};
\node[below,popgreen] at (4.6,0.15) {\small 2-chloropropane};
% flèche minoritaire
\draw[-{Stealth[length=3mm]},thick,popmuted] (0.4,-0.35) to[bend right=15] (2.6,-0.9);
\node[popmuted] at (1.5,-1.35) {\small minoritaire};
\node at (4.7,-0.9) {\chemfig{CH_2(-[2]Cl)-CH_2-CH_3}};
\node[below,popmuted] at (4.7,-1.55) {\small 1-chloropropane};
\end{tikzpicture}
\end{center}
\legende{Addition de \ce{HCl} sur le propène : conformément à la règle de Markovnikov, l'hydrogène se fixe sur le carbone le plus hydrogéné (\ce{CH2}) ; le 2-chloropropane est le produit majoritaire.}
\end{popfigure}

\begin{exemplebox}[Application de la règle de Markovnikov au propène]
Dans \ce{CH2=CH-CH3}, le carbone 1 (\ce{CH2}) porte \textbf{deux} atomes d'hydrogène, le carbone 2 (\ce{CH}) n'en porte qu'\textbf{un}. D'après Markovnikov, le H de \ce{HCl} se fixe sur le carbone 1, et le Cl sur le carbone 2 :
\[ \ce{CH2=CH-CH3 + HCl -> CH3-CHCl-CH3} \]
Le produit majoritaire est le \textbf{2-chloropropane}. Le 1-chloropropane \ce{CH2Cl-CH2-CH3} ne se forme qu'en très faible quantité.
\end{exemplebox}

\begin{attention}[Erreur fréquente au Probatoire]
Beaucoup d'élèves écrivent le \textbf{1-chloropropane} comme produit principal de l'addition de \ce{HCl} sur le propène, par simple « lecture de gauche à droite ». C'est faux ! Applique systématiquement la règle de Markovnikov : repère le carbone de la double liaison qui porte le plus d'atomes d'hydrogène, et donne-lui le H du réactif. Le produit majoritaire est toujours l'halogénoalcane le plus \emph{substitué} (halogène sur le carbone le plus entouré de groupes alkyles).
\end{attention}

\courssub{L'hydratation : des alcènes aux alcools}

\begin{propriete}[Hydratation des alcènes]
En présence d'un catalyseur acide (acide phosphorique \ce{H3PO4} sur support en phase gazeuse industriellement, ou acide sulfurique \ce{H2SO4} en phase liquide), vers \SI{300}{\celsius} sous pression, l'eau s'additionne sur la double liaison : on obtient un \cle{alcool}. La règle de Markovnikov s'applique : le groupe \ce{OH} se fixe sur le carbone le moins hydrogéné.
\[ \ce{CH2=CH2 + H2O ->[H3PO4, 300\,^\circ C, P] CH3-CH2OH} \]
L'éthylène conduit à l'\textbf{éthanol}. C'est une voie industrielle majeure de production d'alcool éthylique (solvants, désinfectants, carburants).
\end{propriete}

\begin{exoresolu}[Hydratation du propène]
\textbf{Énoncé.} On fait réagir le propène \ce{CH2=CH-CH3} avec l'eau en présence d'acide phosphorique à \SI{300}{\celsius}. Écrire l'équation de la réaction et nommer le produit majoritaire obtenu.
\tcblower
\textbf{Solution.} L'eau se décompose formellement en H et \ce{OH}. Le propène est asymétrique : le carbone 1 (\ce{CH2}) porte deux atomes d'hydrogène, le carbone 2 (\ce{CH}) un seul. D'après la règle de Markovnikov, H se fixe sur le carbone 1 et \ce{OH} sur le carbone 2 :
\[ \ce{CH2=CH-CH3 + H2O ->[H3PO4] CH3-CHOH-CH3} \]
Le produit majoritaire est le \textbf{propan-2-ol} (ou isopropanol), alcool secondaire bien connu des pharmacies comme alcool antiseptique. Le propan-1-ol ne se forme qu'en traces.
\end{exoresolu}

\courssub{Les additions sur les alcynes}

Les alcynes, avec leur triple liaison \ce{C#C} (une liaison $\sigma$ et deux liaisons $\pi$), sont encore plus insaturés que les alcènes : ils peuvent fixer \textbf{deux molécules} de réactif, en deux étapes successives. L'hydrogénation de l'acétylène l'illustre parfaitement :

\textbf{Première étape} (une mole de \ce{H2}, catalyseur de Lindlar \ce{Pd/CaCO3} empoisonné) : la triple liaison devient double, on obtient l'éthylène (stéréochimie \emph{cis}) :
\[ \ce{CH#CH + H2 ->[Pd/Lindlar] CH2=CH2} \]

\textbf{Deuxième étape} (une deuxième mole de \ce{H2}, nickel de Raney) : la double liaison devient simple, on obtient l'éthane :
\[ \ce{CH2=CH2 + H2 ->[Ni] CH3-CH3} \]

Bilan global : $\ce{CH#CH + 2 H2 -> CH3-CH3}$. Avec un catalyseur « désactivé » (palladium de Lindlar), on peut arrêter la réaction à l'étape de l'alcène \emph{cis} ; avec du nickel, on va jusqu'à l'alcane.

De même, l'acétylène décolore l'eau de brome en fixant deux molécules de \ce{Br2} pour donner le 1,1,2,2-tétrabromoéthane \ce{CHBr2-CHBr2} (incolore). L'addition d'hydracides (\ce{HCl}, \ce{HBr}) suit aussi la règle de Markovnikov à chaque étape, conduisant aux dihalogénures geminaux (deux halogènes sur le même carbone) comme produits finaux majoritaires.

\begin{aretenir}[L'essentiel de la section]
\begin{itemize}
\item Combustion complète alcène : $\ce{C_{n}H_{2n} + \frac{3n}{2} O2 -> n CO2 + n H2O}$ ; alcynes : $\ce{C_{n}H_{2n-2} + \frac{3n-1}{2} O2 -> n CO2 + (n-1) H2O}$. Flamme éclairante et fuligineuse (richesse en carbone). Acétylène : $\ce{2 C2H2 + 5 O2 -> 4 CO2 + 2 H2O}$, flamme à \SI{3000}{\celsius} (soudure).
\item L'\cle{addition} est la réaction caractéristique des insaturations : la liaison multiple s'ouvre et fixe le réactif sur les deux carbones.
\item Hydrogénation (Ni, Pt, Pd, chauffage) : alcène $\rightarrow$ alcane ; alcyne $\rightarrow$ alcène (Lindlar) $\rightarrow$ alcane.
\item Eau de brome : \textbf{décoloration = insaturation} (alcène ou alcyne) ; l'alcane ne décolore pas.
\item Règle de \cle{Markovnikov} : sur un alcène asymétrique, H se fixe sur le carbone le plus hydrogéné ; propène + \ce{HCl} $\rightarrow$ 2-chloropropane (majoritaire). Valable pour \ce{HBr}, \ce{HI} et \ce{H2O}.
\item Hydratation acide (\ce{H3PO4}, \SI{300}{\celsius}, pression) : alcène + eau $\rightarrow$ alcool (éthylène $\rightarrow$ éthanol ; propène $\rightarrow$ propan-2-ol).
\item Les alcynes additionnent en \textbf{deux étapes} : alcyne $\rightarrow$ alcène $\rightarrow$ alcane (ou dihalogénure).
\end{itemize}
\end{aretenir}

\begin{exercice}[Pour t'entraîner]
\begin{enumerate}
\item Équilibrer l'équation de la combustion complète du but-1-ène \ce{C4H8}.
\item On fait barboter un hydrocarbure gazeux dans l'eau de brome : la solution se décolore. Que peut-on conclure ? Le gaz peut-il être l'éthane ?
\item Écrire l'équation de l'addition de \ce{HCl} sur le but-1-ène \ce{CH2=CH-CH2-CH3} et nommer le produit majoritaire.
\item Quel(s) alcène(s) faut-il hydrater pour obtenir majoritairement le butan-2-ol ?
\end{enumerate}
\end{exercice}

\begin{corrige}[Exercice n°1]
\begin{enumerate}
\item $\ce{C4H8 + 6 O2 -> 4 CO2 + 4 H2O}$ (vérification : 4 C, 8 H, 12 O de chaque côté).
\item Le gaz est insaturé (alcène ou alcyne). Ce ne peut pas être l'éthane, alcane saturé qui ne décolore pas l'eau de brome.
\item D'après Markovnikov, H se fixe sur \ce{CH2} (carbone le plus hydrogéné) : $\ce{CH2=CH-CH2-CH3 + HCl -> CH3-CHCl-CH2-CH3}$, le \textbf{2-chlorobutane} (majoritaire).
\item Le butan-2-ol \ce{CH3-CHOH-CH2-CH3} s'obtient par hydratation du \textbf{but-1-ène} (le \ce{OH} se fixe sur le carbone 2, le moins hydrogéné) \textbf{ou du but-2-ène} (alcène symétrique, un seul produit possible : butan-2-ol).
\end{enumerate}
\end{corrige}

Maintenant que tu maîtrises les réactions qui consomment la double liaison, il reste à découvrir comment des milliers de molécules d'alcène peuvent s'additionner \emph{entre elles} pour former des géants : les polymères. C'est l'objet de la dernière section, qui te mènera aussi au laboratoire pour préparer l'acétylène.

\coursec{Polymérisation et préparation de l'acétylène}

Dans la section précédente, nous avons vu que la double liaison \ce{C=C} est une véritable \og porte ouverte \fg{} : elle peut s'ouvrir pour accueillir du dihydrogène, de l'eau, du chlorure d'hydrogène ou des dihalogènes. Mais que se passe-t-il si, au lieu d'ajouter une petite molécule étrangère, on force les molécules d'alcène à s'ajouter \emph{les unes aux autres} ? Chaque double liaison s'ouvre et les molécules s'enchaînent comme les maillons d'une longue chaîne : c'est la \cle{polymérisation}, réaction qui a donné naissance à tous les plastiques qui nous entourent, des sachets du marché Mokolo aux tuyaux PVC des chantiers de Douala. Nous terminerons cette leçon par la préparation au laboratoire d'un alcyne célèbre : l'\cle{acétylène}, le gaz des soudeurs.

\courssub{La polymérisation : des monomères aux polymères}

\textbf{Intuition.} Prends un collier de perles : chaque perle est identique, et c'est leur enchaînement répété des centaines de fois qui fait le collier. En chimie, la perle s'appelle le \cle{monomère} (du grec \emph{monos}, seul, et \emph{meros}, partie) et le collier s'appelle le \cle{polymère} (\emph{polus}, nombreux).

\begin{definition}[Polymère et polymérisation]
Un \cle{polymère} est une \cle{macromolécule} (molécule géante, de masse molaire pouvant dépasser \SI{100000}{g.mol^{-1}}) obtenue par enchaînement covalent d'un très grand nombre de \cle{motifs} identiques issus d'une petite molécule appelée \cle{monomère}. La réaction qui transforme le monomère en polymère s'appelle la \cle{polymérisation}. Lorsque la polymérisation se fait par ouverture des doubles liaisons sans perte d'atomes, on parle de \cle{polyaddition}.
\end{definition}

\textbf{Le polyéthylène (PE) : le plastique le plus simple.} Sous l'action de la température, de la pression et d'un catalyseur, les molécules d'éthylène \ce{CH2=CH2} ouvrent leur double liaison et s'enchaînent :

\[ n\,\ce{CH2=CH2} \xrightarrow{T,\ P,\ \text{catalyseur}} \ce{-(CH2-CH2)_{n}-} \]

Le groupe \ce{-CH2-CH2-} est le \cle{motif} du polymère ; l'entier $n$ est l'\cle{indice de polymérisation} (ou degré de polymérisation) : il indique combien de motifs sont enchaînés. Il peut valoir plusieurs milliers !

\begin{popfigure}
\begin{center}
$n$ \chemfig{CH_2=[1]CH_2} \hspace{0,6em} \raisebox{0,9em}{$\xrightarrow{T,\ P,\ \text{catalyseur}}$} \hspace{0,6em}
\chemfig{-[@{a}]CH_2-CH_2-[@{b}]}
\chemmove{\draw[popink,thick] ([yshift=6pt]a) -- ([yshift=6pt]b) node[midway,above,popink]{$n$ fois};}
\end{center}
\legende{Polyaddition de l'éthylène : $n$ molécules de monomère donnent une chaîne de $n$ motifs \ce{-CH2-CH2-} : le polyéthylène.}
\end{popfigure}

\begin{methode}[Écrire une équation de polymérisation]
\begin{enumerate}
\item Écris la formule semi-développée du monomère en mettant bien en évidence la double liaison \ce{C=C}.
\item \og Ouvre \fg{} la double liaison : chaque carbone garde une valence libre de chaque côté.
\item Écris le motif obtenu entre parenthèses, avec des tirets de part et d'autre, et l'indice $n$ en bas à droite : \ce{-(...)_n-}.
\item Place le coefficient $n$ devant le monomère et indique les conditions ($T$, $P$, catalyseur) sur la flèche.
\end{enumerate}
\end{methode}

\begin{exemplebox}[D'autres polymères d'addition]
\begin{itemize}
\item \textbf{Propène} \ce{CH2=CH-CH3} $\longrightarrow$ \textbf{polypropylène} (PP), de motif \ce{-(CH2-CH(CH3))-} : chaises et tables plastiques, cuvettes, seaux.
\item \textbf{Chlorure de vinyle} \ce{CH2=CHCl} $\longrightarrow$ \textbf{polychlorure de vinyle} (PVC), de motif \ce{-(CH2-CHCl)-} : tuyaux de plomberie, gaines électriques, gouttières.
\end{itemize}
\[ n\,\ce{CH2=CHCl} \xrightarrow{T,\ P,\ \text{catalyseur}} \ce{-(CH2-CHCl)_{n}-} \]
\end{exemplebox}

\begin{savaistu}[Les plastiques dans la vie camerounaise]
Les sachets d'emballage de l'eau minérale, les jerricans jaunes d'huile de palme, les chaises plastiques des cérémonies, les tuyaux PVC des adductions d'eau : tous sont des polymères d'addition issus d'alcènes. Leur succès vient de leur légèreté, de leur insolubilité dans l'eau et de leur faible coût. En contrepartie, ils sont très peu biodégradables : leur accumulation dans les rues et les rivières de Douala et Yaoundé pose un grave problème environnemental, d'où l'importance du tri et du recyclage.
\end{savaistu}

\begin{attention}[Pièges fréquents]
\begin{itemize}
\item Le motif du polyéthylène est \ce{-CH2-CH2-} et non \ce{CH2=CH2} : dans le polymère, il n'y a \textbf{plus de double liaison} !
\item Ne pas confondre polymérisation et polycondensation : ici, aucune petite molécule (eau, HCl) n'est éliminée.
\item L'indice $n$ n'est pas un coefficient stœchiométrique précis : il est très grand et variable d'une chaîne à l'autre.
\end{itemize}
\end{attention}

\courssub{Préparation de l'acétylène au laboratoire}

\textbf{Principe.} L'acétylène \ce{C2H2} (éthyne) s'obtient facilement par action de l'eau sur le \cle{carbure de calcium} \ce{CaC2}, solide grisâtre. La réaction est \cle{vive} et \cle{exothermique} (elle dégage de la chaleur) :

\[ \ce{CaC2(s) + 2 H2O(l) -> C2H2(g) + Ca(OH)2(s)} \]

Le second produit est l'\cle{hydroxyde de calcium} \ce{Ca(OH)2}, appelé \emph{chaux éteinte}, qui forme un dépôt blanchâtre.

\begin{attention}[Erreur classique au Probatoire]
Beaucoup d'élèves écrivent \ce{CaC2 + H2O -> C2H2} en oubliant le \ce{Ca(OH)2} ! Vérifie toujours la conservation des atomes : le calcium et les deux groupes hydroxyle doivent apparaître quelque part. L'équation correcte est \ce{CaC2 + 2 H2O -> C2H2 + Ca(OH)2}.
\end{attention}

\textbf{Dispositif expérimental.} On place le carbure de calcium dans un flacon ; on verse l'eau goutte à goutte à l'aide d'une ampoule de coulée (pour contrôler la vitesse de la réaction) ; le gaz dégagé est conduit par un tube à dégagement jusqu'à une éprouvette renversée, remplie d'eau, placée dans un cristallisoir : c'est le \cle{recueil par déplacement d'eau}.

\begin{popfigure}
\begin{center}
\begin{tikzpicture}[scale=0.95, every node/.style={font=\small}]
% Flacon de réaction
\draw[thick,popdark] (0,0) -- (0,2.2) -- (0.5,2.7) -- (0.5,3.4);
\draw[thick,popdark] (2.4,0) -- (2.4,2.2) -- (1.9,2.7) -- (1.9,3.4);
\draw[thick,popdark] (0,0) .. controls (0.6,-0.35) and (1.8,-0.35) .. (2.4,0);
% Carbure (morceaux)
\fill[poppurple!60] (0.5,0.15) circle (0.12);
\fill[poppurple!60] (1.0,0.1) circle (0.15);
\fill[poppurple!60] (1.6,0.18) circle (0.11);
\fill[poppurple!60] (2.0,0.08) circle (0.13);
\node[popmuted] at (1.2,-0.7) {carbure de calcium \ce{CaC2}};
% Bouchon
\draw[thick,popgold,fill=popgoldL] (0.5,3.35) rectangle (1.9,3.6);
% Ampoule de coulée
\draw[thick,popdark] (0.95,3.6) -- (0.95,4.1);
\draw[thick,popdark] (1.45,3.6) -- (1.45,4.1);
\draw[thick,popdark,fill=popblue!20] (0.95,4.1) .. controls (0.8,4.9) and (0.85,5.4) .. (1.2,5.5)
  .. controls (1.55,5.4) and (1.6,4.9) .. (1.45,4.1) -- cycle;
\draw[thick,popdark] (1.2,5.5) -- (1.2,5.8);
\draw[thick,popdark] (1.1,5.8) -- (1.3,5.8);
% robinet
\draw[thick,popink] (1.0,3.85) -- (1.4,3.85);
\draw[thick,popink] (1.2,3.75) -- (1.2,3.95);
\node[popblue,align=center] at (-1.6,4.9) {ampoule de coulée\\ (eau \ce{H2O})};
\draw[-{Stealth},popblue] (-0.6,4.7) -- (0.8,4.7);
% Tube à dégagement
\draw[thick,popdark] (1.65,3.6) -- (1.65,4.3) -- (4.6,4.3) -- (4.6,1.6);
\draw[thick,popdark] (1.85,3.6) -- (1.85,4.1) -- (4.8,4.1) -- (4.8,1.6);
\node[popmuted] at (3.2,4.75) {tube à dégagement};
% Cristallisoir
\draw[thick,popdark] (3.2,-0.6) -- (3.2,1.4) -- (6.4,1.4) -- (6.4,-0.6) -- cycle;
\fill[popblue!15] (3.25,-0.55) rectangle (6.35,1.1);
\node[popblue] at (6.9,0.3) {eau};
% Éprouvette renversée
\draw[thick,popdark] (4.3,1.35) -- (4.3,0.2) -- (5.1,0.2) -- (5.1,1.35);
\fill[white] (4.35,0.85) rectangle (5.05,1.3);
\fill[popblue!15] (4.35,0.25) rectangle (5.05,0.85);
\node[popmuted,align=center] at (4.7,2.0) {éprouvette renversée\\ (acétylène recueilli)};
\draw[-{Stealth},popmuted] (4.7,1.75) -- (4.7,1.45);
% bulles
\fill[popblue!50] (4.7,0.5) circle (0.05);
\fill[popblue!50] (4.6,0.7) circle (0.04);
\fill[popblue!50] (4.8,0.65) circle (0.045);
\end{tikzpicture}
\end{center}
\legende{Préparation de l'acétylène au laboratoire : l'eau tombe goutte à goutte sur le carbure de calcium ; le gaz \ce{C2H2} est recueilli par déplacement d'eau.}
\end{popfigure}

\textbf{Pourquoi recueillir l'acétylène sur l'eau ?} Parce que l'acétylène est \cle{peu soluble dans l'eau} : il chasse l'eau de l'éprouvette sans s'y dissoudre de façon notable. C'est une justification classique demandée à l'examen.

\textbf{Propriétés du gaz recueilli.} L'acétylène est un gaz \cle{incolore}, \cle{très inflammable} : il brûle dans l'air avec une flamme très chaude (environ \SI{3000}{\celsius} dans le dioxygène pur), ce qui explique son usage dans les chalumeaux oxyacétyléniques des soudeurs. Comme tout alcyne, il \cle{décolore l'eau de brome} (addition de \ce{Br2} sur la triple liaison) : c'est le test qui prouve son insaturation.

\begin{attention}[Sécurité]
L'acétylène forme avec l'air un \cle{mélange détonant}. Avant d'allumer quoi que ce soit, il faut chasser l'air du dispositif et \textbf{éloigner toute flamme} du dégagement gazeux. Ne jamais boucher le tube à dégagement pendant la réaction : la pression monterait dangereusement dans le flacon.
\end{attention}

\begin{exoresolu}[Masse de carbure de calcium nécessaire]
On veut préparer $V = \SI{2,4}{L}$ d'acétylène, mesurés dans des conditions où le volume molaire vaut $V_m = \SI{24}{L.mol^{-1}}$. Quelle masse de carbure de calcium pur faut-il utiliser ? On donne $M(\ce{CaC2}) = \SI{64}{g.mol^{-1}}$.
\tcblower
\textbf{Étape 1 : quantité d'acétylène.}
\[ n(\ce{C2H2}) = \frac{V}{V_m} = \frac{2,4}{24} = \SI{0,10}{mol} \]

\textbf{Étape 2 : tableau d'avancement.}
\begin{center}
\begin{tabular}{|c|c|c|c|c|}\hline
 & \ce{CaC2(s)} & \ce{2 H2O(l)} & \ce{C2H2(g)} & \ce{Ca(OH)2(s)} \\\hline
État initial & $x$ & excès & 0 & 0 \\\hline
État final & $x-0,10$ & excès & \SI{0,10}{mol} & \SI{0,10}{mol} \\\hline
\end{tabular}
\end{center}
D'après l'équation, 1 mol \ce{CaC2} donne 1 mol \ce{C2H2}, donc $n(\ce{CaC2}) = \SI{0,10}{mol}$.

\textbf{Étape 3 : masse de carbure.}
\[ m(\ce{CaC2}) = n \times M = 0,10 \times 64 = \boxed{\SI{6,4}{g}} \]
Il faut donc \SI{6,4}{g} de carbure de calcium pur (en pratique un peu plus, car le carbure technique contient des impuretés).
\end{exoresolu}

\begin{savaistu}[Le carbure et le mûrissement des fruits : une pratique dangereuse]
Sur certains marchés du Cameroun, on trouve du carbure de calcium utilisé pour faire mûrir rapidement bananes, mangues et avocats : au contact de l'humidité de l'air, il dégage de l'acétylène, qui accélère le mûrissement (l'éthylène naturel des fruits joue le même rôle). Mais le carbure \emph{technique} contient des impuretés toxiques (dérivés du phosphore et de l'arsenic) qui contaminent les fruits. Cet usage alimentaire est dangereux pour la santé et déconseillé : préfère le mûrissement naturel, par exemple en enfermant les fruits avec des bananes déjà mûres qui dégagent de l'éthylène.
\end{savaistu}

\begin{exercice}[Pour s'entraîner]
\begin{enumerate}
\item Écris l'équation de polymérisation du propène \ce{CH2=CH-CH3} et entoure le motif du polymère obtenu.
\item Sur le schéma du dispositif de préparation de l'acétylène, nomme : l'ampoule de coulée, le tube à dégagement, le cristallisoir, l'éprouvette. Justifie le mode de recueil du gaz.
\item Quel volume d'acétylène ($V_m = \SI{24}{L.mol^{-1}}$) obtient-on à partir de \SI{12,8}{g} de carbure de calcium pur ?
\end{enumerate}
\end{exercice}

\begin{corrige}[Exercice]
\begin{enumerate}
\item $n\,\ce{CH2=CH-CH3} \xrightarrow{T,P,\text{catalyseur}} \ce{-(CH2-CH(CH3))_{n}-}$ ; le motif est \ce{-CH2-CH(CH3)-}.
\item L'ampoule de coulée verse l'eau goutte à goutte sur le carbure ; le tube à dégagement conduit le gaz jusqu'à l'éprouvette renversée sur le cristallisoir d'eau. Le recueil se fait par déplacement d'eau car l'acétylène est peu soluble dans l'eau.
\item $n(\ce{CaC2}) = \dfrac{12,8}{64} = \SI{0,20}{mol}$. 
\begin{center}
\begin{tabular}{|c|c|c|c|c|}\hline
 & \ce{CaC2} & \ce{2 H2O} & \ce{C2H2} & \ce{Ca(OH)2} \\\hline
État final & 0 & excès & \SI{0,20}{mol} & \SI{0,20}{mol} \\\hline
\end{tabular}
\end{center}
$n(\ce{C2H2}) = \SI{0,20}{mol}$, d'où $V = n \times V_m = 0,20 \times 24 = \boxed{\SI{4,8}{L}}$.
\end{enumerate}
\end{corrige}

\begin{aretenir}[L'essentiel de la section]
\begin{itemize}
\item La \cle{polymérisation} enchaîne un grand nombre de monomères en une macromolécule ; le motif se répète $n$ fois ($n$ : indice de polymérisation).
\item $n\,\ce{CH2=CH2} \rightarrow \ce{-(CH2-CH2)_{n}-}$ : le polyéthylène ; de même le propène donne le PP et le chlorure de vinyle le PVC.
\item L'acétylène se prépare par action de l'eau sur le carbure de calcium : $\ce{CaC2 + 2 H2O -> C2H2 + Ca(OH)2}$ (réaction vive et exothermique).
\item Le gaz est recueilli par déplacement d'eau car il est peu soluble ; il est incolore, très inflammable et décolore l'eau de brome.
\item Sécurité : mélange détonant avec l'air — aucune flamme à proximité du dégagement gazeux.
\end{itemize}
\end{aretenir}

\newpage

\coursec{Travaux pratiques}

\courssub{Objectifs du TP}

À l'issue de cette séance de travaux pratiques, tu dois être capable de :

\begin{itemize}
\item préparer l'\cle{acétylène} (éthyne, \ce{C2H2}) par action de l'eau sur le carbure de calcium ;
\item recueillir un gaz insoluble dans l'eau par \cle{déplacement d'eau} ;
\item mettre en évidence expérimentalement les \cle{insaturations} (liaisons multiples) d'un hydrocarbure par le test à l'\cle{eau de brome} ;
\item observer la \cle{combustion} de l'acétylène et interpréter l'aspect de sa flamme ;
\item écrire et équilibrer les équations-bilans des trois réactions observées ;
\item appliquer les règles de sécurité liées à la manipulation d'un gaz inflammable formant avec l'air un \cle{mélange détonant}.
\end{itemize}

Ce TP illustre directement la famille de situations du programme : \emph{utilisation de l'acétylène}. Rappelle-toi que le carbure de calcium, vendu dans les quincailleries et sur les marchés de Yaoundé ou de Douala, a longtemps servi dans les lampes à acétylène des mineurs et des pêcheurs, et que le chalumeau oxyacétylénique des soudeurs de nos quartiers utilise ce même gaz.

\courssub{Matériel et produits}

\begin{center}
\rowcolors{2}{popblueL}{white}
\begin{tabularx}{\linewidth}{|l|X|l|}
\hline
\rowcolor{popblue!40} \textbf{Matériel / produit} & \textbf{Rôle} & \textbf{Quantité} \\
\hline
Carbure de calcium \ce{CaC2} en morceaux & Réactif solide (source d'acétylène) & \SI{5}{g} environ \\
\hline
Eau distillée (ou eau de pluie filtrée) & Réactif liquide & \SI{50}{mL} \\
\hline
Flacon de Woulfe (ou flacon à col large + bouchon percé à 2 trous) & Réacteur (lieu de la réaction) & 1 \\
\hline
Ampoule de coulée (ou entonnoir + robinet, ou pipette compte-gouttes à travers le bouchon) & Introduire l'eau goutte à goutte & 1 \\
\hline
Tube à dégagement coudé en verre + tuyau souple & Conduire le gaz & 1 \\
\hline
Cristallisoir (ou bassine, ou cuvette en plastique) & Cuve à eau pour le recueil & 1 \\
\hline
Éprouvette graduée de \SI{250}{mL} (ou bouteille plastique graduée renversée) & Recueillir et mesurer le gaz & 2 \\
\hline
Tubes à essais + portoir & Tests d'identification & 3 \\
\hline
Eau de brome (solution orange de \ce{Br2}) & Réactif des insaturations & \SI{5}{mL} \\
\hline
Bec Bunsen (ou bougie, ou allumettes) & Test de combustion & 1 \\
\hline
Blouse, lunettes de protection, gants & Sécurité individuelle & 1 jeu \\
\hline
\end{tabularx}
\end{center}

\textbf{Alternatives locales.} Si l'eau de brome manque, on peut la remplacer par une solution de permanganate de potassium \ce{KMnO4} acidifiée (violette), qui se décolore aussi en présence d'insaturations ; le permanganate se trouve en pharmacie. À défaut d'éprouvette graduée, une bouteille en plastique transparente, préalablement jaugée avec un verre mesureur, convient parfaitement. Le carbure de calcium s'achète couramment en quincaillerie ; il doit être conservé dans un pot \emph{hermétiquement fermé}, car il réagit avec l'humidité de l'air.

\courssub{Sécurité}

\begin{attention}[Consignes de sécurité impératives]
\begin{itemize}
\item \textbf{Pictogramme « gaz inflammable »} (losange rouge avec flamme) : l'acétylène est un gaz \emph{extrêmement inflammable}. Aucune flamme ne doit approcher le dispositif pendant la préparation.
\item \textbf{Danger du mélange détonant} : le mélange acétylène + air \emph{explose} à l'allumage. On n'enflamme le gaz à la sortie du tube \textbf{qu'après avoir chassé tout l'air} du dispositif, c'est-à-dire après avoir recueilli et rejeté au moins une première éprouvette de gaz, ou après avoir vérifié qu'un échantillon prélevé brûle calmement.
\item \textbf{Pictogramme « corrosif »} : le carbure de calcium et la chaux \ce{Ca(OH)2} formée sont irritants pour la peau et les yeux. Port de la \textbf{blouse} et des \textbf{lunettes} obligatoires ; ne jamais toucher le carbure à mains nues.
\item \textbf{Eau de brome} : toxique et corrosive, elle tache la peau en brun. Manipuler avec des gants, sous le hotteau ou près d'une fenêtre ouverte ; ne jamais la porter à la bouche.
\item Verser l'eau \textbf{goutte à goutte} : la réaction est vive et exothermique ; un ajout brutal provoque un dégagement incontrôlable de gaz et une forte mousse.
\item La flamme de l'acétylène est très éclairante et fuligineuse : ne pas approcher le visage du tube à dégagement lors de l'allumage ; utiliser une allumette longue ou une allumette fixée à une pince.
\item En fin de séance, les résidus de chaux éteinte sont dilués dans beaucoup d'eau avant d'être jetés ; le carbure non utilisé est rendu au professeur, jamais jeté dans un évier humide.
\end{itemize}
\end{attention}

\courssub{Schéma du dispositif}

\begin{popfigure}
\begin{center}
\begin{tikzpicture}[scale=0.95, every node/.style={font=\small}]
% Cristallisoir
\draw[thick, popblue] (-1.2,0) -- (-1.2,1.6) (5.2,0) -- (5.2,1.6);
\draw[thick, popblue] (-1.2,0) -- (5.2,0);
\draw[popblue!50, fill=popblue!15] (-1.15,0.05) rectangle (5.15,1.3);
\node[popblue!70!black] at (4.3,0.35) {eau};
% Éprouvette renversée
\draw[thick, popdark] (1.4,1.15) -- (1.4,3.4) -- (2.6,3.4) -- (2.6,1.15);
\draw[popblue!40, fill=popblue!25] (1.45,1.2) rectangle (2.55,1.9);
\draw[fill=white] (1.45,1.95) rectangle (2.55,3.35);
\node[align=center] at (2.0,2.7) {gaz\\ \ce{C2H2}};
\foreach \y in {2.2,2.5,2.8,3.1} \draw[popmuted] (2.35,\y) -- (2.55,\y);
\node[align=center, popdark] at (2.0,4.0) {éprouvette renversée};
% Flacon de Woulfe
\draw[thick, popdark] (-4.6,2.2) .. controls (-4.9,2.6) and (-4.9,4.2) .. (-4.4,4.6)
   -- (-4.4,5.2) (-3.4,5.2) -- (-3.4,4.6)
   .. controls (-2.9,4.2) and (-2.9,2.6) .. (-3.2,2.2) -- cycle;
\draw[thick, popdark] (-4.6,2.2) -- (-3.2,2.2);
\draw[popgold!60, fill=popgold!40] (-4.55,2.3) rectangle (-3.25,2.9);
\node[popgold!50!black] at (-3.9,2.6) {\ce{CaC2}};
% Bouchon
\draw[fill=popmuted!40] (-4.5,5.2) rectangle (-3.3,5.5);
% Ampoule de coulée
\draw[thick, popdark] (-4.2,5.5) -- (-4.2,6.0);
\draw[thick, popdark] (-4.45,6.0) -- (-4.45,7.3) -- (-3.95,7.3) -- (-3.95,6.0);
\draw[thick, popdark] (-4.35,6.0) -- (-4.05,6.0);
\draw[popblue!50, fill=popblue!20] (-4.42,6.35) rectangle (-3.98,7.25);
\node[popblue!70!black] at (-4.2,6.8) {\footnotesize eau};
\draw[thick, popdark] (-4.2,6.0) -- (-4.2,4.0);
\draw[thick] (-4.32,5.75) -- (-4.08,5.75);
\node[right, align=left] at (-3.9,7.5) {ampoule de coulée\\ (robinet)};
% Tube à dégagement
\draw[thick, popdark] (-3.55,5.5) -- (-3.55,6.3) -- (2.0,6.3) -- (2.0,1.6);
\draw[thick, popdark] (-3.45,5.5) -- (-3.45,6.4) -- (2.1,6.4) -- (2.1,1.6);
\node[above] at (-0.8,6.4) {tube à dégagement};
% Bulles
\foreach \p in {(1.8,1.5),(2.2,1.7),(2.0,1.4)} \draw[popblue] \p circle (0.06);
% Légendes flèches
\draw[->, thick, popink] (-3.9,3.4) -- (-3.9,3.0);
\node[left, popink] at (-4.6,3.2) {réaction vive};
\end{tikzpicture}
\end{center}
\legende{Dispositif de préparation et de recueil de l'acétylène : l'eau coule goutte à goutte sur le carbure de calcium ; le gaz formé est recueilli par déplacement d'eau dans l'éprouvette renversée.}
\end{popfigure}

\courssub{Protocole}

\begin{enumerate}
\item \textbf{Montage.} Placer environ \SI{5}{g} de carbure de calcium en morceaux (3 à 4 morceaux de la taille d'un haricot) au fond du flacon de Woulfe. Adapter le bouchon percé muni de l'ampoule de coulée (robinet fermé) et du tube à dégagement. Vérifier l'étanchéité du montage.
\item \textbf{Mise en place du recueil.} Remplir le cristallisoir d'eau aux trois quarts. Remplir complètement l'éprouvette graduée d'eau, la boucher avec le pouce, la renverser sur la cuve et la fixer, ouverture sous l'eau. Introduire l'extrémité du tube à dégagement sous l'ouverture de l'éprouvette.
\item \textbf{Chasse de l'air (étape critique).} Verser \SI{50}{mL} d'eau dans l'ampoule de coulée. Ouvrir délicatement le robinet pour laisser tomber l'eau \textbf{goutte à goutte} (environ 1 goutte par seconde) sur le carbure. Le dégagement gazeux est immédiat. \textbf{Rejeter la première éprouvette de gaz} (elle contient surtout l'air du dispositif) : la vider à l'air libre, loin de toute flamme, puis la remplir d'eau et la remettre en place.
\item \textbf{Recueil.} Recueillir l'acétylène dans la deuxième éprouvette jusqu'à obtention d'environ \SI{200}{mL} de gaz. Refermer le robinet. Boucher l'éprouvette sous l'eau, la sortir et la conserver renversée sur la paillasse.
\item \textbf{Test à l'eau de brome.} Introduire \SI{2}{mL} d'eau de brome dans un tube à essais. À l'aide d'un tube en verre relié au dispositif (ou en débouchant l'éprouvette sous l'eau et en y plongeant le tube), faire barboter lentement l'acétylène dans l'eau de brome pendant 30 secondes. Observer et noter l'évolution de la coloration.
\item \textbf{Test de combustion.} L'air du dispositif ayant été chassé (étape 3), approcher une flamme de l'extrémité du tube à dégagement pendant que le gaz s'écoule doucement. Observer l'aspect de la flamme et déposer éventuellement une soucoupe froide au-dessus de la flamme pour recueillir le dépôt.
\item \textbf{Fin de séance.} Fermer le robinet, démonter, diluer le contenu du flacon (bouillie blanche de chaux) dans un grand volume d'eau, laver la verrerie.
\end{enumerate}

\courssub{Observations et résultats attendus}

\begin{center}
\rowcolors{2}{popgreenL}{white}
\begin{tabularx}{\linewidth}{|X|X|}
\hline
\rowcolor{popgreen!45} \textbf{Étape} & \textbf{Observation attendue} \\
\hline
Action de l'eau sur \ce{CaC2} & Dégagement gazeux immédiat et régulier, bouillonnement, échauffement du flacon, formation d'une bouillie blanchâtre (chaux) ; légère odeur caractéristique. \\
\hline
Recueil du gaz & Le niveau d'eau descend dans l'éprouvette ; environ \SI{200}{mL} de gaz incolore recueillis en 4 à 6 minutes pour \SI{5}{g} de carbure commercial. \\
\hline
Test à l'eau de brome & La solution orange se \textbf{décolore} progressivement jusqu'à devenir incolore. \\
\hline
Test de combustion & Le gaz s'enflamme avec une flamme \textbf{très éclairante, jaune et fuligineuse}, dégageant une fumée noire ; dépôt noir de carbone sur la soucoupe froide. \\
\hline
\end{tabularx}
\end{center}

\courssub{Exploitation}

\begin{exoresolu}[Questions d'exploitation du TP]
\textbf{1.} Écrire l'équation-bilan de la réaction de préparation de l'acétylène. De quel type de réaction s'agit-il ?

\textbf{2.} Pourquoi peut-on recueillir l'acétylène par déplacement d'eau ? Pourquoi rejeter la première éprouvette ?

\textbf{3.} Écrire l'équation-bilan de la réaction entre l'acétylène et le dibrome. Que prouve la décoloration ?

\textbf{4.} Écrire l'équation-bilan de la combustion complète de l'acétylène. Comment expliquer que la flamme soit fuligineuse ?

\textbf{5.} Calculer le volume théorique d'acétylène dégagé par \SI{5,0}{g} de carbure de calcium pur, dans les conditions du laboratoire ($V_m = \SI{24}{L.mol^{-1}}$). On donne $M(\ce{CaC2}) = \SI{64}{g.mol^{-1}}$. Comparer au volume recueilli et commenter.

\tcblower
\textbf{1.} L'eau réagit sur le carbure de calcium selon :
\[ \ce{CaC2(s) + 2 H2O(l) -> C2H2(g) + Ca(OH)2(s)} \]
C'est une réaction d'\cle{hydrolyse} (réaction acido-basique au sens large), vive et exothermique. Le résidu blanc est l'hydroxyde de calcium (chaux éteinte).

\textbf{2.} L'acétylène est \emph{très peu soluble dans l'eau} : il peut donc chasser l'eau de l'éprouvette sans s'y dissoudre de façon notable. La première éprouvette est rejetée car elle contient le \textbf{mélange air + acétylène} initialement présent dans le flacon et les tubes : ce mélange est \emph{détonant} et ne doit jamais être enflammé.

\textbf{3.} L'acétylène fixe le dibrome par \cle{addition} sur la triple liaison (jusqu'à saturation) :
\[ \ce{HC#CH + 2 Br2 -> Br2CH-CHBr2} \]
La décoloration de l'eau de brome prouve la présence d'\textbf{insaturations} (liaisons multiples) : le dibrome orange est consommé et le produit formé (1,1,2,2-tétrabromoéthane) est incolore. C'est le test caractéristique des alcènes et des alcynes.

\textbf{4.} Combustion complète dans le dioxygène :
\[ \ce{2 C2H2(g) + 5 O2(g) -> 4 CO2(g) + 2 H2O(g)} \]
L'acétylène est l'hydrocarbure le plus riche en carbone ($92,3\,\%$ en masse). Dans l'air, l'apport de dioxygène est insuffisant : la combustion est \emph{incomplète} et produit des particules de \textbf{carbone} incandescentes, d'où la flamme jaune très éclairante et la fumée noire. Dans le chalumeau oxyacétylénique, on apporte du dioxygène pur : la combustion devient complète et la flamme, bleue, atteint environ \SI{3000}{\celsius}.

\textbf{5.} Quantité de carbure : $n(\ce{CaC2}) = \dfrac{m}{M} = \dfrac{5,0}{64} = \SI{7,8e-2}{mol}$.

D'après l'équation, $n(\ce{C2H2}) = n(\ce{CaC2})$, donc :
\[ V(\ce{C2H2}) = n \times V_m = 7,8\times 10^{-2} \times 24 \approx \SI{1,9}{L} \]
Le volume réellement recueilli ($\approx \SI{200}{mL}$) est bien inférieur : le carbure commercial est \emph{impur} (il contient de la chaux et des sulfures, d'où l'odeur), on n'a utilisé qu'une partie de l'eau, et une fraction du gaz s'échappe lors des manipulations. Le rendement expérimental est de l'ordre de $10\,\%$ dans les conditions d'un TP de lycée, ce qui reste largement suffisant pour les tests.
\end{exoresolu}

\courssub{Conclusion}

Ce TP t'a permis de préparer l'\cle{acétylène} par hydrolyse du carbure de calcium, de le recueillir par déplacement d'eau, puis de caractériser ses insaturations par deux tests complémentaires : la \textbf{décoloration de l'eau de brome}, qui prouve l'existence de liaisons multiples capables de fixer le dibrome par addition, et la \textbf{combustion à flamme fuligineuse}, signature d'un hydrocarbure très riche en carbone. Tu as également appliqué la règle de sécurité fondamentale : ne jamais enflammer un gaz avant d'avoir chassé l'air du dispositif. Ces savoir-faire expérimentaux concrétisent les propriétés chimiques des hydrocarbures insaturés étudiées dans la leçon et expliquent l'usage industriel de l'acétylène, du chalumeau des soudeurs camerounais à la synthèse des matières plastiques.

\begin{savaistu}[Le carbure et la maturation des fruits]
Au Cameroun comme dans de nombreux pays, certains commerçants utilisent du carbure de calcium pour faire mûrir rapidement bananes et mangues : l'acétylène dégagé imite l'\emph{éthylène}, l'hormone végétale naturelle de maturation. Cette pratique est \textbf{dangereuse et interdite} : le carbure technique contient des impuretés toxiques (arsine, phosphine) qui contaminent les fruits. La chimie des alcènes et des alcynes n'est donc pas qu'une affaire de laboratoire : elle touche directement à la santé publique !
\end{savaistu}

\newpage

\coursec{Méthodes et exercices résolus}

\courssub{Méthodes et exercices résolus}

\begin{methode}[Déterminer la géométrie et l'hybridation des atomes de carbone]
\textbf{Objectif :} Prédire la géométrie locale (linéaire, trigonale plane, tétraédrique) et l'hybridation (\textit{sp}, \textit{sp}$^2$, \textit{sp}$^3$) de chaque carbone d'un alcène ou d'un alcyne.

\begin{enumerate}
    \item \textbf{Compter les domaines électroniques} (liaisons $\sigma$ + doublets non liants) autour du carbone étudié.
    \item \textbf{Déduire l'hybridation} :
        \begin{itemize}
            \item 2 domaines $\rightarrow$ \textit{sp} (digonal, angle 180$^\circ$) \quad \textit{Ex : Carbone d'un triple lien C$\equiv$C}.
            \item 3 domaines $\rightarrow$ \textit{sp}$^2$ (trigonal plan, angle $\approx$ 120$^\circ$) \quad \textit{Ex : Carbone d'un double lien C=C}.
            \item 4 domaines $\rightarrow$ \textit{sp}$^3$ (tétraédrique, angle $\approx$ 109,5$^\circ$) \quad \textit{Ex : Carbone saturé C$-$C}.
        \end{itemize}
    \item \textbf{Représenter en Cram} : 
        \begin{itemize}
            \item \textit{sp}$^2$ : 3 liaisons $\sigma$ dans le plan (triangle), 1 liaison $\pi$ perpendiculaire (non représentée en Cram standard, le nuage $\pi$ est au-dessus/au-dessous du plan).
            \item \textit{sp} : 2 liaisons $\sigma$ colinéaires (axe), 2 liaisons $\pi$ perpendiculaires entre elles et à l'axe.
        \end{itemize}
\end{enumerate}
\textbf{Réflexe :} Un carbone d'alcène a \textbf{3 liaisons $\sigma$} (1 $\sigma$ C=C + 2 $\sigma$ C$-$H ou C$-$C) $\rightarrow$ \textit{sp}$^2$. Un carbone d'alcyne a \textbf{2 liaisons $\sigma$} (1 $\sigma$ C$\equiv$C + 1 $\sigma$ C$-$H ou C$-$C) $\rightarrow$ \textit{sp}.
\end{methode}

\begin{exoresolu}[Géométrie du but-1-en-3-yne (CH$\equiv$C$-$CH=CH$_2$)]
\textbf{Énoncé :} On considère la molécule de \textbf{but-1-en-3-yne} de formule semi-développée \ce{HC#C-CH=CH2}.
\begin{enumerate}
    \item Déterminer l'hybridation de chacun des quatre atomes de carbone.
    \item Préciser la géométrie au voisinage de chaque carbone et les angles de liaison approximatifs.
    \item Faire une représentation de Cram de la molécule (on ne représente que les liaisons $\sigma$).
\end{enumerate}

\tcblower
\textbf{Solution détaillée}

\begin{enumerate}
    \item \textbf{Analyse de la structure :} La chaîne carbonée comporte 4 atomes de carbone : C$_1$$\equiv$C$_2$$-$C$_3$=C$_4$H$_2$ avec H sur C$_1$ et H sur C$_3$.
    \begin{itemize}
        \item \textbf{C$_1$ (CH$\equiv$) :} 2 liaisons $\sigma$ (1 C$\equiv$C $\sigma$ + 1 C$-$H) $\rightarrow$ \textbf{\textit{sp}}.
        \item \textbf{C$_2$ ($\equiv$C$-$) :} 2 liaisons $\sigma$ (1 C$\equiv$C $\sigma$ + 1 C$-$C simple) $\rightarrow$ \textbf{\textit{sp}}.
        \item \textbf{C$_3$ ($-$CH=) :} 3 liaisons $\sigma$ (1 C$-$C simple + 1 C=C $\sigma$ + 1 C$-$H) $\rightarrow$ \textbf{\textit{sp}$^2$}.
        \item \textbf{C$_4$ (=CH$_2$) :} 3 liaisons $\sigma$ (1 C=C $\sigma$ + 2 C$-$H) $\rightarrow$ \textbf{\textit{sp}$^2$}.
    \end{itemize}

    \item \textbf{Géométries et angles :}
    \begin{center}
    \begin{tabularx}{\linewidth}{|c|c|c|X|}
    \hline
    \rowcolor{popblueL}
    \textbf{Carbone} & \textbf{Hybridation} & \textbf{Géométrie} & \textbf{Angles} \\ \hline
    C$_1$, C$_2$ & \textit{sp} & Linéaire (digonale) & 180$^\circ$ \\ \hline
    C$_3$, C$_4$ & \textit{sp}$^2$ & Trigonal plan & $\approx$ 120$^\circ$ \\ \hline
    \end{tabularx}
    \end{center}

    \item \textbf{Représentation de Cram (liaisons $\sigma$ seulement) :}
    \begin{popfigure}
\begin{tikzpicture}[scale=1.2, every node/.style={transform shape}]
        % Coordinates
        \coordinate (C1) at (0,0);
        \coordinate (C2) at (1.5,0);
        \coordinate (C3) at (3,0);
        % C4 attached to C3 at 60 deg (trigonal planar geometry)
        \coordinate (C4) at ($(C3) + (60:1.5)$);
        % H on C1 (linear, opposite to C2)
        \coordinate (H_C1) at ($(C1) + (180:1)$);
        % H on C3 (linear sp, opposite to C2-C3 axis? No, sp has 180 deg. C2-C3-H is 180. So H is at 0 deg from C3? Wait.)
        % C2-C3 is along x-axis. C3 is sp. Sigma bonds: C2-C3 and C3-C4 and C3-H? No, sp carbon has 2 sigma bonds.
        % Structure: H-C1#C2-C3=C4H2. C1 is sp (H-C#C). C2 is sp (C#C-C). C3 is sp2 (C=C). C4 is sp2.
        % Correction: C1#C2-C3=C4.
        % C1: bonded to H and C2 (triple). Sigma: C1-H, C1-C2.
        % C2: bonded to C1 (triple) and C3 (single). Sigma: C2-C1, C2-C3.
        % C3: bonded to C2 (single), C4 (double), H. Sigma: C3-C2, C3-C4, C3-H. Trigonal planar ~120 deg.
        % C2-C3 bond is along x-axis (0 deg). C3-C4 sigma at +120 deg (or 60 deg). C3-H at -120 deg (or 240 deg).
        % Let's set C3-C4 at +120 deg (pointing up-left) or +60 deg (up-right).
        % Original code had C4 at 60 deg from C3. So C3-C4 axis is 60 deg.
        % Then C3-H should be at 60 - 120 = -60 deg (or 60 + 120 = 180 deg).
        % Original code used -60 deg for H_C3. This is consistent.
        \coordinate (H_C3) at ($(C3) + (-60:1)$);
        % H on C4: C4 is sp2. Sigma bonds: C3-C4 (axis 60 deg), C4-Ha, C4-Hb. Angles 120 deg apart.
        % Ha at 60 + 120 = 180 deg. Hb at 60 - 120 = -60 deg.
        \coordinate (H_C4a) at ($(C4) + (180:1)$);
        \coordinate (H_C4b) at ($(C4) + (-60:1)$);
        
        % Draw sigma bonds
        \draw[thick, popdark] (H_C1) -- (C1) node[midway, below] {$\sigma$};
        \draw[thick, popdark] (C1) -- (C2) node[midway, below] {$\sigma$};
        \draw[thick, popdark] (C2) -- (C3) node[midway, below] {$\sigma$};
        \draw[thick, popdark] (C3) -- (H_C3) node[midway, below right] {$\sigma$};
        \draw[thick, popdark] (C3) -- (C4) node[midway, above left] {$\sigma$};
        \draw[thick, popdark] (C4) -- (H_C4a) node[midway, left] {$\sigma$};
        \draw[thick, popdark] (C4) -- (H_C4b) node[midway, below right] {$\sigma$};
        
        % Atom labels
        \node[popdark] at (C1) {\textbf{C$_1$}};
        \node[popdark] at (C2) {\textbf{C$_2$}};
        \node[popdark] at (C3) {\textbf{C$_3$}};
        \node[popdark] at (C4) {\textbf{C$_4$}};
        \node[popdark] at (H_C1) {H};
        \node[popdark] at (H_C3) {H};
        \node[popdark] at (H_C4a) {H};
        \node[popdark] at (H_C4b) {H};
        
        % Pi bonds indication (dashed pink)
        % Two pi bonds for C1#C2
        \draw[poppink, dashed, thick] (C1) -- (C2) node[midway, above=2mm] {$\pi, \pi$};
        % One pi bond for C3=C4
        \draw[poppink, dashed, thick] (C3) -- (C4) node[midway, above right] {$\pi$};
    \end{tikzpicture}
    \legende{Représentation de Cram du squelette $\sigma$ du but-1-en-3-yne. Les liaisons $\pi$ (pointillées roses) sont perpendiculaires au plan des $\sigma$.}
    \end{popfigure}
    \textit{Note :} Les liaisons $\pi$ (pointillées roses) sont perpendiculaires au plan des $\sigma$. En C$_1$ et C$_2$ (\textit{sp}), il y a deux nuages $\pi$ perpendiculaires entre eux. En C$_3$ et C$_4$ (\textit{sp}$^2$), un seul nuage $\pi$ perpendiculaire au plan trigonal.
\end{enumerate}
\textbf{Conclusion :} La molécule possède deux centres \textit{sp} (linéaires) et deux centres \textit{sp}$^2$ (trigonaux plans). La chaîne $\sigma$ C$_1$-C$_2$-C$_3$ est linéaire, puis le carbone C$_3$ impose un coude de 120$^\circ$ vers C$_4$.
\end{exoresolu}

\begin{methode}[Écrire les formules et nommer un alcène ou un alcyne (Nomenclature UICPA)]
\textbf{Objectif :} Passer de la formule semi-développée au nom systématique et inversement.

\begin{enumerate}
    \item \textbf{Identifier la chaîne principale :} La plus longue chaîne carbonée \textbf{contenant le double ou le triple lien}. 
    \item \textbf{Numéroter la chaîne :} Donner les numéros les plus bas possibles au multiple lien (règle du « premier point de différence » appliquée au multiple lien en priorité sur les substituants). 
    \item \textbf{Nommer les substituants :} Alkyles (méthyle, éthyle...), halogènes, etc., avec leur locant.
    \item \textbf{Construire le nom :} \textit{locants-substituants}- \textit{locant-du-multiple-lien}-\textit{préfixe-de-la-chaîne}-\textbf{-ène} (ou \textbf{-yne}). 
    \item \textbf{Stéréoisomérie (Alcènes) :} Si le double lien est substitué différemment sur chaque carbone (C$_{a}$$\neq$C$_{b}$ et C$_{c}$$\neq$C$_{d}$), préciser \textbf{(Z)} (ensemble) ou \textbf{(E)} (entgegen) selon les priorités de Cahn-Ingold-Prelog (numéro atomique).
\end{enumerate}
\textbf{Formules générales :} Alcènes cycliques : C$_n$H$_{2n}$ ; Alcènes acycliques : C$_n$H$_{2n}$ ; Alkynes acycliques : C$_n$H$_{2n-2}$.
\end{methode}

\begin{exoresolu}[Nomenclature et isomérie Z/E d'un alcène ramifié]
\textbf{Énoncé :} On donne la formule semi-développée : \ce{CH3-CH=C(CH3)-CH2-CH2-Cl}.
\begin{enumerate}
    \item Écrire la formule brute et la formule développée.
    \item Donner le nom UICPA complet (avec stéréochimie si applicable).
    \item Dessiner l'isomère de configuration (Z) et (E) de cet alcène en représentation de Cram.
\end{enumerate}

\tcblower
\textbf{Solution détaillée}

\begin{enumerate}
    \item \textbf{Formules :}
    \begin{itemize}
        \item \textbf{Développée :} \chemfig{CH_3-CH=C(-[2]CH_3)-CH_2-CH_2-Cl} (Vérification : C$_6$H$_{11}$Cl).
        \item \textbf{Brute :} \textbf{C$_6$H$_{11}$Cl} (Alcène monochloré : C$_n$H$_{2n}$Cl $\rightarrow$ C$_6$H$_{12}$Cl moins 1 H pour le Cl ? Non. Alcène C$_n$H$_{2n}$. Remplacer H par Cl $\rightarrow$ C$_n$H$_{2n-1}$Cl. Pour n=6 : C$_6$H$_{11}$Cl. Correct).
    \end{itemize}

    \item \textbf{Nomenclature UICPA :}
    \begin{enumerate}
        \item \textbf{Chaîne principale :} 5 atomes de carbone contenant le double lien (pentène). La chaîne à 6C (hexène) ne contient pas le double lien si on part du Cl ? Si on prend la chaîne à 6C : Cl-CH$_2$-CH$_2$-C(CH$_3$)=CH-CH$_3$. Le double lien est en C$_3$-C$_4$. La chaîne à 5C : le double lien est en C$_2$-C$_3$. \textbf{Règle :} La chaîne principale doit contenir le double lien. La plus longue chaîne contenant le double lien a 5 atomes de carbone (le carbone du méthyle sur C$_3$ est un substituant). $\rightarrow$ \textbf{Pent-2-ène}.
        \item \textbf{Numérotation :} Pour donner l'indice le plus bas au double lien. 
            \begin{itemize}
                \item Sens gauche $\rightarrow$ droite : Double lien en 2. Substituants : Méthyle en 3, Chloro en 5. $\rightarrow$ 3-méthyle-5-chloropent-2-ène.
                \item Sens droite $\rightarrow$ gauche : Double lien en 3. Substituants : Méthyle en 3, Chloro en 1. $\rightarrow$ 3-méthyle-1-chloropent-3-ène.
            \end{itemize}
            \textbf{Règle du premier point de différence sur le multiple lien :} 2 < 3. On numérote de gauche à droite (le CH$_3$ terminal est C$_1$).
        \item \textbf{Substituants :} En C$_3$ : un \textbf{méthyle}. En C$_5$ : un \textbf{chloro}.
        \item \textbf{Stéréochimie :} Le double lien est en C$_2$=C$_3$.
            \begin{itemize}
                \item C$_2$ : substituants = H et CH$_3$ (C$_1$). Priorité : \textbf{CH$_3$ > H} (Z atomique C=6 > H=1).
                \item C$_3$ : substituants = CH$_3$ (ramification) et CH$_2$CH$_2$Cl (chaîne). Priorité : \textbf{CH$_2$CH$_2$Cl > CH$_3$} (premier atome C pour les deux, on regarde le suivant : C, H, H pour CH$_2$Cl vs H, H, H pour CH$_3$ $\rightarrow$ C > H).
            \end{itemize}
            Sur la figure donnée (formule semi-développée linéaire), les groupes prioritaires (CH$_3$ en C$_2$ et CH$_2$CH$_2$Cl en C$_3$) sont \textbf{de part et d'autre} du double lien (trans) $\rightarrow$ Configuration \textbf{(E)}.
        \item \textbf{Nom complet :} \textbf{(E)-3-méthyle-5-chloropent-2-ène} (ou 5-chloro-3-méthylepent-2-ène par ordre alphabétique : chloro avant méthyle).
    \end{enumerate}
    \textbf{Nom retenu :} \cle{(E)-5-chloro-3-méthylepent-2-ène}.

    \item \textbf{Représentations de Cram (Isomères Z et E) :}
    On fixe la chaîne principale dans le plan. Le double lien C$_2$=C$_3$ est rigide.
    \begin{popfigure}
    \begin{center}
    \begin{tabular}{cc}
    \textbf{(E)-5-chloro-3-méthylepent-2-ène} & \textbf{(Z)-5-chloro-3-méthylepent-2-ène} \\
    \begin{tikzpicture}[scale=0.8, baseline=(C2.base)]
        \coordinate (C2) at (0,0);
        \coordinate (C3) at (2,0);
        \draw[thick, popdark] (C2) -- (C3);
        \draw[poppink, dashed, thick] (C2) -- (C3);
        % C2: Pri=CH3 (UP), H (DOWN)
        \draw[thick, popdark] (C2) -- ++(90:1) node[above] {CH$_3$ (Pri)};
        \draw[thick, popdark, densely dashed] (C2) -- ++(-90:1) node[below] {H};
        % C3: Pri=CH2CH2Cl (DOWN), CH3 (UP)
        \draw[thick, popdark] (C3) -- ++(-90:1) node[below] {CH$_2$CH$_2$Cl (Pri)};
        \draw[thick, popdark, densely dashed] (C3) -- ++(90:1) node[above] {CH$_3$};
        \node[popdark] at (C2) {\textbf{C$_2$}};
        \node[popdark] at (C3) {\textbf{C$_3$}};
    \end{tikzpicture}
    &
    \begin{tikzpicture}[scale=0.8, baseline=(C2.base)]
        \coordinate (C2) at (0,0);
        \coordinate (C3) at (2,0);
        \draw[thick, popdark] (C2) -- (C3);
        \draw[poppink, dashed, thick] (C2) -- (C3);
        % C2: Pri=CH3 (UP), H (DOWN)
        \draw[thick, popdark] (C2) -- ++(90:1) node[above] {CH$_3$ (Pri)};
        \draw[thick, popdark, densely dashed] (C2) -- ++(-90:1) node[below] {H};
        % C3: Pri=CH2CH2Cl (UP), CH3 (DOWN) -> Z (Zusammen)
        \draw[thick, popdark] (C3) -- ++(90:1) node[above] {CH$_2$CH$_2$Cl (Pri)};
        \draw[thick, popdark, densely dashed] (C3) -- ++(-90:1) node[below] {CH$_3$};
        \node[popdark] at (C2) {\textbf{C$_2$}};
        \node[popdark] at (C3) {\textbf{C$_3$}};
    \end{tikzpicture}
    \end{tabular}
    \end{center}
    \legende{Représentations de Cram des isomères (E) et (Z). Trait plein : liaison dans le plan (ou wedge sortant), trait pointillé : liaison rentrante (dash). Les groupes prioritaires sont en gras.}
    \end{popfigure}
    \textbf{Conclusion :} La numérotation privilégie le double lien (indice 2). La priorité CIP place le groupe CH$_2$CH$_2$Cl devant CH$_3$ sur C$_3$. L'isomère donné correspond à la configuration (E).
\end{enumerate}
\end{exoresolu}

\begin{methode}[Écrire les équations des réactions d'addition sur les alcènes (Règle de Markovnikov)]
\textbf{Objectif :} Prédire le(s) produit(s) majeur(s) de l'addition de HX, H$_2$O (hydratation), X$_2$, H$_2$ sur un alcène asymétrique.

\begin{enumerate}
    \item \textbf{Identifier l'alcène :} Repérer le double lien C=C et les substituants sur chaque carbone (C$_a$ et C$_b$).
    \item \textbf{Addition de HX / H$_2$O (milieu acide) :} Réaction ionique via \textbf{carbocation intermédiaire}.
        \begin{itemize}
            \item L'H$^+$ s'ajoute sur le carbone \textbf{le plus hydrogéné} (le moins substitué) pour former le carbocation \textbf{le plus stable} (tertiaire $>$ secondaire $>$ primaire) sur l'autre carbone.
            \item L'anion (X$^-$ ou H$_2$O) attaque le carbocation.
        \end{itemize}
        \cle{Règle de Markovnikov :} « Le riche devient plus riche » (L'H va sur le C qui a déjà le plus de H ; le X/OH va sur le C qui a le moins de H).
    \item \textbf{Addition de X$_2$ (Br$_2$, Cl$_2$) :} Réaction via \textbf{ion halonium cyclique} (pont halogène). 
        \begin{itemize}
            \item Ajout \textit{anti} (stéréospécifique). 
            \item Pas de réarrangement de carbocation. Produit : vicinal dihalogéné.
        \end{itemize}
    \item \textbf{Hydrogénation (H$_2$, Pt/Ni/Pd) :} Ajout \textit{syn} de 2 H. Alcène $\rightarrow$ Alcane.
    \item \textbf{Écrire l'équation :} Réactifs $\rightarrow$ Produit(s) majeur(s) + (éventuel mineur).
\end{enumerate}
\textbf{Attention :} Pour l'hydratation, écrire H$_2$O / H$^+$ (catalyseur) au-dessus de la flèche. Le produit est un alcool.
\end{methode}

\begin{exoresolu}[Additions sur le 3-méthylepent-1-ène : régiosélectivité et stéréochimie]
\textbf{Énoncé :} On fait réagir le 3-méthylepent-1-ène (\ce{CH2=CH-CH(CH3)-CH2-CH3}) avec :
\begin{enumerate}
    \item Du chlorure d'hydrogène gazeux (HCl).
    \item De l'eau en milieu acide (H$_2$O / H$^+$).
    \item Du dibrome (Br$_2$) dans le tétrachlorure de carbone (CCl$_4$).
    \item Du dihydrogène (H$_2$) en présence de platine (Pt).
\end{enumerate}
Pour chaque réaction : écrire l'équation de la réaction, nommer le(s) produit(s) majeur(s) et justifier la régiosélectivité le cas échéant.

\tcblower
\textbf{Solution détaillée}

\textbf{Structure du réactif :} \ce{CH2=CH-CH(CH3)-CH2-CH3}. 
Carbone C$_1$ (CH$_2$=) : 2 H (le plus hydrogéné, primaire).
Carbone C$_2$ (=CH$-$) : 1 H, lié à C$_1$, C$_3$ (secondaire).
Carbone C$_3$ : CH(CH$_3$)-CH$_2$CH$_3$ (ramifié). \textbf{Attention :} C$_3$ est un carbone chiral (asymétrique) dans le réactif (4 substituants différents : H, CH$_3$, CH$_2$CH$_3$, CH=CH$_2$). Le réactif est généralement un mélange racémique.

\begin{enumerate}
    \item \textbf{Addition de HCl (Addition ionique, Markovnikov) :}
    \begin{itemize}
        \item \textbf{Mécanisme :} Attaque du double lien sur H$^+$. 
        \item \textbf{Choix du carbocation :} 
            \begin{itemize}
                \item Si H$^+$ sur C$_1$ $\rightarrow$ Carbocation sur C$_2$ (Secondaire, stabilisé par effet inductif donneur du groupe alkyle C$_3$). \textbf{Favorisé}.
                \item Si H$^+$ sur C$_2$ $\rightarrow$ Carbocation sur C$_1$ (Primaire, très instable). \textbf{Défavorisé}.
            \end{itemize}
        \item \textbf{Produit majeur :} Cl$^-$ attaque le carbocation C$_2^+$.
        \item \textbf{Équation :} 
        \[ \ce{CH2=CH-CH(CH3)CH2CH3 + HCl -> CH3-CH(Cl)-CH(CH3)CH2CH3} \]
        \item \textbf{Nom :} \textbf{2-chloro-3-méthylepentane}. 
        \item \textbf{Stéréochimie :} Création d'un centre chiral en C$_2$. Le réactif est déjà chiral en C$_3$ (racémique). On obtient un mélange de \textbf{diastéréoisomères} (2 paires d'énantiomères, 4 stéréoisomères au total).
    \end{itemize}

    \item \textbf{Hydratation (H$_2$O / H$^+$) :}
    \begin{itemize}
        \item Même régiosélectivité (Markovnikov). L'H$^+$ va sur C$_1$, carbocation sur C$_2$. H$_2$O attaque C$_2^+$, puis déprotonation.
        \item \textbf{Équation :} 
        \[ \ce{CH2=CH-CH(CH3)CH2CH3 + H2O ->[H^+] CH3-CH(OH)-CH(CH3)CH2CH3} \]
        \item \textbf{Nom :} \textbf{3-méthylepentan-2-ol} (Alcool secondaire). 
        \item \textbf{Stéréochimie :} Création d'un centre chiral en C$_2$. Deux centres chiraux (C$_2$, C$_3$) $\rightarrow$ Mélange de \textbf{diastéréoisomères}.
    \end{itemize}

    \item \textbf{Addition de Br$_2$ (CCl$_4$) :}
    \begin{itemize}
        \item Pas de carbocation $\rightarrow$ Pas de règle de Markovnikov. Ajout vicinal \textit{anti} via ion bromonium.
        \item Les deux carbones C$_1$ et C$_2$ reçoivent un atome de brome. L'attaque de Br$^-$ se fait sur le carbone le plus substitué (C$_2$) par ouverture de l'ion bromonium (régiosélectivité faible mais orientée).
        \item \textbf{Équation :} 
        \[ \ce{CH2=CH-CH(CH3)CH2CH3 + Br2 ->[CCl4] BrCH2-CH(Br)-CH(CH3)CH2CH3} \]
        \item \textbf{Nom :} \textbf{1,2-dibromo-3-méthylepentane}.
        \item \textbf{Stéréochimie :} C$_1$ devient CH$_2$Br (achiral). C$_2$ devient chiral. C$_3$ est chiral (inchangé). L'addition \textit{anti} fixe la configuration relative entre les deux atomes de Br sur C$_1$ et C$_2$, mais C$_1$ n'est pas chiral. Le centre C$_2$ est créé sans stéréosélectivité (attaque des deux faces de l'ion bromonium équiprobalables). On obtient un mélange de \textbf{diastéréoisomères} (car C$_3$ est chiral et racémique).
    \end{itemize}

    \item \textbf{Hydrogénation (H$_2$, Pt) :}
    \begin{itemize}
        \item Ajout \textit{syn} de 2 H sur le double lien. Saturation totale.
        \item \textbf{Équation :} 
        \[ \ce{CH2=CH-CH(CH3)CH2CH3 + H2 ->[Pt] CH3-CH2-CH(CH3)CH2CH3} \]
        \item \textbf{Nom :} \textbf{3-méthylepentane} (Alcane ramifié). 
        \item \textbf{Stéréochimie :} C$_1$ et C$_2$ deviennent sp$^3$ non chiraux (CH$_3$ et CH$_2$). C$_3$ porte désormais deux groupes éthyle identiques (CH$_2$CH$_3$ de part et d'autre) : \textbf{C$_3$ n'est plus chiral}. Le produit final est \textbf{achiral}.
    \end{itemize}
\end{enumerate}

\textbf{Conclusion :} Les additions ioniques (HCl, H$_2$O/H$^+$) suivent la règle de Markovnikov via le carbocation le plus stable (secondaire ici). L'addition de Br$_2$ est stéréospécifique \textit{anti} et crée un centre chiral en C$_2$ (mélange de diastéréoisomères car C$_3$ est chiral). L'hydrogénation est \textit{syn}, saturante, et fait perdre la chiralité au C$_3$ par symétrisation.
\end{exoresolu}

\begin{methode}[Réaliser le dosage d'un alcène par bromométrie (Détermination de l'insaturation)]
\textbf{Objectif :} Doser une solution d'alcène par ajout excédentaire de bromure/bromate (ou solution de Br$_2$ titrée) et dosage de l'excès de brome par thiosulfate de sodium.

\begin{enumerate}
    \item \textbf{Réaction de consommation :} Alcène + Br$_2$ $\rightarrow$ Dibromé (rapide, totale). 
        \[ \ce{C=C + Br2 -> CBr-CBr} \quad (1 \text{ mol alcène consomme } 1 \text{ mol Br}_2) \]
    \item \textbf{Réaction de dosage de l'excès (Iodométrie) :} 
        \[ \ce{Br2 + 2 I- -> 2 Br- + I2} \quad \text{(dans l'excès de KI)} \]
        \[ \ce{I2 + 2 S2O3^2- -> 2 I- + S4O6^2-} \quad \text{(Dosage au thiosulfate)} \]
        \textbf{Bilan :} 1 mol Br$_2$ $\equiv$ 2 mol S$_2$O$_3^{2-}$.
    \item \textbf{Protocole :}
        \begin{enumerate}
            \item Introduire un volume connu $V_A$ de solution d'alcène (concentration $C_A$ inconnue) dans un erlenmeyer.
            \item Ajouter un volume connu $V_{Br}$ de solution de brome (concentration $C_{Br}$ connue) \textbf{en excès certain} ($n_{Br_2}^{init} > n_{alcène}$).
            \item Laisser réagir (quelques minutes, à l'abri de la lumière pour Br$_2$).
            \item Ajouter un excès de KI solide (ou solution concentrée) $\rightarrow$ libération d'I$_2$ proportionnelle au Br$_2$ restant.
            \item Doser l'I$_2$ libéré par une solution de thiosulfate de sodium $C_{S_2O_3}$ (burette) jusqu'à couleur jaune paille.
            \item Ajouter quelques gouttes d'amidon (bleu vif) et continuer le dosage jusqu'à décoloration complète (point d'équivalence).
        \end{enumerate}
    \item \textbf{Calculs (Tableau d'avancement recommandé) :}
        \begin{align*}
            n_{Br_2}^{init} &= C_{Br} \times V_{Br} \\
            n_{S_2O_3}^{eq} &= C_{S_2O_3} \times V_{S_2O_3}^{eq} \\
            n_{Br_2}^{restant} &= \frac{1}{2} n_{S_2O_3}^{eq} \\
            n_{Br_2}^{consommé} &= n_{Br_2}^{init} - n_{Br_2}^{restant} \\
            n_{alcène} &= n_{Br_2}^{consommé} \quad (\text{stoïchiométrie 1/1}) \\
            C_A &= \frac{n_{alcène}}{V_A}
        \end{align*}
\end{enumerate}
\end{methode}

\begin{exoresolu}[Dosage de l'insaturation d'un mélange d'hex-1-ène]
\textbf{Énoncé :} On souhaite doser l'hex-1-ène (C$_6$H$_{12}$, $M = \SI{84,16}{g.mol^{-1}}$) contenu dans un échantillon commercial. 
On prélève $V_A = \SI{10,0}{mL}$ de cette solution d'hex-1-ène (densité $\rho = \SI{0,673}{g.mL^{-1}}$, on suppose la solution pure pour le calcul de la concentration molaire théorique, mais on veut la concentration réelle par dosage).
On ajoute $V_{Br} = \SI{25,0}{mL}$ d'une solution de brome $C_{Br} = \SI{0,100}{mol.L^{-1}}$.
Après réaction complète, on ajoute un excès de KI. Le mélange libère de l'iode que l'on dose par une solution de thiosulfate de sodium $C_{S_2O_3} = \SI{0,100}{mol.L^{-1}}$. 
Volume de thiosulfate versé à l'équivalence : $V_{S_2O_3} = \SI{18,0}{mL}$.
\begin{enumerate}
    \item Écrire les équations-bilan des réactions mises en jeu.
    \item Calculer la quantité de matière d'hex-1-ène présente dans l'échantillon dosé.
    \item En déduire la concentration molaire réelle de la solution d'hex-1-ène.
    \item Calculer le titre massique (en \%) de l'hex-1-ène dans le liquide commercial (on rappelle $\rho = \SI{0,673}{g.mL^{-1}}$).
\end{enumerate}

\tcblower
\textbf{Solution détaillée}

\begin{enumerate}
    \item \textbf{Équations-bilan :}
    \begin{enumerate}
        \item \textbf{Addition sur l'alcène (consommation de Br$_2$) :}
        \[ \ce{C6H12 + Br2 -> C6H12Br2} \quad (\text{1,2-dibromohexane}) \]
        \item \textbf{Réaction du brome en excès avec l'iodure (libération d'iode) :}
        \[ \ce{Br2 + 2 I- -> 2 Br- + I2} \]
        \item \textbf{Dosage de l'iode par le thiosulfate :}
        \[ \ce{I2 + 2 S2O3^2- -> 2 I- + S4O6^2-} \]
    \end{enumerate}
    \textbf{Bilan global Br$_2$ / S$_2$O$_3^{2-}$ :} 1 mol Br$_2$ $\rightarrow$ 1 mol I$_2$ $\rightarrow$ 2 mol S$_2$O$_3^{2-}$.

    \item \textbf{Calcul des quantités de matière :}
    \begin{itemize}
        \item $n_{Br_2}^{init} = C_{Br} \times V_{Br} = \SI{0,100}{mol.L^{-1}} \times \SI{25,0E-3}{L} = \SI{2,50E-3}{mol}$.
        \item $n_{S_2O_3}^{eq} = C_{S_2O_3} \times V_{S_2O_3} = \SI{0,100}{mol.L^{-1}} \times \SI{18,0E-3}{L} = \SI{1,80E-3}{mol}$.
        \item $n_{Br_2}^{restant} = \frac{1}{2} n_{S_2O_3}^{eq} = \frac{1,80 \times 10^{-3}}{2} = \SI{0,90E-3}{mol}$.
        \item $n_{Br_2}^{consommé} = n_{Br_2}^{init} - n_{Br_2}^{restant} = (2,50 - 0,90) \times 10^{-3} = \SI{1,60E-3}{mol}$.
    \end{itemize}

    \item \textbf{Quantité de matière d'hex-1-ène :}
    Stoïchiométrie : 1 mol hex-1-ène $\equiv$ 1 mol Br$_2$ consommé.
    \[ n_{C_6H_{12}} = n_{Br_2}^{consommé} = \boxed{\SI{1,60E-3}{mol}} \]

    \item \textbf{Concentration molaire de la solution :}
    \[ C_{C_6H_{12}} = \frac{n_{C_6H_{12}}}{V_A} = \frac{\SI{1,60E-3}{mol}}{\SI{10,0E-3}{L}} = \boxed{\SI{0,160}{mol.L^{-1}}} \]

    \item \textbf{Titre massique (pureté) :}
    \begin{itemize}
        \item Masse d'hex-1-ène dans $V_A = \SI{10,0}{mL}$ : 
        $m_{C_6H_{12}} = n \times M = \SI{1,60E-3}{mol} \times \SI{84,16}{g.mol^{-1}} = \SI{0,13466}{g}$.
        \item Masse de la solution dosée ($V_A = \SI{10,0}{mL}$) :
        $m_{sol} = \rho \times V_A = \SI{0,673}{g.mL^{-1}} \times \SI{10,0}{mL} = \SI{6,73}{g}$.
        \item Titre massique $w$ :
        \[ w = \frac{m_{C_6H_{12}}}{m_{sol}} \times 100 = \frac{0,13466}{6,73} \times 100 = \boxed{2,00\%} \]
        \textit{(Note : Ce titre faible suggère une solution diluée dans un solvant, ou une erreur dans l'hypothèse "solution pure". Si c'était de l'hex-1-ène pur, $C \approx \frac{673}{84,16} \approx 8,0$ mol.L$^{-1}$. Le dosage montre une solution à 0,16 M, soit environ 2\% massique).}
    \end{itemize}
\end{enumerate}
\textbf{Conclusion :} La solution commerciale contient 2,00\% en masse d'hex-1-ène (concentration molaire \SI{0,160}{mol.L^{-1}}). Le dosage par bromométrie indirecte (iodométrie) est une méthode classique de détermination de l'indice d'iode / insaturation.
\end{exoresolu}

\begin{methode}[Préparer l'acétylène (éthyne) au laboratoire]
\textbf{Objectif :} Mettre en œuvre la réaction de l'éthylénure de calcium (carbure de calcium) avec l'eau pour produire de l'éthyne gazeux, le purifier et le recueillir.

\begin{enumerate}
    \item \textbf{Réaction chimique :} 
    \[ \ce{CaC2(s) + 2 H2O(l) -> Ca(OH)2(aq/s) + C2H2(g)} \]
    \item \textbf{Dispositif (Générateur de Kipp ou montage simple à entonnoir/burette) :}
        \begin{itemize}
            \item \textbf{Réacteur :} Creuset ou tube à essai (petites quantités) ou générateur de Kipp (production continue/contrôlée). Le solide (CaC$_2$) est au fond, l'eau est ajoutée goutte à goutte.
            \item \textbf{Purification (indispensable) :} Le gaz brut contient des impuretés toxiques et odorantes : \ce{H2S} (odeur œuf pourri), \ce{PH3} (phosphine, inflammable spontané), \ce{NH3}, \ce{AsH3}, \ce{SiH4}.
            \item \textbf{Lavage successif (bouteilles à laver / tubes à laver) :}
                \begin{enumerate}
                    \item \textbf{Eau :} Dissout NH$_3$, une partie des hydrures, refroidit le gaz.
                    \item \textbf{Solution de CuSO$_4$ + HCl (ou AgNO$_3$) :} Piège \ce{H2S} (précipité CuS noir) et \ce{PH3}.
                    \item \textbf{Solution de NaOH concentrée :} Piège \ce{CO2} (si air aspiré), \ce{H2S} résiduel, \ce{Cl2} (si eau de Javel utilisée). \textit{Attention : Ne pas utiliser NaOH si on veut garder l'acétylène pur pour réaction ultérieure car il peut réagir ? Non, C2H2 ne réagit pas avec NaOH froid.}
                \end{enumerate}
            \item \textbf{Séchage :} Tube à U rempli de \textbf{chlorure de calcium anhydre CaCl$_2$} (ou P$_2$O$_5$, H$_2$SO$_4$ conc.). \textit{Interdit :} Pentoxyde de phosphore (risque phosphine), Acide sulfurique concentré (risque polymérisation/charbonnage si chaud). CaCl$_2$ est le standard.
            \item \textbf{Recueil :} Par déplacement de l'eau (gaz peu soluble, $\approx$ 1,2 g/L à 20°C) ou par déplacement d'air (siphon inversé) si besoin de gaz sec pur.
        \end{itemize}
    \item \textbf{Sécurité :} \ce{C2H2} \textbf{TRÈS INFLAMMABLE} (limites 2,5 - 82\% vol). \textbf{EXPLOSIF} sous pression ($> \SI{1,5}{bar}$) ou pur comprimé (décomposition exothermique $\ce{C2H2 -> 2 C + H2}$). \textbf{Jamais} comprimer l'acétylène pur. Utiliser des bouteilles spéciales (acétone + masse poreuse) pour le stockage industriel. Au labo : produire et consommer immédiatement. Port des lunettes, manipulation sous hotte.
\end{enumerate}
\textbf{Schéma type (Générateur de Kipp) :} Trois corps emboîtés. Corps inférieur : CaC$_2$. Corps central : Eau. Corps supérieur : Réservoir eau + tuyau de sortie gaz. Robinet sur sortie gaz. Ouverture robinet $\rightarrow$ pression gaz fait descendre eau $\rightarrow$ contact eau/CaC$_2$ $\rightarrow$ réaction. Fermeture robinet $\rightarrow$ pression gaz remonte eau $\rightarrow$ arrêt réaction.
\end{methode}

\begin{exoresolu}[Préparation et dosage de l'acétylène au laboratoire]
\textbf{Énoncé :} On souhaite préparer de l'éthyne (acétylène) au laboratoire à partir de \SI{5,00}{g} d'éthylénure de calcium (CaC$_2$, $M = \SI{64,10}{g.mol^{-1}}$) technique (pureté 80\% en masse). Le gaz est lavé, séché sur CaCl$_2$ et recueilli par déplacement de l'eau dans une burette graduée retournée (\SI{25}{\celsius}, $P_{atm} = \SI{1,013 \times 10^5}{Pa}$). Tension de vapeur de l'eau à 25°C : $P_{H_2O}^\circ = \SI{3,17 \times 10^3}{Pa}$.
\begin{enumerate}
    \item Écrire l'équation de la réaction de préparation. Préciser l'état physique des espèces.
    \item Calculer le volume théorique d'éthyne sec récupérable (en L) dans ces conditions.
    \item Décrire le dispositif de lavage/séchage en justifiant le rôle de chaque bouteille.
    \item Pourquoi ne recueille-t-on pas l'acétylène par déplacement de l'air dans un ballon ? Pourquoi le déplacement d'eau est-il préféré pour le dosage volumétrique ?
\end{enumerate}

\tcblower
\textbf{Solution détaillée}

\begin{enumerate}
    \item \textbf{Équation de la réaction :}
    \[ \ce{CaC2(s) + 2 H2O(l) -> Ca(OH)2(aq) + C2H2(g)} \]
    (L'hydroxyde de calcium est peu soluble, il se dépose en grande partie sous forme de lait de chaux / solide).

    \item \textbf{Calcul du volume théorique :}
    \begin{itemize}
        \item Masse de CaC$_2$ pur : $m_{pur} = \SI{5,00}{g} \times 0,80 = \SI{4,00}{g}$.
        \item Quantité de matière CaC$_2$ : $n_{CaC_2} = \frac{\SI{4,00}{g}}{\SI{64,10}{g.mol^{-1}}} = \SI{6,240E-2}{mol}$.
        \item Stoïchiométrie : 1 mol CaC$_2$ $\rightarrow$ 1 mol C$_2$H$_2$.
        $n_{C_2H_2} = \SI{6,240E-2}{mol}$.
        \item Conditions de recueil (déplacement d'eau) : Gaz \textbf{humide} (saturé en vapeur d'eau).
        Pression partielle de l'éthyne sec : $P_{C_2H_2} = P_{atm} - P_{H_2O}^\circ = \SI{1,013E5}{Pa} - \SI{3,17E3}{Pa} = \SI{0,9813E5}{Pa} = \SI{9,813E4}{Pa}$.
        \item Loi des gaz parfaits : $V = \frac{nRT}{P}$.
        $R = \SI{8,314}{J.mol^{-1}.K^{-1}}$, $T = \SI{298,15}{K}$.
        \[ V = \frac{\SI{6,240E-2}{mol} \times \SI{8,314}{J.mol^{-1}.K^{-1}} \times \SI{298,15}{K}}{\SI{9,813E4}{Pa}} \]
        \[ V \approx \frac{154,5}{98130} \approx \SI{1,575E-3}{m^3} = \boxed{\SI{1,58}{L}} \]
        \textit{(Volume d'éthyne sec. Le volume lu sur la burette sera celui du mélange gazeux humide, légèrement supérieur : $V_{humide} = n_{tot}RT/P_{atm}$. Mais on demande le volume d'éthyne sec récupérable, c'est bien $V$ calculé ci-dessus).}
    \end{itemize}

    \item \textbf{Dispositif de lavage/séchage (3 bouteilles à laver en série) :}
    \begin{center}
    \begin{tabularx}{\linewidth}{|l|X|l|}
    \hline
    \rowcolor{popblueL}
    \textbf{Bouteille} & \textbf{Contenu} & \textbf{Rôle / Impuretés éliminées} \\ \hline
    1\textsuperscript{re} & Eau distillée & Refroidissement, dissolution \ce{NH3}, \ce{H2S} partiel, \ce{PH3} partiel. Enlève poussières Ca(OH)2. \\ \hline
    2\textsuperscript{e} & Solution \ce{CuSO4} acidifiée (HCl) \textbf{OU} \ce{AgNO3} & \textbf{Piégeage sélectif \ce{H2S} (CuS noir) et \ce{PH3} (Cu3P2 / Ag3P).} Indispensable pour toxicité/odeur. \\ \hline
    3\textsuperscript{e} & \ce{NaOH} aqueux concentré (ou \ce{KOH}) & Neutralisation \ce{CO2} (air), \ce{H2S} résiduel, \ce{Cl2} (si eau de Javel). \\ \hline
    \textbf{Tube à U} & \textbf{CaCl2 anhydre} (granulats) & \textbf{Séchage final} (absorption H2O gazeuse). \textit{Pas H2SO4 conc (polymérisation), pas P2O5 (poussière/phosphine)}. \\ \hline
    \end{tabularx}
    \end{center}

    \item \textbf{Mode de recueil :}
    \begin{itemize}
        \item \textbf{Pas de ballon / déplacement d'air :} L'acétylène forme des mélanges \textbf{explosifs} avec l'air (limites 2,5\% - 82\%). Un ballon peut se percer, créer une poche explosive. De plus, le dosage volumétrique est imprécis (pression variable, fuite).
        \item \textbf{Déplacement d'eau (burette graduée) :} 
            \begin{itemize}
                \item Permet un \textbf{dosage volumétrique précis} (lecture du volume liquide).
                \item La pression du gaz est connue ($P_{atm} - P_{H_2O}^\circ$).
                \item L'eau ne dissout que très peu d'acétylène ($\approx$ 1,2 g/L, perte $< 2\%$).
                \item Sécurité : Le gaz est confiné dans l'appareillage verre/eau, pas de mélange air/gaz libre dans la pièce.
            \end{itemize}
    \end{itemize}
\end{enumerate}
\textbf{Conclusion :} On récupère théoriquement \SI{1,58}{L} d'éthyne sec. La purification par CuSO$_4$/HCl et le séchage au CaCl$_2$ sont obligatoires pour obtenir un gaz pur utilisable en synthèse (ex: vinylation) ou pour l'analyse.
\end{exoresolu}

\begin{methode}[Écrire l'équation de polymérisation et identifier le monomère / le polymère]
\textbf{Objectif :} Passer de la formule du monomère (alcène/alcyne) à la formule de la chaîne polymère (unité répétitive) et inversement.

\begin{enumerate}
    \item \textbf{Polymérisation par addition (alcènes) :} Ouverture du double lien. 
\begin{itemize}
            \item Monomère : $CH_2=CH-R$ (alcène vinylène) ou $CH_2=CR-R'$ etc.
            \item Polymère : $-[CH_2-CHR]_n$ (chaîne carbonée saturée, pendants R).
            \item Écriture : $n\ CH_2=CH-R \xrightarrow{\text{initiateur}} -[CH_2-CHR]_n$.
        \end{itemize}
    \item \textbf{Polymérisation des alkynes :} Possible mais donne des polymères conjugués (polyacétylène) conducteurs. Moins courant au programme standard (focus sur alcènes).
    \item \textbf{Identification du monomère depuis le polymère :}
        \begin{enumerate}
            \item Repérer l'unité répétitive (motif entre crochets \ce{[- ... -]_n}).
            \item Couper la chaîne C$-$C pour reformer un double lien C=C.
            \item Chaque carbone de l'ancien double lien récupère une liaison $\pi$ : on « déplie » la chaîne.
        \end{enumerate}
    \item \textbf{Nomenclature polymère :} \textit{poly(nom du monomère)}. Ex: \ce{CH2=CHCl} (chloroéthène) $\rightarrow$ \textit{polychloroéthène} (PVC).
\end{enumerate}
\textbf{Point de vigilance :} Stéréochimie (tactique) : isotactique, syndiotactique, atactique (non demandé au Probatoire mais bon de savoir que le carbone chiral créé lors de la polymérisation du propène donne des mélanges).
\end{methode}

\begin{exoresolu}[Du monomère au polymère : cas du chloroprène et du styrène]
\textbf{Énoncé :}
\begin{enumerate}
    \item Le chloroprène (2-chlorobuta-1,3-diène) est le monomère du néoprène (caoutchouc synthétique). Sa formule semi-développée : \ce{CH2=C(Cl)-CH=CH2}.
        Écrire l'équation de sa polymérisation par addition 1,4 (la plus courante) en donnant la formule de l'unité répétitive du polychloroprène.
    \item On donne l'unité répétitive d'un polymère : \ce{-[CH2-CH(C6H5)]-_n} (où \ce{C6H5} = phényle).
        Identifier le monomère, donner son nom usuel et son nom UICPA.
    \item Pour le polychloroprène, préciser la nature des doubles liaisons résiduelles dans la chaîne (configuration \textit{cis} ou \textit{trans} majoritaire ?).
\end{enumerate}

\tcblower
\textbf{Solution détaillée}

\begin{enumerate}
    \item \textbf{Polymérisation du chloroprène (Addition 1,4) :}
    Le chloroprène est un diène conjugué. La polymérisation 1,4 ouvre les deux doubles liaisons pour former une chaîne avec un double lien résiduel en position 2 (entre C$_2$ et C$_3$ de l'ancien monomère).
    \begin{itemize}
        \item Monomère : \ce{CH2=C(Cl)-CH=CH2} (Carbones numérotés 1, 2, 3, 4).
        \item Attaque radicalaire/ionique sur C$_1$ ou C$_4$. Propagation par ajout sur C$_4$ ou C$_1$.
        \item Unité répétitive issue de l'addition 1,4 : Les liaisons se font entre C$_1$ d'un monomère et C$_4$ du suivant.
        \item Structure de l'unité : \ce{-CH2-C(Cl)=CH-CH2-}.
        \item \textbf{Équation :}
        \[ n \ce{CH2=C(Cl)-CH=CH2} \xrightarrow{\text{initiateur}} \ce{-[CH2-C(Cl)=CH-CH2]-_n} \]
        \item \textbf{Nom du polymère :} \textbf{Polychloroprène} (Néoprène).
    \end{itemize}

    \item \textbf{Identification du monomère depuis \ce{-[CH2-CH(C6H5)]-_n} :}
    \begin{itemize}
        \item Unité répétitive : \ce{-CH2-CH(Ph)-} (Ph = phényle \ce{C6H5}).
        \item La chaîne polymère est saturée (liaisons C$-$C simples). Le double lien du monomère a été ouvert.
        \item Pour retrouver le monomère : On coupe une liaison C$-$C de la chaîne principale et on remet un double lien entre les deux carbones concernés.
        \item Carbone 1 (CH$_2$) et Carbone 2 (CH-Ph) $\rightarrow$ Double lien entre C$_1$ et C$_2$.
        \item Monomère : \ce{CH2=CH-Ph} (\ce{CH2=CH-C6H5}).
        \item \textbf{Nom usuel :} \textbf{Styrène}.
        \item \textbf{Nom UICPA :} \textbf{Phényéthène} (ou Éthénylbenzène).
        \item \textbf{Polymère :} \textbf{Polystyrène} (PS).
    \end{itemize}

    \item \textbf{Configuration du double lien résiduel (Polychloroprène) :}
    Lors de la polymérisation 1,4 d'un diène, le double lien résiduel (C$_2$=C$_3$) peut être \textit{cis} ou \textit{trans}.
    \begin{itemize}
        \item Polymérisation radicalaire (standard) : Mélange \textit{cis}/\textit{trans} (souvent majoritairement \textit{trans} pour la stabilité thermodynamique, mais cinétique favorise \textit{cis} parfois).
        \item Polymérisation de Ziegler-Natta / coordination : Peut être stéréosélective.
        \item \textbf{Néoprène commercial (polychloroprène) :} Majoritairement \textbf{\textit{trans}-1,4} (environ 90\% \textit{trans}). Cela lui donne une cristallinité élevée (renforcement auto-renforçant à l'étirement), propriété clé pour un caoutchouc.
    \end{itemize}
    \textbf{Conclusion :} Le polychloroprène a une unité répétitive \ce{-[CH2-C(Cl)=CH-CH2]-_n} avec un double lien résiduel majoritairement \textit{trans}. Le styrène (\ce{CH2=CHPh}) donne le polystyrène \ce{-[CH2-CHPh]-_n}.
\end{enumerate}
\end{exoresolu}

\begin{methode}[Résoudre un problème de synthèse organique multi-étapes (Alcène $\rightarrow$ Alcool / Halogéné)]
\textbf{Objectif :} Proposer un chemin de synthèse pour obtenir une molécule cible à partir d'un alcène donné, en utilisant les réactions d'addition étudiées (Markovnikov, addition Br$_2$, hydrogénation).

\begin{enumerate}
    \item \textbf{Analyse rétrosynthétique (Raisonnement à l'envers) :}
        \begin{itemize}
            \item Cible : Alcool secondaire ? $\leftarrow$ Hydratation d'alcène (Markovnikov). 
            \item Cible : Halogénure d'alkyle secondaire ? $\leftarrow$ Addition HX sur alcène (Markovnikov).
            \item Cible : Dihalogéné vicinal ? $\leftarrow$ Addition X$_2$ sur alcène.
            \item Cible : Alcane ? $\leftarrow$ Hydrogénation alcène.
        \end{itemize}
    \item \textbf{Vérification de la régiosélectivité :} Pour chaque étape d'addition sur alcène asymétrique, appliquer la règle de Markovnikov (H sur le C le plus hydrogéné, X/OH sur le plus substitué $\rightarrow$ carbocation stable).
    \item \textbf{Gestion de la stéréochimie :} 
        \begin{itemize}
            \item Addition HX / H$_2$O : Création centre chiral $\rightarrow$ Mélange racémique (pas de contrôle stéréochimique au programme).
            \item Addition Br$_2$ : Ajout \textit{anti} $\rightarrow$ Si alcène cyclique ou (Z)/(E) défini, prédire diastéréoisomères (\textit{meso} ou racémique).
        \end{itemize}
    \item \textbf{Écriture de la synthèse directe :} Réactifs, conditions (solvant, catalyseur), équation pour chaque étape. Rendement global = produit des rendements.
\end{enumerate}
\end{methode}

\begin{exoresolu}[Synthèse du 2-bromo-3-méthylebutane à partir du 3-méthylebut-1-ène]
\textbf{Énoncé :} On part du 3-méthylebut-1-ène (isopentène, \ce{CH2=CH-CH(CH3)2}). On souhaite obtenir le 2-bromo-3-méthylebutane (\ce{CH3-CH(Br)-CH(CH3)2}).
\begin{enumerate}
    \item Proposer une synthèse en une étape. Écrire l'équation et les conditions.
    \item L'expérience est réalisée avec \SI{10,0}{g} d'alcène (pureté 95\%) et un excès d'HCl gazeux (erreur volontaire dans l'énoncé : on veut du \textbf{brome}, pas du chlore). Corriger le réactif nécessaire et calculer la masse théorique de produit bromé obtenue ($M_{alcène}=\SI{70,13}{g.mol^{-1}}$, $M_{produit}=\SI{149,04}{g.mol^{-1}}$).
    \item Le produit obtenu est-il chiral ? Justifier. Quel mélange obtient-on ?
\end{enumerate}

\tcblower
\textbf{Solution détaillée}

\begin{enumerate}
    \item \textbf{Synthèse en une étape :}
    \begin{itemize}
        \item Cible : \ce{CH3-CH(Br)-CH(CH3)2} (Brome sur C$_2$).
        \item Départ : \ce{CH2=CH-CH(CH3)2} (Double lien C$_1$=C$_2$).
        \item Analyse : Il faut ajouter H et Br sur le double lien. Le Br doit se retrouver sur C$_2$ (le carbone le moins hydrogéné du double lien, secondaire). L'H sur C$_1$ (le plus hydrogéné, primaire).
        \item Cela correspond \textbf{exactement à la règle de Markovnikov} pour l'addition de HBr.
        \item \textbf{Réactif :} \textbf{HBr} (chlorure d'hydrogène \textbf{incorrect}, donne le chloro). \textbf{Acide bromhydrique} (gaz ou solution aqueuse/concentrée). Pas de peroxyde (règle anti-Markovnikov).
        \item \textbf{Équation :}
        \[ \ce{CH2=CH-CH(CH3)2 + HBr -> CH3-CH(Br)-CH(CH3)2} \]
        \item \textbf{Conditions :} Milieu inerte (éther, CCl$_4$) ou aqueux, température ambiante, à l'abri de la lumière (pour éviter radicaux/peroxydes).
    \end{itemize}

    \item \textbf{Calcul de la masse théorique (correction : HBr) :}
    \begin{itemize}
        \item Masse alcène pur : $m_{pur} = \SI{10,0}{g} \times 0,95 = \SI{9,50}{g}$.
        \item $n_{alcène} = \frac{\SI{9,50}{g}}{\SI{70,13}{g.mol^{-1}}} = \SI{0,1355}{mol}$.
        \item HBr en excès $\rightarrow$ Alcène réactif limitant.
        \item Stoïchiométrie 1/1 $\rightarrow$ $n_{produit} = n_{alcène} = \SI{0,1355}{mol}$.
        \item Masse théorique $m_{th} = n_{produit} \times M_{produit} = \SI{0,1355}{mol} \times \SI{149,04}{g.mol^{-1}} = \boxed{\SI{20,19}{g}}$.
    \end{itemize}

    \item \textbf{Chiralité du produit :}
    \begin{itemize}
        \item Molécule produit : \ce{CH3-CH(Br)-CH(CH3)2}.
        \item Carbone C$_2$ (porteur du Br) : 4 substituants différents ?
            \begin{enumerate}
                \item H
                \item Br
                \item CH$_3$ (méthyle)
                \item CH(CH$_3$)$_2$ (isopropyle)
            \end{enumerate}
            Oui, 4 substituants différents. \textbf{C$_2$ est un carbone chiral (asymétrique)}.
        \item Carbone C$_3$ (CH) : Porte H, CH$_3$, CH$_3$, CH(Br)CH$_3$. Deux méthyles identiques $\rightarrow$ \textbf{Pas chiral}.
        \item \textbf{Mécanisme :} Addition ionique via carbocation secondaire sur C$_2$ (plan, \textit{sp}$^2$). Attaque de Br$^-$ sur les deux faces équiprobables.
        \item \textbf{Résultat :} Formation d'un \textbf{mélange racémique (50\% R, 50\% S)} du 2-bromo-3-méthylebutane. Pouvoir rotateur nul.
    \end{itemize}
\end{enumerate}
\textbf{Conclusion :} L'addition de HBr sur le 3-méthylebut-1-ène suit la règle de Markovnikov et donne le 2-bromo-3-méthylebutane sous forme de mélange racémique (création d'un centre chiral). Masse théorique : \SI{20,2}{g}.
\end{exoresolu}

\begin{attention}[Les 5 erreurs qui coûtent des points au Probatoire]
\begin{enumerate}
    \item \textbf{Confondre Markovnikov et Anti-Markovnikov :} Écrire que l'H va sur le carbone \textbf{le moins hydrogéné} (le plus substitué) pour l'addition de HX/H$_2$O. \textit{Rappel :} « Le riche devient plus riche » $\rightarrow$ H sur le C qui a le plus de H (le plus pauvre en C), X/OH sur l'autre (carbocation stable). \textbf{Perte : 2 à 3 pts sur l'équation et la justification.}
    
    \item \textbf{Oublier la stéréochimie \textit{anti} de l'addition du Br$_2$ :} Dessiner les deux atomes de brome du même côté (wedge/wedge ou dash/dash) sur un cycloalcène ou un alcène (Z)/(E) défini. \textit{Règle :} Ion halonium $\rightarrow$ attaque \textit{anti} obligatoire. \textbf{Perte : 1 à 2 pts sur la représentation de Cram.}
    
    \item \textbf{Numérotation UICPA incorrecte (Priorité au multiple lien) :} Numéroter la chaîne pour donner l'indice le plus bas à un substituant (ex: méthyle en 2) au détriment du double lien (qui se retrouve en 3 au lieu de 2). \textit{Règle absolue :} Le multiple lien a la priorité sur \textbf{tout} substituant pour la numérotation. \textbf{Perte : 1 pt sur le nom, souvent éliminatoire si la chaîne principale est mal choisie.}
    
    \item \textbf{Erreur sur l'hybridation / géométrie :} Dire qu'un carbone d'alcyne (C$\equiv$C) est \textit{sp}$^2$ (trigonal) ou qu'un carbone d'alcène est \textit{sp} (linéaire). Confondre le nombre de liaisons $\pi$ et le nombre de domaines $\sigma$. \textit{Réflexe :} Compter les liaisons $\sigma$ (simple + 1 par double/triple). \textbf{Perte : 1 à 2 pts sur la structure.}
    
    \item \textbf{Oublier la correction de la pression de vapeur d'eau (Dosage gaz / Préparation acétylène) :} Calculer le volume de gaz sec en utilisant $P_{atm}$ directement au lieu de $P_{gaz} = P_{atm} - P_{H_2O}^\circ(T)$. Ou utiliser $P_{H_2O}^\circ$ à une mauvaise température. \textbf{Perte : 2 pts sur le calcul final (volume ou quantité de matière).}
\end{enumerate}
\textbf{Astuce finale :} Relisez toujours votre équation de réaction : atomes conservés ? Charges ? États physiques ? Stéréochimie cohérente avec le mécanisme ?
\end{attention}

\newpage

\coursec{Exercices}

\courssub{Je vérifie mes connaissances}

\begin{exercice}[QCM : structure et nomenclature]
Pour chaque question, choisis la (ou les) bonne(s) réponse(s) en justifiant brièvement.
\begin{enumerate}
\item La formule générale d'un alcène non cyclique est :
\begin{enumerate}
\item $C_nH_{2n+2}$ \quad \item $C_nH_{2n}$ \quad \item $C_nH_{2n-2}$ \quad \item $C_nH_{2n-4}$
\end{enumerate}
\item Dans la molécule d'éthylène \ce{CH2=CH2}, chaque atome de carbone est :
\begin{enumerate}
\item tétragonal \quad \item trigonal \quad \item digonal \quad \item linéaire
\end{enumerate}
\item Le composé \ce{CH3-CH2-C#CH} se nomme :
\begin{enumerate}
\item but-1-yne \quad \item but-2-yne \quad \item prop-1-yne \quad \item but-1-ène
\end{enumerate}
\item L'angle de liaison $\widehat{H-C-H}$ dans l'éthylène vaut environ :
\begin{enumerate}
\item \ang{90} \quad \item \ang{109} \quad \item \ang{120} \quad \item \ang{180}
\end{enumerate}
\end{enumerate}
\end{exercice}

\begin{exercice}[Vrai ou faux justifié]
Réponds par \textbf{vrai} ou \textbf{faux} en justifiant chaque réponse par une phrase ou une équation-bilan.
\begin{enumerate}
\item Tous les alcènes présentent l'isomérie Z/E.
\item Le dibrome \ce{Br2} permet de distinguer un alcane d'un alcène.
\item Selon la règle de Markovnikov, lors de l'addition de \ce{HCl} sur un alcène dissymétrique, l'atome d'hydrogène se fixe sur le carbone de la double liaison le moins hydrogéné.
\item L'acétylène \ce{C2H2} est obtenu au laboratoire par action de l'eau sur le carbure de calcium \ce{CaC2}.
\end{enumerate}
\end{exercice}

\begin{exercice}[Texte à trous]
Complète le texte suivant avec les mots ou expressions : \emph{trigonal}, \emph{digonal}, \emph{addition}, \emph{polymérisation}, \emph{double}, \emph{triple}, \emph{\ang{180}}, \emph{\ang{120}}.

Dans les alcènes, les deux atomes de carbone de la liaison \dotfill{} sont des carbones \dotfill{} : les angles de liaison autour d'eux valent environ \dotfill{}. Dans les alcynes, les atomes de carbone de la liaison \dotfill{} sont des carbones \dotfill{} : la molécule d'éthyne est linéaire, avec des angles de \dotfill{}. La réaction caractéristique des alcènes et des alcynes est la réaction d'\dotfill{}, au cours de laquelle la liaison multiple se transforme. L'enchaînement de nombreuses molécules de monomère insaturé constitue une réaction de \dotfill{}.
\end{exercice}

\begin{exercice}[Définitions]
Définis en une ou deux phrases chacun des termes suivants, en illustrant par un exemple :
\begin{enumerate}
\item alcène ; \item alcyne ; \item isomérie Z/E ; \item règle de Markovnikov ; \item monomère et polymère.
\end{enumerate}
\end{exercice}

\courssub{Je m'entraîne}

\begin{exoresolu}[Nomenclature et isomérie Z/E]
Nomme les composés suivants et précise ceux qui présentent l'isomérie Z/E (dans ce cas, représente les deux stéréoisomères) :
\begin{enumerate}
\item \ce{CH3-CH=CH-CH3} ; \quad
\item \chemfig{CH_2=C(-[2]CH_3)(-[6]CH_3)} ; \quad
\item \ce{CH3-CH2-C#CH} ; \quad
\item \chemfig{CH_3-CH=C(-[2]CH_3)-CH_2-CH_3}.
\end{enumerate}
\tcblower
\textbf{Solution détaillée :}
\begin{enumerate}
\item \textbf{But-2-ène.} Les deux carbones de la double liaison portent chacun un H et un $\ce{CH3}$ (deux substituants différents sur chaque carbone) $\rightarrow$ \textbf{Isomérie Z/E possible}.
  \begin{center}
  \begin{tabular}{cc}
  \chemfig{CH_3-[:30]=_[:30]CH-[:-30]CH_3} & \chemfig{CH_3-[:30]=_[:-30]CH-[:30]CH_3} \\
  \textbf{(Z)-but-2-ène} & \textbf{(E)-but-2-ène} \\
  (groupes \ce{CH3} du même côté) & (groupes \ce{CH3} côtés opposés)
  \end{tabular}
  \end{center}
\item \textbf{2-méthylprop-1-ène} (nom usuel : isobutène). Le carbone C1 de la double liaison porte deux atomes d'hydrogène $\rightarrow$ \textbf{Pas d'isomérie Z/E}.
\item \textbf{But-1-yne.} C'est un alcyne : la triple liaison impose une géométrie linéaire (digonal) $\rightarrow$ \textbf{Pas d'isomérie Z/E}.
\item \textbf{Pent-2-ène.} Chaîne principale à 5 C, double liaison en 2. Carbone 2 : H et $\ce{CH3}$. Carbone 3 : $\ce{CH3}$ et $\ce{CH2CH3}$. Deux substituants différents sur chaque carbone $\rightarrow$ \textbf{Isomérie Z/E possible}.
  \begin{center}
  \begin{tabular}{cc}
  \chemfig{CH_3-[:30]=_[:30]C(-[2]CH_3)-[:-30]CH_2-CH_3} & \chemfig{CH_3-[:30]=_[:-30]C(-[2]CH_3)-[:30]CH_2-CH_3} \\
  \textbf{(Z)-pent-2-ène} & \textbf{(E)-pent-2-ène} \\
  (priorités \ce{CH3} et \ce{CH2CH3} même côté) & (priorités côtés opposés)
  \end{tabular}
  \end{center}
\end{enumerate}
\end{exoresolu}

\begin{exoresolu}[Addition de HCl sur le but-1-ène]
On fait réagir le chlorure d'hydrogène \ce{HCl} sur le but-1-ène.
\begin{enumerate}
\item Écris l'équation-bilan de la réaction conduisant au produit majoritaire.
\item Nomme les deux produits possibles et précise lequel est majoritaire en justifiant par la règle de Markovnikov.
\item Quel type de réaction s'agit-il ? Cite un test expérimental permettant de vérifier que le but-1-ène a disparu.
\end{enumerate}
\tcblower
\textbf{Solution détaillée :}
\begin{enumerate}
\item But-1-ène : $\ce{CH2=CH-CH2-CH3}$. Addition de \ce{HCl} selon Markovnikov : H sur C1 (2 H), Cl sur C2 (1 H).
  \[ \ce{CH2=CH-CH2-CH3 + HCl -> CH3-CHCl-CH2-CH3} \]
\item Produits possibles :
  \begin{itemize}
  \item \textbf{2-chlorobutane} ($\ce{CH3-CHCl-CH2-CH3}$) : \textbf{Produit majoritaire} (règle de Markovnikov : H se fixe sur le carbone le plus hydrogéné, C1).
  \item 1-chlorobutane ($\ce{CH2Cl-CH2-CH2-CH3}$) : produit minoritaire (anti-Markovnikov).
  \end{itemize}
\item C'est une \textbf{réaction d'addition électrophile}.
  Test : Faire bouillonner le gaz dans de l'\textbf{eau de brome} (ou solution de dibrome dans le tétrachlorure de carbone / cyclohexane). \textbf{Observation :} Décoloration rapide de la couleur orange/brune de \ce{Br2} (formation de \ce{CH2Br-CHBr-CH2-CH3}, incolore).
\end{enumerate}
\end{exoresolu}

\begin{exoresolu}[Polymérisation du chlorure de vinyle]
Le chlorure de vinyle \ce{CH2=CHCl} est le monomère du polychlorure de vinyle (PVC), utilisé pour fabriquer les tuyaux d'adduction d'eau.
\begin{enumerate}
\item Écris l'équation de la réaction de polymérisation du chlorure de vinyle.
\item Calcule la masse molaire du motif du polymère.
\item Un échantillon de PVC a une masse molaire moyenne de \SI{62500}{g.mol^{-1}}. Calcule son indice de polymérisation moyen (nombre de motifs).
\end{enumerate}
\emph{Données :} $M(\ce{C}) = \SI{12}{g.mol^{-1}}$ ; $M(\ce{H}) = \SI{1}{g.mol^{-1}}$ ; $M(\ce{Cl}) = \SI{35,5}{g.mol^{-1}}$.
\tcblower
\textbf{Solution détaillée :}
\begin{enumerate}
\item Polymérisation par addition (ouverture de la double liaison) :
  \[ n\,\ce{CH2=CHCl} \xrightarrow{\text{initiateur}} \ce{-[CH2-CHCl]-_n} \]
  (On écrit souvent : $n \ce{CH2=CHCl} \rightarrow \ce{(-CH2-CHCl-)_n}$)
\item Motif répétitif : $\ce{-CH2-CHCl-}$.
  $M_{\text{motif}} = 2 \times M(\ce{C}) + 3 \times M(\ce{H}) + M(\ce{Cl}) = 2\times12 + 3\times1 + 35,5 = \SI{62,5}{g.mol^{-1}}$.
\item Indice de polymérisation moyen $\bar{n} = \dfrac{M_{\text{polymère}}}{M_{\text{motif}}} = \dfrac{62500}{62,5} = \textbf{1000}$.
\end{enumerate}
\end{exoresolu}

\begin{exoresolu}[Combustion d'un alcène]
La combustion complète de \SI{0,10}{mol} d'un alcène gazeux dans le dioxygène fournit \SI{13,2}{g} de dioxyde de carbone et de l'eau.
\begin{enumerate}
\item Écris l'équation générale de combustion complète d'un alcène $C_nH_{2n}$.
\item Détermine la valeur de $n$ et la formule brute de l'alcène.
\item Écris l'équation de combustion équilibrée de cet alcène.
\end{enumerate}
\emph{Données :} $M(\ce{C}) = \SI{12}{g.mol^{-1}}$ ; $M(\ce{O}) = \SI{16}{g.mol^{-1}}$.
\tcblower
\textbf{Solution détaillée :}
\begin{enumerate}
\item $\ce{C_nH_{2n} + \frac{3n}{2} O2 -> n CO2 + n H2O}$ (ou multipliée par 2 : $\ce{2 C_nH_{2n} + 3n O2 -> 2n CO2 + 2n H2O}$).
\item $n(\ce{CO2}) = \dfrac{m(\ce{CO2})}{M(\ce{CO2})} = \dfrac{13,2}{44,0} = \SI{0,30}{mol}$.
  D'après l'équation : 1 mol alcène $\rightarrow$ $n$ mol $\ce{CO2}$.
  $n_{\text{alcène}} \times n = n(\ce{CO2}) \Rightarrow 0,10 \times n = 0,30 \Rightarrow \textbf{n = 3}$.
  Formule brute : \textbf{$\ce{C3H6}$} (propène ou cyclopropane ; « alcène gazeux » non cyclique $\rightarrow$ propène).
\item $\ce{C3H6 + \frac{9}{2} O2 -> 3 CO2 + 3 H2O}$ \quad (ou $\ce{2 C3H6 + 9 O2 -> 6 CO2 + 6 H2O}$).
\end{enumerate}
\end{exoresolu}

\courssub{Je me prépare au Probatoire}

\begin{exoresolu}[L'hydrocarbure mystérieux]
Un hydrocarbure gazeux A a une densité de vapeur $d = 0{,}97$ par rapport à l'air. Un échantillon de A décolore rapidement l'eau de brome, mais ne réagit pas avec le chlorure d'hydrogène en l'absence de catalyseur adapté.
\begin{enumerate}
\item Rappelle la définition de la densité de vapeur d'un gaz et montre que la masse molaire de A vaut environ \SI{28}{g.mol^{-1}}. On donne $M_{\text{air}} = \SI{29}{g.mol^{-1}}$.
\item Que traduit la décoloration de l'eau de brome ? Écris l'équation-bilan correspondante avec la formule brute de A.
\item Détermine la formule brute de A sachant qu'il s'agit d'un alcène. Écris sa formule semi-développée et nomme-le.
\item A présente-t-il l'isomérie Z/E ? Justifie.
\item On fait réagir A avec le dihydrogène en présence de nickel. Écris l'équation-bilan et nomme le produit obtenu.
\item La combustion complète de \SI{5,6}{g} de A dans le dioxygène produit un volume $V$ de dioxyde de carbone mesuré dans les conditions où $V_m = \SI{24}{L.mol^{-1}}$. Calcule $V$.
\end{enumerate}
\emph{Données :} $M(\ce{C}) = \SI{12}{g.mol^{-1}}$ ; $M(\ce{H}) = \SI{1}{g.mol^{-1}}$.
\tcblower
\textbf{Solution détaillée :}
\begin{enumerate}
\item Densité de vapeur $d = \dfrac{M_{\text{gaz}}}{M_{\text{air}}}$ (même $T, P$).
  $M_A = d \times M_{\text{air}} = 0,97 \times 29 = \SI{28,13}{g.mol^{-1}} \approx \SI{28}{g.mol^{-1}}$.
\item La décoloration de l'eau de brome traduit une \textbf{réaction d'addition} du dibrome $\ce{Br2}$ sur une liaison multiple (insaturation). L'alcène A réagit : $\ce{C2H4 + Br2 -> C2H4Br2}$ (1,2-dibromoéthane).
\item Masse molaire $\approx 28 \text{ g/mol}$. Alcène : formule générale $\ce{C_nH_{2n}}$, $M = 14n$.
  $14n = 28 \Rightarrow n=2$. Formule brute : \textbf{$\ce{C2H4}$}.
  Formule semi-développée : \textbf{$\ce{CH2=CH2}$}. Nom : \textbf{éthène} (nom usuel : éthylène).
\item \textbf{Non.} Les deux carbones de la double liaison portent chacun deux atomes d'hydrogène (deux substituants identiques sur chaque carbone). Condition d'existence de l'isomérie Z/E non remplie.
\item Hydrogénation (addition de \ce{H2}) catalysée par le nickel (Ni) :
  \[ \ce{CH2=CH2 + H2 ->[Ni] CH3-CH3} \]
  Produit : \textbf{éthane}.
\item $n_A = \dfrac{5,6}{28} = \SI{0,20}{mol}$.
  Équation combustion : $\ce{C2H4 + 3 O2 -> 2 CO2 + 2 H2O}$.
  $n(\ce{CO2}) = 2 \times n_A = \SI{0,40}{mol}$.
  $V(\ce{CO2}) = n(\ce{CO2}) \times V_m = 0,40 \times 24 = \textbf{\SI{9,6}{L}}$.
\end{enumerate}
\end{exoresolu}

\begin{exoresolu}[L'acétylène du soudeur]
Un atelier de soudure de Douala utilise le chalumeau oxyacétylénique. L'acétylène \ce{C2H2} est préparé au laboratoire par action de l'eau sur le carbure de calcium \ce{CaC2} selon une réaction qui produit également de la chaux éteinte \ce{Ca(OH)2}.
\begin{enumerate}
\item Écris l'équation-bilan de la préparation de l'acétylène.
\item Schématise et légende le dispositif expérimental de préparation et de recueil de l'acétylène au laboratoire (on précisera le mode de recueil du gaz).
\item Propose un test chimique simple permettant de prouver que le gaz obtenu est insaturé ; décris l'observation attendue et écris l'équation-bilan du test.
\item Le soudeur consomme \SI{48}{L} d'acétylène mesurés dans les conditions où $V_m = \SI{24}{L.mol^{-1}}$. Calcule la quantité de matière d'acétylène utilisée.
\item En déduis la masse de carbure de calcium pur qu'il a fallu décomposer.
\item Calcule la masse de chaux éteinte \ce{Ca(OH)2} formée au cours de cette préparation.
\item En réalité, le carbure de calcium commercial n'est pur qu'à \SI{80}{\%} en masse. Quelle masse de carbure commercial faut-il utiliser ?
\end{enumerate}
\emph{Données :} $M(\ce{Ca}) = \SI{40}{g.mol^{-1}}$ ; $M(\ce{C}) = \SI{12}{g.mol^{-1}}$ ; $M(\ce{O}) = \SI{16}{g.mol^{-1}}$ ; $M(\ce{H}) = \SI{1}{g.mol^{-1}}$.
\tcblower
\textbf{Solution détaillée :}
\begin{enumerate}
\item $\ce{CaC2(s) + 2 H2O(l) -> Ca(OH)2(aq) + C2H2(g)}$
\item \textbf{Dispositif :} Voir schéma ci-dessous. C'est un montage à « goutte-à-goutte » (entonnoir à robinet ou burette) vers un creuset/ballon contenant le solide \ce{CaC2}. Le gaz est recueilli par \textbf{déplacement d'eau} (bac pneumatique + tube à essai ou jauge graduée renversée) car l'acétylène est peu soluble dans l'eau et moins dense que l'air.
  \begin{popfigure}
  \begin{tikzpicture}[scale=0.9, >=Stealth]
    % Support
    \draw[thick] (0,0) -- (10,0);
    % Entonnoir à robinet (burette simplifiée)
    \draw[thick] (2,4) -- (2,2.5) -- (1.5,2) -- (2.5,2) -- (2,2.5);
    \draw[thick] (1.5,2) -- (1.5,1.5) -- (2.5,1.5) -- (2.5,2);
    \draw[thick] (2,1.5) -- (2,1) node[midway, right] {Robinet};
    \draw[fill=popblueL] (2,3.5) rectangle (2.4,4) node[midway, above] {Eau};
    % Ballon/Creuset
    \draw[thick] (5,1) .. controls (5,0.5) and (7,0.5) .. (7,1);
    \draw[thick] (5,1) -- (5,2.5);
    \draw[thick] (7,1) -- (7,2.5);
    \draw[thick] (5,2.5) -- (7,2.5);
    \draw[fill=poporangeL] (5.2,1.1) rectangle (6.8,1.8) node[midway] {$\ce{CaC2}$};
    % Tubulure entrée
    \draw[thick] (2,1) -- (5,1.8);
    % Tubulure sortie
    \draw[thick] (7,2.5) -- (8,2.5) -- (8,0.5);
    % Bac pneumatique
    \draw[thick] (7.5,-0.5) -- (7.5,0) -- (9.5,0) -- (9.5,-0.5);
    \draw[fill=popblueL!30] (7.5,0) -- (9.5,0) -- (9.5,-0.5) -- (7.5,-0.5);
    % Tube à essai renversé
    \draw[thick] (8.2,0) -- (8.2,2) -- (8.8,2) -- (8.8,0);
    \draw[fill=popblueL!30] (8.2,0.2) -- (8.8,0.2) -- (8.8,0) -- (8.2,0);
    \draw[<->] (8.2,2) -- (8.2,0.2) node[midway, left] {Gaz};
    % Légendes
    \node[align=center] at (2,4.5) {Entonnoir à robinet\\(Eau distillée)};
    \node[align=center] at (6,3) {Creuset / Ballon\\(\ce{CaC2} solide)};
    \node[align=center] at (9.5,1) {Recueil par\\déplacement d'eau};
    \node at (8.5,2.3) {$\ce{C2H2}$};
    \draw[->, popdark] (2,3.5) -- (3,3);
    \draw[->, popdark] (7,2.5) -- (7.5,3);
  \end{tikzpicture}
  \legende{Dispositif de préparation de l'acétylène par action de l'eau sur le carbure de calcium (recueil par déplacement d'eau).}
  \end{popfigure}
\item Test : Faire buller le gaz dans de l'\textbf{eau de brome} (ou solution de \ce{KMnO4} acide diluée).
  \textbf{Observation :} Décoloration rapide de la couleur orange (eau de brome) ou violette (\ce{KMnO4}).
  \textbf{Équation-bilan (eau de brome) :} $\ce{C2H2 + 2 Br2 -> C2H2Br4}$ (tétrabromoéthane, incolore).
\item $n(\ce{C2H2}) = \dfrac{V}{V_m} = \dfrac{48}{24} = \textbf{\SI{2,0}{mol}}$.
\item Équation : 1 mol $\ce{CaC2} \rightarrow$ 1 mol $\ce{C2H2}$. Il faut \SI{2,0}{mol} de $\ce{CaC2}$ pur.
  $M(\ce{CaC2}) = 40 + 2\times12 = \SI{64}{g.mol^{-1}}$.
  $m_{\text{pur}}(\ce{CaC2}) = 2,0 \times 64 = \textbf{\SI{128}{g}}$.
\item Équation : 1 mol $\ce{CaC2} \rightarrow$ 1 mol $\ce{Ca(OH)2}$. $n(\ce{Ca(OH)2}) = \SI{2,0}{mol}$.
  $M(\ce{Ca(OH)2}) = 40 + 2\times(16+1) = \SI{74}{g.mol^{-1}}$.
  $m(\ce{Ca(OH)2}) = 2,0 \times 74 = \textbf{\SI{148}{g}}$.
\item Pureté 80\% : $m_{\text{commercial}} = \dfrac{m_{\text{pur}}}{0,80} = \dfrac{128}{0,80} = \textbf{\SI{160}{g}}$.
\end{enumerate}
\end{exoresolu}

\begin{exoresolu}[Du pétrole de la SONARA au plastique]
Le vapocraquage des fractions pétrolières fournit l'éthylène \ce{C2H4}, matière première de l'industrie des plastiques. On étudie quelques transformations de l'éthylène.
\begin{enumerate}
\item Donne la formule développée de l'éthylène et précise la géométrie autour de chaque atome de carbone (type de carbone, angles de liaison).
\item L'éthylène subit une hydratation en présence d'un catalyseur acide pour donner un composé B utilisé comme désinfectant dans les hôpitaux. Écris l'équation-bilan, nomme B et précise sa famille chimique.
\item L'éthylène réagit avec le dichlore pour donner un composé C. Écris l'équation-bilan et nomme C.
\item L'éthylène polymérise pour donner le polyéthylène (PE), utilisé pour fabriquer les sachets et les jerricans. Écris l'équation de polymérisation.
\item Un échantillon de polyéthylène a un indice de polymérisation moyen $n = 2000$. Calcule sa masse molaire moyenne, puis la masse d'un échantillon contenant \SI{0,050}{mol} de macromolécules.
\item La combustion complète de \SI{1,0}{kg} de polyéthylène rejette du dioxyde de carbone. En assimilant le polymère à une succession de motifs \ce{CH2}, calcule la masse de \ce{CO2} rejetée. Commente l'impact environnemental de l'incinération des sachets plastiques.
\end{enumerate}
\emph{Données :} $M(\ce{C}) = \SI{12}{g.mol^{-1}}$ ; $M(\ce{H}) = \SI{1}{g.mol^{-1}}$ ; $M(\ce{O}) = \SI{16}{g.mol^{-1}}$.
\tcblower
\textbf{Solution détaillée :}
\begin{enumerate}
\item Formule développée : $\ce{H2C=CH2}$ (ou $\chemfig{H_2C=CH_2}$).
  Chaque carbone est \textbf{trigonal} (hybridation $sp^2$), géométrie plane, angles de liaison $\approx \mathbf{\ang{120}}$.
\item Hydratation (addition d'eau, catalyse acide $\ce{H3O+}$ / $\ce{H2SO4}$ concentré) :
  \[ \ce{CH2=CH2 + H2O ->[H+/cat] CH3-CH2-OH} \]
  Composé B : \textbf{éthanol} (alcool éthylique). Famille : \textbf{alcool} (alcool primaire, saturé). Utilisé comme antiseptique (désinfectant) à \SI{70}{\%} massique.
\item Addition de dichlore :
  \[ \ce{CH2=CH2 + Cl2 -> CH2Cl-CH2Cl} \]
  Composé C : \textbf{1,2-dichloroéthane} (dichlorure d'éthylène).
\item Polymérisation :
  \[ n\,\ce{CH2=CH2} \xrightarrow{\text{pression, temp, catalyseur}} \ce{-[CH2-CH2]-_n} \]
  (Polyéthylène, motif $\ce{-CH2-CH2-}$ ou $\ce{CH2}$).
\item Motif $\ce{C2H4}$ : $M_{\text{motif}} = 2\times12 + 4\times1 = \SI{28}{g.mol^{-1}}$.
  $M_{\text{PE}} = n \times M_{\text{motif}} = 2000 \times 28 = \SI{56000}{g.mol^{-1}} = \SI{56}{kg.mol^{-1}}$.
  Masse de \SI{0,050}{mol} de macromolécules : $m = n \times M = 0,050 \times 56000 = \textbf{\SI{2800}{g}} = \textbf{\SI{2,8}{kg}}$.
\item Motif $\ce{CH2}$ : $M_{\text{motif}} = \SI{14}{g.mol^{-1}}$.
  Quantité de matière de motifs dans 1,0 kg = \SI{1000}{g} : $n_{\text{motifs}} = \dfrac{1000}{14} \approx \SI{71,43}{mol}$.
  Combustion du motif : $\ce{CH2 + 3/2 O2 -> CO2 + H2O}$. 1 mol motif $\rightarrow$ 1 mol $\ce{CO2}$.
  $n(\ce{CO2}) = \SI{71,43}{mol}$.
  $m(\ce{CO2}) = 71,43 \times 44 = \textbf{\SI{3143}{g} \approx \SI{3,14}{kg}}$.
  \textbf{Commentaire :} L'incinération de 1 kg de sachets plastiques (PE) rejette plus de 3 kg de $\ce{CO2}$, gaz à effet de serre majeur contribuant au réchauffement climatique. De plus, une combustion incomplète (feu ouvert) produit des suies, du monoxyde de carbone \ce{CO} et des composés toxiques (dioxines si présence de chlore, bien que le PE n'en contienne pas). La gestion des déchets plastiques par recyclage ou valorisation énergétique contrôlée (avec filtration des fumées) est préférable au brûlage sauvage.
\end{enumerate}
\end{exoresolu}

\newpage

\coursec{Corrigés des exercices}

\begin{corrige}[Exercice 1 — QCM : structure et nomenclature]
\begin{enumerate}
\item \textbf{Bonne réponse : b) $C_nH_{2n}$.} Un alcène non cyclique possède une double liaison \ce{C=C} ; par rapport à l'alcane de même chaîne ($C_nH_{2n+2}$), la formation de la double liaison supprime 2 atomes d'hydrogène : $2n+2-2 = 2n$.
\item \textbf{Bonne réponse : b) trigonal.} Chaque carbone de l'éthylène est hybridé $sp^2$ : il est lié à 3 atomes (2 H et 1 C) disposés dans un plan ; c'est un \cle{carbone trigonal}.
\item \textbf{Bonne réponse : a) but-1-yne.} La chaîne compte 4 atomes de carbone (préfixe \emph{but-}) et la triple liaison est en position 1 si l'on numérote depuis l'extrémité la plus proche de la triple liaison : suffixe \emph{-1-yne}.
\item \textbf{Bonne réponse : c) \ang{120}.} Autour d'un carbone trigonal $sp^2$, la géométrie est plane triangulaire : angles voisins de \ang{120}.
\end{enumerate}
\textbf{Points de vigilance :} ne pas confondre carbone \emph{trigonal} (alcène, $sp^2$, \ang{120}) et carbone \emph{digonal} (alcyne, $sp$, \ang{180}) ; ne pas oublier que la numérotation d'un alcyne donne toujours à la triple liaison l'indice le plus petit.
\end{corrige}

\begin{corrige}[Exercice 2 — Vrai ou faux justifié]
\begin{enumerate}
\item \textbf{Faux.} L'isomérie Z/E n'existe que si \emph{chacun} des deux carbones de la double liaison porte deux substituants \emph{différents}. Contre-exemples : \ce{CH2=CH2} (éthène) et \ce{CH2=CH-CH3} (propène) n'ont pas d'isomères Z/E car l'un des carbones porte deux atomes d'hydrogène identiques.
\item \textbf{Vrai.} Un alcène décolore instantanément l'eau de brome par addition : $\ce{CH2=CH2 + Br2 -> CH2Br-CH2Br}$ (incolore). Un alcane ne réagit pas avec \ce{Br2} à froid et à l'obscurité (substitution lente, photochimique). C'est le \cle{test d'insaturation}.
\item \textbf{Faux.} C'est l'inverse : selon la règle de Markovnikov, l'hydrogène de \ce{HCl} se fixe sur le carbone de la double liaison \emph{le plus hydrogéné} ; le chlore se fixe sur le carbone le plus substitué. Exemple : $\ce{CH2=CH-CH3 + HCl -> CH3-CHCl-CH3}$ (2-chloropropane majoritaire).
\item \textbf{Vrai.} L'action de l'eau sur le carbure de calcium produit l'éthyne (acétylène) et la chaux éteinte : \[ \ce{CaC2(s) + 2 H2O(l) -> Ca(OH)2(aq) + C2H2(g)} \]
\end{enumerate}
\textbf{Points de vigilance :} dans la règle de Markovnikov, repérer d'abord le carbone le plus hydrogéné ; le mot « décoloration » doit être accompagné du nom du réactif (eau de brome) et du type de réaction (addition).
\end{corrige}

\begin{corrige}[Exercice 3 — Texte à trous]
Dans les alcènes, les deux atomes de carbone de la liaison \textbf{double} sont des carbones \textbf{trigonaux} : les angles de liaison autour d'eux valent environ \textbf{\ang{120}}. Dans les alcynes, les atomes de carbone de la liaison \textbf{triple} sont des carbones \textbf{digonaux} : la molécule d'éthyne est linéaire, avec des angles de \textbf{\ang{180}}. La réaction caractéristique des alcènes et des alcynes est la réaction d'\textbf{addition}, au cours de laquelle la liaison multiple se transforme. L'enchaînement de nombreuses molécules de monomère insaturé constitue une réaction de \textbf{polymérisation}.

\textbf{Points de vigilance :} « trigonal » qualifie le carbone $sp^2$ (3 voisins), « digonal » le carbone $sp$ (2 voisins) ; la polymérisation des alcènes est une \emph{polyaddition} : il n'y a pas d'élimination de petite molécule.
\end{corrige}

\begin{corrige}[Exercice 4 — Définitions]
\begin{enumerate}
\item \textbf{Alcène :} hydrocarbure insaturé non cyclique comportant une double liaison \ce{C=C}, de formule générale $C_nH_{2n}$ ($n \geq 2$). Exemple : l'éthène \ce{CH2=CH2}.
\item \textbf{Alcyne :} hydrocarbure insaturé non cyclique comportant une triple liaison \ce{C#C}, de formule générale $C_nH_{2n-2}$ ($n \geq 2$). Exemple : l'éthyne \ce{HC#CH} (acétylène).
\item \textbf{Isomérie Z/E :} stéréoisomérie de configuration due à l'absence de rotation autour de la double liaison \ce{C=C} ; elle existe lorsque chaque carbone de la double liaison porte deux substituants différents. L'isomère \textbf{Z} a les groupes prioritaires du même côté, l'isomère \textbf{E} de part et d'autre. Exemple : le but-2-ène \ce{CH3-CH=CH-CH3}.
\item \textbf{Règle de Markovnikov :} lors de l'addition d'un hydracide \ce{HX} sur un alcène dissymétrique, l'atome d'hydrogène se fixe préférentiellement sur le carbone de la double liaison le plus hydrogéné. Exemple : $\ce{CH2=CH-CH3 + HCl -> CH3-CHCl-CH3}$.
\item \textbf{Monomère :} petite molécule insaturée capable de s'enchaîner à de nombreuses molécules identiques. \textbf{Polymère :} macromolécule formée par la répétition d'un motif issu du monomère. Exemple : le chloroéthène \ce{CH2=CHCl} (monomère) donne le PVC \ce{(-CH2-CHCl-)_n} (polymère).
\end{enumerate}
\textbf{Points de vigilance :} une définition sans exemple est incomplète au Probatoire ; préciser toujours la formule générale avec la condition $n \geq 2$.
\end{corrige}

\begin{corrige}[Exercice 5 — Nomenclature et isomérie Z/E]
\begin{enumerate}
\item \ce{CH3-CH=CH-CH3} : chaîne de 4 C, double liaison en position 2 $\rightarrow$ \textbf{but-2-ène}. Chaque carbone de la double liaison porte H et \ce{CH3} (deux substituants différents) : \textbf{isomérie Z/E}.
\begin{center}
\begin{tabular}{cc}
\chemfig{CH_3-[:30]=_[:30]CH-[:-30]CH_3} & \chemfig{CH_3-[:30]=_[:-30]CH-[:30]CH_3} \\
\textbf{(Z)-but-2-ène} & \textbf{(E)-but-2-ène} \\
(\ce{CH3} du même côté) & (\ce{CH3} de part et d'autre)
\end{tabular}
\end{center}
\item \chemfig{CH_2=C(-[2]CH_3)(-[6]CH_3)} : chaîne de 3 C avec un méthyle en position 2 $\rightarrow$ \textbf{2-méthylprop-1-ène}. Le carbone 1 porte deux H identiques : \textbf{pas d'isomérie Z/E}.
\item \ce{CH3-CH2-C#CH} : \textbf{but-1-yne}. Les alcynes sont linéaires au niveau de la triple liaison (carbones digonaux) : \textbf{pas d'isomérie Z/E}.
\item \chemfig{CH_3-CH=C(-[2]CH_3)-CH_2-CH_3} : chaîne la plus longue contenant la double liaison : 5 C, insaturation en 2, méthyle en 2 $\rightarrow$ \textbf{2-méthylpent-2-ène}. C2 porte H et \ce{CH3} ; C3 porte \ce{CH3} et \ce{C2H5} : \textbf{isomérie Z/E}.
\begin{center}
\begin{tabular}{cc}
\chemfig{CH_3-[:30]=_[:30]C(-[2]CH_3)-[:-30]CH_2-CH_3} & \chemfig{CH_3-[:30]=_[:-30]C(-[2]CH_3)-[:30]CH_2-CH_3} \\
\textbf{(Z)-2-méthylpent-2-ène} & \textbf{(E)-2-méthylpent-2-ène}
\end{tabular}
\end{center}
\end{enumerate}
\textbf{Points de vigilance :} choisir la chaîne la plus longue \emph{contenant la double liaison} (5 C et non 4 pour le composé 4) ; vérifier la condition Z/E sur \emph{chaque} carbone de la double liaison avant de conclure.
\end{corrige}

\begin{corrige}[Exercice 6 — Addition de HCl sur le but-1-ène]
\begin{enumerate}
\item Le but-1-ène est \ce{CH2=CH-CH2-CH3}. C1 porte 2 H (le plus hydrogéné), C2 porte 1 H. Selon Markovnikov, H se fixe sur C1 et Cl sur C2 :
\[ \ce{CH2=CH-CH2-CH3 + HCl -> CH3-CHCl-CH2-CH3} \]
\item Les deux produits possibles :
\begin{itemize}
\item \textbf{2-chlorobutane} \ce{CH3-CHCl-CH2-CH3} : produit \textbf{majoritaire} (règle de Markovnikov : H sur le carbone le plus hydrogéné) ;
\item 1-chlorobutane \ce{CH2Cl-CH2-CH2-CH3} : produit minoritaire.
\end{itemize}
\item Il s'agit d'une \textbf{réaction d'addition} (électrophile) sur la double liaison. \textbf{Test de disparition de l'alcène :} faire buller le gaz restant dans l'eau de brome ; si le but-1-ène a disparu, l'eau de brome \textbf{reste orange} (pas de décoloration). Inversement, tant qu'il reste de l'alcène, on observe la décoloration : $\ce{CH2=CH-CH2-CH3 + Br2 -> CH2Br-CHBr-CH2-CH3}$.
\end{enumerate}
\textbf{Points de vigilance :} équilibrer l'équation (1 molécule de HCl pour 1 molécule d'alcène) ; ne pas écrire le produit minoritaire comme produit unique ; le test à l'eau de brome caractérise l'insaturation, pas le produit chloré.
\end{corrige}

\begin{corrige}[Exercice 7 — Polymérisation du chlorure de vinyle]
\begin{enumerate}
\item Équation de polymérisation (polyaddition, ouverture de la double liaison) :
\[ n\,\ce{CH2=CHCl} \ce{->[\text{initiateur}]} \ce{(-CH2-CHCl-)_n} \]
\item Motif \ce{-CH2-CHCl-} de formule brute \ce{C2H3Cl} :
\[ M_{\text{motif}} = 2 \times 12 + 3 \times 1 + 35{,}5 = 24 + 3 + 35{,}5 = \SI{62,5}{g.mol^{-1}} \]
\item Indice de polymérisation moyen :
\[ \bar{n} = \frac{M_{\text{polymère}}}{M_{\text{motif}}} = \frac{62500}{62{,}5} = \boxed{1000} \]
La macromolécule comporte en moyenne 1000 motifs.
\end{enumerate}
\textbf{Vérification :} $62{,}5 \times 1000 = 62500$ \checkmark. \textbf{Points de vigilance :} le motif contient 3 H (2 sur le \ce{CH2}, 1 sur le \ce{CHCl}) et non 4 ; l'indice de polymérisation est un nombre sans unité.
\end{corrige}

\begin{corrige}[Exercice 8 — Combustion d'un alcène]
\begin{enumerate}
\item Équation générale de combustion complète d'un alcène :
\[ \ce{C_nH_{2n} + \frac{3n}{2} O2 -> n CO2 + n H2O} \]
Vérification : C : $n = n$ \checkmark ; H : $2n = 2n$ \checkmark ; O : $2 \times \frac{3n}{2} = 3n = 2n + n$ \checkmark.
\item Quantité de \ce{CO2} : $M(\ce{CO2}) = 12 + 2 \times 16 = \SI{44}{g.mol^{-1}}$.
\[ n(\ce{CO2}) = \frac{13{,}2}{44} = \SI{0,30}{mol} \]
D'après l'équation, 1 mol d'alcène donne $n$ mol de \ce{CO2} :
\[ 0{,}10 \times n = 0{,}30 \quad \Rightarrow \quad n = 3 \]
Formule brute : \textbf{\ce{C3H6}} (alcène non cyclique : le propène \ce{CH2=CH-CH3}).
\item Équation équilibrée :
\[ \ce{C3H6 + \frac{9}{2} O2 -> 3 CO2 + 3 H2O} \quad \text{ou} \quad \ce{2 C3H6 + 9 O2 -> 6 CO2 + 6 H2O} \]
\end{enumerate}
\textbf{Points de vigilance :} utiliser $M(\ce{CO2})$ et non $M(\ce{C})$ pour convertir la masse de \ce{CO2} ; une combustion \emph{complète} ne produit que \ce{CO2} et \ce{H2O} (jamais \ce{CO} ni \ce{C}).
\end{corrige}

\begin{corrige}[Exercice 9 — L'hydrocarbure mystérieux]
\begin{enumerate}
\item La densité de vapeur d'un gaz est le rapport de sa masse molaire à celle de l'air (mêmes conditions de $T$ et $P$) : $d = \dfrac{M_A}{M_{\text{air}}}$.
\[ M_A = d \times M_{\text{air}} = 0{,}97 \times 29 = \SI{28,1}{g.mol^{-1}} \approx \SI{28}{g.mol^{-1}} \]
\item La décoloration de l'eau de brome traduit une \textbf{réaction d'addition} de \ce{Br2} sur une liaison multiple : A est un hydrocarbure \textbf{insaturé}.
\[ \ce{C2H4 + Br2 -> C2H4Br2} \quad \text{(1,2-dibromoéthane, incolore)} \]
\item A est un alcène : $M = 14n = 28 \Rightarrow n = 2$. Formule brute : \textbf{\ce{C2H4}} ; formule semi-développée : \textbf{\ce{CH2=CH2}} ; nom : \textbf{éthène} (éthylène).
\item \textbf{Non.} Chaque carbone de la double liaison porte deux atomes d'hydrogène identiques : la condition d'existence de l'isomérie Z/E n'est pas remplie.
\item Hydrogénation catalytique :
\[ \ce{CH2=CH2 + H2 ->[Ni] CH3-CH3} \quad \text{produit : \textbf{éthane}} \]
\item $n_A = \dfrac{5{,}6}{28} = \SI{0,20}{mol}$. Combustion : $\ce{C2H4 + 3 O2 -> 2 CO2 + 2 H2O}$.
\begin{center}
\begin{tabularx}{\linewidth}{|l|X|X|X|X|}
\hline
\rowcolor{popblueL} Équation & \ce{C2H4} & \ce{+ 3 O2} & \ce{-> 2 CO2} & \ce{+ 2 H2O} \\
\hline
État initial (mol) & 0,20 & excès & 0 & 0 \\
\hline
État final (mol) & 0 & excès & 0,40 & 0,40 \\
\hline
\end{tabularx}
\end{center}
$n(\ce{CO2}) = 2 \times 0{,}20 = \SI{0,40}{mol}$ ; $V(\ce{CO2}) = 0{,}40 \times 24 = \boxed{\SI{9,6}{L}}$.
\end{enumerate}
\textbf{Barème indicatif (10 pts) :} (1) 1,5 pt ; (2) 1,5 pt ; (3) 2 pts ; (4) 1 pt ; (5) 1,5 pt ; (6) 2,5 pts. \textbf{Points de vigilance :} la densité est sans unité ; ne pas oublier le catalyseur Ni au-dessus de la flèche ; vérifier le facteur stœchiométrique 2 pour \ce{CO2}.
\end{corrige}

\begin{corrige}[Exercice 10 — L'acétylène du soudeur]
\begin{enumerate}
\item Équation-bilan de la préparation :
\[ \ce{CaC2(s) + 2 H2O(l) -> Ca(OH)2(aq) + C2H2(g)} \]
\item Dispositif : l'eau tombe goutte à goutte (ampoule de coulée) sur le carbure de calcium solide placé dans un ballon ; le gaz dégagé est recueilli par \textbf{déplacement d'eau} (cuve à eau, tube renversé), car \ce{C2H2} est très peu soluble dans l'eau.
\begin{popfigure}
\begin{tikzpicture}[scale=0.85, >=Stealth]
% support
\draw[thick] (0,0) -- (11,0);
% ampoule de coulée
\draw[thick] (2,4.6) ellipse (0.55 and 0.25);
\draw[thick] (1.45,4.6) -- (1.45,3.4) .. controls (1.45,3) and (2.55,3) .. (2.55,3.4) -- (2.55,4.6);
\draw[fill=popblueL] (1.5,4.55) -- (1.5,3.6) .. controls (1.6,3.2) and (2.4,3.2) .. (2.5,3.6) -- (2.5,4.55) -- cycle;
\draw[thick] (1.9,3.05) -- (1.9,2.3); \draw[thick] (2.1,3.05) -- (2.1,2.3);
\draw[thick] (1.8,2.9) -- (2.2,2.9);
\node[right] at (2.3,2.9) {\small robinet};
% ballon
\draw[thick] (4.2,1.2) .. controls (4.2,0.4) and (6.8,0.4) .. (6.8,1.2);
\draw[thick] (4.2,1.2) -- (4.2,2.4) -- (4.9,2.4) -- (4.9,2.9);
\draw[thick] (6.8,1.2) -- (6.8,2.4) -- (6.1,2.4) -- (6.1,2.9);
\draw[fill=poporangeL] (4.4,0.9) rectangle (6.6,1.5);
\node at (5.5,1.2) {\small \ce{CaC2} (solide)};
% tube d'arrivée d'eau
\draw[thick] (2,2.3) -- (2,1.9) -- (4.5,1.9) -- (4.5,1.55);
% tube de sortie
\draw[thick] (6.5,2.9) -- (6.5,3.3) -- (8.6,3.3) -- (8.6,0.6);
% cuve à eau
\draw[thick] (7.6,-0.4) -- (7.6,0.4) -- (10,0.4) -- (10,-0.4);
\draw[fill=popblueL!40] (7.65,-0.35) rectangle (9.95,0.2);
% tube renversé
\draw[thick] (8.3,0.2) -- (8.3,2.2) -- (9.1,2.2) -- (9.1,0.2);
\draw[fill=white] (8.35,0.6) rectangle (9.05,2.15);
\node at (8.7,1.5) {\small \ce{C2H2}};
% légendes
\node[align=center] at (2,5.3) {\small ampoule de coulée\\ \small (eau)};
\node[align=center] at (5.5,3.6) {\small ballon};
\node[align=center] at (10.6,1.1) {\small recueil par\\ \small déplacement d'eau};
\end{tikzpicture}
\legende{Préparation de l'acétylène : action de l'eau sur le carbure de calcium ; recueil du gaz par déplacement d'eau.}
\end{popfigure}
\item \textbf{Test d'insaturation :} faire buller le gaz dans l'eau de brome. \textbf{Observation :} décoloration de la solution orange. Équation-bilan :
\[ \ce{C2H2 + 2 Br2 -> CHBr2-CHBr2} \quad \text{(1,1,2,2-tétrabromoéthane, incolore)} \]
\item Quantité d'acétylène : $n(\ce{C2H2}) = \dfrac{V}{V_m} = \dfrac{48}{24} = \boxed{\SI{2,0}{mol}}$.
\item D'après l'équation, $n(\ce{CaC2}) = n(\ce{C2H2}) = \SI{2,0}{mol}$. $M(\ce{CaC2}) = 40 + 2 \times 12 = \SI{64}{g.mol^{-1}}$.
\[ m_{\text{pur}}(\ce{CaC2}) = 2{,}0 \times 64 = \boxed{\SI{128}{g}} \]
\item $n(\ce{Ca(OH)2}) = \SI{2,0}{mol}$ ; $M(\ce{Ca(OH)2}) = 40 + 2 \times (16 + 1) = \SI{74}{g.mol^{-1}}$.
\[ m(\ce{Ca(OH)2}) = 2{,}0 \times 74 = \boxed{\SI{148}{g}} \]
\item Masse de carbure commercial à 80 \% :
\[ m_{\text{commercial}} = \frac{128}{0{,}80} = \boxed{\SI{160}{g}} \]
\end{enumerate}
\textbf{Barème indicatif (12 pts) :} (1) 1,5 pt ; (2) 2,5 pts ; (3) 2 pts ; (4) 1 pt ; (5) 1,5 pt ; (6) 1,5 pt ; (7) 2 pts. \textbf{Points de vigilance :} équilibrer avec 2 \ce{H2O} ; dans le test, il faut \textbf{2 \ce{Br2}} par molécule d'alcyne (deux additions successives) ; diviser (et non multiplier) par la pureté pour trouver la masse commerciale ; l'acétylène ne doit jamais être recueilli dans un récipient en cuivre (acétylures explosifs).
\end{corrige}

\begin{corrige}[Exercice 11 — Du pétrole de la SONARA au plastique]
\begin{enumerate}
\item Formule développée de l'éthylène : \chemfig{H_2C=CH_2}. Chaque carbone est un \cle{carbone trigonal} (hybridation $sp^2$) : la molécule est \textbf{plane}, angles de liaison voisins de \textbf{\ang{120}}.
\item Hydratation catalytique :
\[ \ce{CH2=CH2 + H2O ->[H+] CH3-CH2-OH} \]
B est l'\textbf{éthanol}, de la famille des \textbf{alcools} (alcool primaire), utilisé comme antiseptique dans les hôpitaux.
\item Addition du dichlore :
\[ \ce{CH2=CH2 + Cl2 -> CH2Cl-CH2Cl} \]
C est le \textbf{1,2-dichloroéthane}.
\item Polymérisation de l'éthylène :
\[ n\,\ce{CH2=CH2} \ce{->[\text{catalyseur}]} \ce{(-CH2-CH2-)_n} \quad \text{(polyéthylène, PE)} \]
\item Motif \ce{C2H4} : $M_{\text{motif}} = 2 \times 12 + 4 \times 1 = \SI{28}{g.mol^{-1}}$.
\[ M_{\text{PE}} = 2000 \times 28 = \boxed{\SI{56000}{g.mol^{-1}}} \]
\[ m = n \times M = 0{,}050 \times 56000 = \boxed{\SI{2800}{g} = \SI{2,8}{kg}} \]
\item Motif assimilé à \ce{CH2} : $M_{\text{motif}} = \SI{14}{g.mol^{-1}}$.
\[ n_{\text{motifs}} = \frac{1000}{14} = \SI{71,4}{mol} \]
Combustion du motif : $\ce{CH2 + \frac{3}{2} O2 -> CO2 + H2O}$, donc $n(\ce{CO2}) = \SI{71,4}{mol}$.
\[ m(\ce{CO2}) = 71{,}4 \times 44 = \boxed{\SI{3143}{g} \approx \SI{3,1}{kg}} \]
\textbf{Commentaire :} brûler 1 kg de sachets plastiques rejette plus de 3 kg de \ce{CO2}, gaz à effet de serre qui aggrave le réchauffement climatique ; la combustion sauvage, souvent incomplète, dégage en outre du monoxyde de carbone \ce{CO} (toxique) et des suies. Il faut privilégier la réduction de l'usage des sachets, le tri, le recyclage et la valorisation contrôlée.
\end{enumerate}
\textbf{Barème indicatif (12 pts) :} (1) 2 pts ; (2) 2 pts ; (3) 1,5 pt ; (4) 1,5 pt ; (5) 2 pts ; (6) 3 pts. \textbf{Points de vigilance :} distinguer la masse molaire du \emph{motif} (\SI{28}{g.mol^{-1}} pour \ce{C2H4}) de celle du \emph{polymère} ; dans la question 6, le raisonnement se fait sur le motif \ce{CH2} (\SI{14}{g.mol^{-1}}) et non sur le monomère entier ; toujours conclure un calcul par une valeur encadrée avec son unité.
\end{corrige}

\newpage

\coursec{Fiche bilan}

\coursec{Les alcènes et les alcynes}

\courssub{Carte mentale de la leçon}

\begin{popfigure}
\begin{tikzpicture}[
  central/.style={ellipse, draw=popdark, fill=popblue!25, thick, align=center, font=\bfseries\small, inner sep=3pt},
  branche/.style={rounded corners=3pt, draw=#1, fill=#1!18, thick, align=center, font=\scriptsize, inner sep=3pt},
  lien/.style={thick, draw=#1},
  level distance=2.6cm]
\node[central] (C) at (0,0) {Alcènes\\et alcynes};
\node[branche=popblue] (S) at (-4.6,2.2) {\textbf{Structure}\\double liaison : trigonal\\triple liaison : digonal};
\node[branche=poporange] (F) at (0,3.0) {\textbf{Formules générales}\\$C_nH_{2n}$ (alcène)\\$C_nH_{2n-2}$ (alcyne)};
\node[branche=popgreen] (N) at (4.6,2.2) {\textbf{Nomenclature}\\suffixe -ène / -yne\\indice le plus petit};
\node[branche=poppurple] (I) at (5.0,-1.2) {\textbf{Isomérie}\\chaîne, position\\Z / E};
\node[branche=popgold] (A) at (2.6,-3.0) {\textbf{Additions}\\$H_2$, $X_2$, $HX$, $H_2O$\\règle de Markovnikov};
\node[branche=popink] (P) at (-2.6,-3.0) {\textbf{Polymérisation}\\$n$ monomères $\to$ polymère\\PE, PVC, PP};
\node[branche=popmuted] (Ac) at (-5.0,-1.2) {\textbf{Acétylène}\\$\ce{CaC2 + 2 H2O}$\\soudure oxyacétylénique};
\draw[lien=popblue] (C) -- (S);
\draw[lien=poporange] (C) -- (F);
\draw[lien=popgreen] (C) -- (N);
\draw[lien=poppurple] (C) -- (I);
\draw[lien=popgold] (C) -- (A);
\draw[lien=popink] (C) -- (P);
\draw[lien=popmuted] (C) -- (Ac);
\end{tikzpicture}
\legende{Carte mentale de synthèse : les sept idées-forces de la leçon.}
\end{popfigure}

\courssub{Structure et géométrie des alcènes et des alcynes}

\begin{definition}[Alcène et alcyne]
Un \cle{alcène} est un hydrocarbure insaturé à chaîne ouverte contenant une \cle{double liaison} carbone-carbone \ce{C=C}. Sa formule générale est \textbf{$C_nH_{2n}$} ($n \geq 2$). Le plus simple est l'\cle{éthylène} (éthène) \ce{CH2=CH2}.
Un \cle{alcyne} est un hydrocarbure insaturé à chaîne ouverte contenant une \cle{triple liaison} carbone-carbone \ce{C#C}. Sa formule générale est \textbf{$C_nH_{2n-2}$} ($n \geq 2$). Le plus simple est l'\cle{acétylène} (éthyne) \ce{HC#CH}.
\end{definition}

\begin{propriete}[Hybridation et géométrie]
\begin{itemize}
  \item \textbf{Carbone trigonal (alcène) :} Le carbone de la double liaison est hybridé \cle{sp$^2$}. Les trois orbitaux hybrides forment trois liaisons $\sigma$ dans un plan (angle \SI{120}{\degree}). L'orbitale $p_z$ restante forme la liaison $\pi$ par recouvrement latéral, au-dessus et au-dessous du plan.
  \item \textbf{Carbone digonal (alcyne) :} Le carbone de la triple liaison est hybridé \cle{sp}. Les deux orbitaux hybrides forment deux liaisons $\sigma$ alignées (angle \SI{180}{\degree}). Les deux orbitales $p$ restantes forment deux liaisons $\pi$ perpendiculaires, créant un nuage électronique cylindrique autour de l'axe internucléaire.
\end{itemize}
\end{propriete}

\begin{exemplebox}[Éthylène et acétylène]
\textbf{Éthylène (\ce{C2H4}) :} Molécule plane. Liaison \ce{C=C} : 1 $\sigma$ + 1 $\pi$. Longueur \ce{C=C} $\approx$ \SI{134}{pm}. Liaisons \ce{C-H} dans le plan.
\textbf{Acétylène (\ce{C2H2}) :} Molécule linéaire. Liaison \ce{C#C} : 1 $\sigma$ + 2 $\pi$. Longueur \ce{C#C} $\approx$ \SI{120}{pm} (plus courte et plus forte que \ce{C=C}).
\end{exemplebox}

\begin{popfigure}
\begin{tikzpicture}[scale=1.2, every node/.style={font=\small}]
% Ethylene geometry
\begin{scope}[xshift=-4cm]
\draw[thick, popblue] (0,0) -- (1.5,0) node[midway, below] {$\sigma$};
\draw[thick, popblue, decorate, decoration={snake, amplitude=1.5pt, segment length=4pt}] (0,0.15) -- (1.5,0.15) node[midway, above] {$\pi$};
\draw[thick, popblue, decorate, decoration={snake, amplitude=1.5pt, segment length=4pt}] (0,-0.15) -- (1.5,-0.15);
\foreach \angle/\label in {0/H, 120/H, 240/H} {
  \draw[thick, popdark] (0,0) -- (\angle:0.8) node[shift=(\angle:0.2)] {\label};
  \draw[thick, popdark] (1.5,0) -- ++(\angle:0.8) node[shift=(\angle:0.2)] {\label};
}
\node[popblue, align=center] at (0.75, -1.5) {Éthylène : planaire\\angles \SI{120}{\degree}\\hybridation \textbf{sp$^2$}};
\end{scope}

% Acetylene geometry
\begin{scope}[xshift=4cm]
\draw[thick, poporange] (-1.5,0) -- (1.5,0) node[midway, below] {$\sigma$};
\draw[thick, poporange, decorate, decoration={snake, amplitude=1pt, segment length=3pt}] (-1.5,0.2) -- (1.5,0.2) node[midway, above] {$\pi_y$};
\draw[thick, poporange, decorate, decoration={snake, amplitude=1pt, segment length=3pt}] (-1.5,-0.2) -- (1.5,-0.2);
\draw[thick, poporange, decorate, decoration={snake, amplitude=1pt, segment length=3pt}] (-1.5,0) -- (1.5,0) node[midway, above=4pt] {$\pi_z$};
\draw[thick, popdark] (-1.5,0) -- (-2.5,0) node[left] {H};
\draw[thick, popdark] (1.5,0) -- (2.5,0) node[right] {H};
\node[poporange, align=center] at (0, -1.5) {Acétylène : linéaire\\angle \SI{180}{\degree}\\hybridation \textbf{sp}};
\end{scope}
\end{tikzpicture}
\legende{Géométrie des liaisons multiples : liaisons $\sigma$ (traits pleins) et $\pi$ (ondulées).}
\end{popfigure}

\courssub{Formules générales et écriture des formules}

\begin{propriete}[Formules générales]
\begin{center}
\begin{tabularx}{\linewidth}{|l|X|X|}
\hline
\rowcolor{popblueL} \textbf{Famille} & \textbf{Formule générale} & \textbf{Indice de degré d'insaturation (IDI)} \\
\hline
Alcanes (saturés) & $C_nH_{2n+2}$ & 0 \\
\hline
Alcènes (1 double liaison) & $C_nH_{2n}$ & 1 \\
\hline
Alcynes (1 triple liaison) & $C_nH_{2n-2}$ & 2 \\
\hline
Cycloalcanes (1 cycle) & $C_nH_{2n}$ & 1 \\
\hline
\end{tabularx}
\end{center}
\emph{Rappel :} L'IDI = $\frac{2n_C + 2 - n_H}{2}$. Un alcène a un IDI de 1, un alcyne un IDI de 2.
\end{propriete}

\begin{methode}[Écrire la formule semi-développée d'un alcène/alcyne]
\begin{enumerate}
  \item Écrire la chaîne carbonée principale ($n$ atomes de C).
  \item Placer la liaison multiple (\ce{C=C} ou \ce{C#C}) aux positions indiquées par les locants.
  \item Saturer les valences restantes par des atomes d'hydrogène (4 liaisons par carbone).
  \item Ajouter les groupes substituants (alkyl, halogènes, etc.) aux positions indiquées.
\end{enumerate}
\end{methode}

\begin{exemplebox}[Exemples d'écriture]
\begin{itemize}
  \item \textbf{But-2-ène} ($n=4$, double liaison en 2) : \chemfig{CH_3-CH=CH-CH_3}
  \item \textbf{3-Méthylpent-1-ène} : chaîne 5C, double liaison en 1, méthyle en C3.
  \chemfig{CH_2=CH-CH(-[2]CH_3)-CH_2-CH_3}
  \item \textbf{But-1-yne} ($n=4$, triple liaison en 1) : \chemfig{HC#C-CH_2-CH_3}
  \item \textbf{Pent-2-yne} : \chemfig{CH_3-C#C-CH_2-CH_3}
\end{itemize}
\end{exemplebox}

\courssub{Nomenclature des alcènes et des alcynes}

\begin{definition}[Règles IUPAC (français) pour les alcènes et alcynes]
\begin{enumerate}
  \item \textbf{Chaîne principale :} La plus longue chaîne carbonée \textbf{contenant la liaison multiple} (même si une chaîne plus longue sans liaison multiple existe).
  \item \textbf{Numérotation :} On numérote la chaîne pour donner le \textbf{plus petit indice possible à la liaison multiple} (règle du « premier point de différence » appliquée à la liaison multiple en priorité).
  \item \textbf{Suffixe :} \textbf{-ène} pour alcène, \textbf{-yne} pour alcyne (remplace le -ane de l'alcane correspondant).
  \item \textbf{Substituants :} Cité(s) par ordre alphabétique, préfixés de leur locant.
  \item \textbf{Plusieurs liaisons multiples :} \textit{diène, triène} ; \textit{diyn, triyn} ; \textit{èn-yne} (double liaison prioritaire pour la numérotation si choix).
\end{enumerate}
\end{definition}

\begin{attention}[Pièges fréquents en nomenclature]
\begin{itemize}
  \item Ne pas choisir la chaîne la plus longue \emph{sans} la liaison multiple. Ex: \chemfig{CH_2=CH-CH_2-CH_2-CH_3} est le \textbf{pent-1-ène} (5C), pas un butène substitué.
  \item Numéroter \emph{toujours} depuis l'extrémité la plus proche de la liaison multiple. Ex: \chemfig{CH_3-CH=CH-CH_2-CH_3} est le \textbf{pent-2-ène} (indices 2), pas pent-3-ène.
  \item Pour les alcynes terminaux (\ce{#C-H}), l'hydrogène acide n'est pas indiqué dans le nom (ex: \textbf{but-1-yne} et non but-1-yn-1-ol).
\end{itemize}
\end{attention}

\begin{exemplebox}[Exercices de nomenclature]
\begin{enumerate}
  \item \chemfig{CH_3-CH_2-CH=CH-CH_3} $\to$ \textbf{Pent-2-ène}
  \item \chemfig{CH_2=CH-CH(CH_3)-CH_2-CH_3} $\to$ \textbf{3-Méthylpent-1-ène}
  \item \chemfig{CH_3-C#C-CH_3} $\to$ \textbf{But-2-yne}
  \item \chemfig{HC#C-CH_2-CH(CH_3)_2} $\to$ \textbf{4-Méthylpent-1-yne}
  \item \chemfig{CH_2=CH-CH=CH_2} $\to$ \textbf{Buta-1,3-diène}
\end{enumerate}
\end{exemplebox}

\courssub{Isomérie Z et E (stéréoisomérie)}

\begin{definition}[Isomérie de configuration (Z/E)]
L'absence de rotation libre autour de la double liaison \ce{C=C} (liaison $\pi$ rigide) engendre une \cle{stéréoisomérie} lorsque \textbf{chacun des deux carbones de la double liaison porte deux substituants différents}.
On utilise les règles de \textbf{Cahn-Ingold-Prelog (CIP)} pour assigner la priorité sur chaque carbone : l'atome de plus grand numéro atomique directement lié au carbone a la priorité.
\begin{itemize}
  \item \textbf{Configuration Z (zusammen = ensemble) :} Les deux groupes de priorité forte sont du \textbf{même côté} du plan de la double liaison.
  \item \textbf{Configuration E (entgegen = opposé) :} Les deux groupes de priorité forte sont de \textbf{part et d'autre} du plan de la double liaison.
\end{itemize}
\end{definition}

\begin{methode}[Déterminer Z ou E]
\begin{enumerate}
  \item Identifier les deux carbones de la double liaison ($C_a$ et $C_b$).
  \item Sur $C_a$, comparer les numéros atomiques des deux atomes liés directement : le plus grand a la priorité 1.
  \item Faire de même sur $C_b$.
  \item Observer l'orientation relative des deux groupes prioritaires (1 sur $C_a$ et 1 sur $C_b$).
  \item Même côté $\to$ \textbf{Z} ; Côtés opposés $\to$ \textbf{E}.
\end{enumerate}
\end{methode}

\begin{exemplebox}[But-2-ène : les deux stéréoisomères]
\begin{center}
\begin{tabular}{cc}
\textbf{(Z)-But-2-ène} & \textbf{(E)-But-2-ène} \\
\chemfig{CH_3-[:30](=[:90]H)-[:-30]CH_3} & \chemfig{CH_3-[:30](=[:-90]H)-[:-30]CH_3} \\
Priorités : C > H sur chaque C. & Priorités : C > H sur chaque C. \\
Groupes CH$_3$ du même côté $\to$ \textbf{Z}. & Groupes CH$_3$ opposés $\to$ \textbf{E}.
\end{tabular}
\end{center}
\emph{Note :} Pour les alcènes simples, Z $\approx$ \textit{cis} et E $\approx$ \textit{trans}, mais la notation Z/E est universelle et rigoureuse.
\end{exemplebox}

\begin{attention}[Cas particuliers]
\begin{itemize}
  \item Si un carbone de la double liaison porte deux groupes identiques (ex: \ce{CH2=} ou \ce{C(CH3)2=}), \textbf{pas d'isomérie Z/E possible}.
  \item Pour les cycles (cycloalcènes), la contrainte du cycle impose souvent la configuration Z pour les petits cycles ($n < 8$).
\end{itemize}
\end{attention}

\courssub{Propriétés chimiques : combustion et réactions d'addition}

\begin{propriete}[Combustion complète]
Les alcènes et alcynes brûlent en présence de dioxygène pour donner du dioxyde de carbone et de l'eau. La flamme est plus lumineuse (suie) que pour les alcanes car le rapport C/H est plus élevé.
\begin{align*}
\text{Alcène } (C_nH_{2n}) &: \quad \ce{C_nH_{2n} + \frac{3n}{2} O2 -> n CO2 + n H2O} \\
\text{Alcyne } (C_nH_{2n-2}) &: \quad \ce{C_nH_{2n-2} + \frac{3n-1}{2} O2 -> n CO2 + (n-1) H2O}
\end{align*}
\textbf{Exemple (éthylène) :} $\ce{C2H4 + 3 O2 -> 2 CO2 + 2 H2O}$ \\
\textbf{Exemple (acétylène) :} $\ce{C2H2 + 5/2 O2 -> 2 CO2 + H2O}$ (flamme très chaude, \SI{3000}{\celsius} avec \ce{O2} pur).
\end{propriete}

\begin{propriete}[Réactions d'addition sur la double liaison \ce{C=C}]
La liaison $\pi$ se rompt, deux nouvelles liaisons $\sigma$ se forment. Le carbone passe de l'hybridation sp$^2$ (trigonal) à sp$^3$ (tétraédrique).
\begin{center}
\begin{tabularx}{\linewidth}{|l|l|X|}
\hline
\rowcolor{popblueL} \textbf{Réaction} & \textbf{Réactifs / Conditions} & \textbf{Équation type (éthylène)} \\
\hline
Hydrogénation & \ce{H2}, catalyseur Ni, Pt ou Pd, chauffage & $\ce{CH2=CH2 + H2 ->[Ni] CH3-CH3}$ \\
\hline
Halogénation (test) & \ce{Cl2} ou \ce{Br2} (dans \ce{CCl4} ou eau de brome), temp. ambiante & $\ce{CH2=CH2 + Br2 -> CH2Br-CH2Br}$ \\
\hline
Hydrohalogénation & \ce{HCl}, \ce{HBr}, \ce{HI} (gaz ou sol. aqueuse), froid & $\ce{CH2=CH2 + HBr -> CH3-CH2Br}$ \\
\hline
Hydratation & \ce{H2O}, catalyseur acide (\ce{H3PO4} solide, \SI{300}{\celsius} indus. ; \ce{H2SO4} conc. labo) & $\ce{CH2=CH2 + H2O ->[H^+] CH3-CH2OH}$ \\
\hline
\end{tabularx}
\end{center}
\end{propriete}

\begin{definition}[Règle de Markovnikov]
Lors de l'addition d'un réactif \textbf{dissymétrique} (\ce{HX}, \ce{H2O}, \ce{H2SO4}) sur un \textbf{alcène dissymétrique} (les deux carbones de la double liaison n'ont pas le même nombre d'hydrogènes), l'atome d'\textbf{hydrogène se fixe sur le carbone déjà le plus hydrogéné} (celui qui porte le plus d'atomes d'H), et l'autre atome/groupe (\ce{X}, \ce{OH}) se fixe sur l'autre carbone.
\emph{Origine (niveau supérieur) :} Stabilité de l'intermédiaire carbocationique tertiaire $>$ secondaire $>$ primaire.
\end{definition}

\begin{methode}[Prévoir le produit majoritaire (Markovnikov)]
\begin{enumerate}
  \item Identifier les deux carbones de la double liaison : $C_\alpha$ et $C_\beta$.
  \item Compter les H sur chaque carbone : celui qui en a \textbf{le plus} reçoit le \textbf{H} du réactif \ce{H-X} (ou \ce{H-OH}).
  \item L'autre carbone reçoit \textbf{X} (ou \ce{OH}).
\end{enumerate}
\end{methode}

\begin{exemplebox}[Addition de HBr sur le propène]
\chemfig{CH_3-CH=CH_2 + H-Br -> CH_3-CH(Br)-CH_3}
\begin{itemize}
  \item Carbone de gauche (C2) : 1 H. Carbone de droite (C3) : 2 H.
  \item Le H de HBr va sur C3 (le plus hydrogéné). Le Br va sur C2.
  \item Produit : \textbf{2-bromopropane} (majoritaire). Le 1-bromopropane est minoritaire (anti-Markovnikov, nécessite des peroxydes).
\end{itemize}
\end{exemplebox}

\begin{experience}[Test de l'eau de brome (mise en évidence de l'insaturation)]
\begin{itemize}
  \item \textbf{Dispositif :} Tube à essai, eau de brome (\ce{Br2}/\ce{H2O}, couleur orange), alcène (ex: cyclohexène) ou alcyne.
  \item \textbf{Protocole :} Verser \SI{2}{mL} d'eau de brome. Ajouter quelques gouttes d'alcène. Agiter.
  \item \textbf{Observation :} La couleur orange \textbf{disparaît} (décoloration) instantanément. Formation de dibromure incolore.
  \item \textbf{Interprétation :} Addition de \ce{Br2} sur la liaison multiple. \textbf{Test positif = insaturation (alcène ou alcyne).}
  \item \textbf{Contraste :} Avec un alcane (cyclohexane), l'eau de brome conserve sa couleur (pas de réaction à l'obscurité).
\end{itemize}
\end{experience}

\begin{propriete}[Réactions d'addition sur la triple liaison \ce{C\#C}]
L'alcyne peut ajouter \textbf{2 équivalents} de réactif.
\begin{itemize}
  \item \textbf{1er équivalent :} Formation d'un alcène (stéréochimie \textit{anti} pour \ce{X2}, \textit{syn} pour \ce{H2}/Lindlar).
  \item \textbf{2ème équivalent :} Formation d'un alcane (ou dérivé saturé).
\end{itemize}
\textbf{Exemple (acétylène + 2 \ce{HBr}) :}
\[
\ce{HC#CH + HBr -> CH2=CHBr} \quad\text{(bromure de vinyle)}\quad \ce{->[HBr] CH3-CHBr2} \quad\text{(1,1-dibromoéthane)}
\]
\textbf{Hydratation de l'acétylène (procédé Koutchakov-Répiquet) :}
\[
\ce{HC#CH + H2O ->[Hg^{2+}, H^+] CH2=CH-OH} \text{ (énol) } \ce{->[\text{tautomérisation}] CH3-CHO} \text{ (éthanal)}
\]
\end{propriete}

\courssub{Polymérisation et préparation de l'acétylène}

\begin{definition}[Polymérisation par addition]
Réaction dans laquelle $n$ molécules de \cle{monomère} (alcène ou dérivé) s'additionnent les unes aux autres par ouverture de la double liaison pour former une \cle{macromolécule} (polymère) sans perte d'atomes.
\[ \ce{n CH2=CH-R ->[-CH2-CH(R)-]_n} \]
L'unité répétitive est \ce{[-CH2-CH(R)-]}.
\end{definition}

\begin{exemplebox}[Polymères industriels majeurs]
\begin{center}
\begin{tabularx}{\linewidth}{|l|l|l|X|}
\hline
\rowcolor{popblueL} \textbf{Monomère} & \textbf{Nom usuel} & \textbf{Polymère} & \textbf{Usages courants} \\
\hline
\ce{CH2=CH2} & Éthylène & Polyéthylène (PE) & Sachets, bouteilles, films, tuyaux (CIMENCAM) \\
\hline
\ce{CH2=CH-CH3} & Propylène & Polypropylène (PP) & Emballages rigides, textiles, pare-chocs \\
\hline
\ce{CH2=CH-Cl} & Chloroéthène (chlorure de vinyle) & Polychlorure de vinyle (PVC) & Tuyaux évacuation, fenêtres, câbles électriques \\
\hline
\ce{CH2=CH-CN} & Acrylonitrile & Polyacrylonitrile (PAN) & Fibres textiles (acrylique), précurseur fibres carbone \\
\hline
\end{tabularx}
\end{center}
\end{exemplebox}

\begin{savaistu}[Polymères et Cameroun]
Au Cameroun, les industries plastiques (usines de Douala, Yaoundé) transforment le PE et le PVC en tuyaux pour l'adduction d'eau (CAMWATER), en sachets pour l'eau « \textit{pure water} » et en emballages. La gestion des déchets plastiques (non biodégradables) est un défi environnemental majeur : incinération contrôlée, recyclage mécanique (broyage-refusion) ou chimique (pyrolyse). La SONARA (Société Nationale de Raffinage) produit les coupes légères (naphta) qui, par vapocraquage, donnent l'éthylène et le propylène, matières premières de ces polymères.
\end{savaistu}

\begin{experience}[Préparation de l'acétylène au laboratoire]
\begin{itemize}
  \item \textbf{Principe :} Hydrolyse du carbure de calcium \ce{CaC2} (pierre grise, aspect métallique).
  \item \textbf{Équation :} \[ \ce{CaC2(s) + 2 H2O(l) -> C2H2(g) + Ca(OH)2(aq)} \]
  \item \textbf{Dispositif :} Générateur à acétylène (type Kipp ou ballon à col large avec entonnoir à robinet / tubulure de dégagement). \textbf{Pas de chauffage.}
  \item \textbf{Purification :} Le gaz passe dans un laveur à \ce{CuSO4}/\ce{NH4OH} (élimine \ce{PH3}, \ce{H2S}, \ce{AsH3} toxiques) puis dans un tube à \ce{CaCl2} anhydre ou \ce{H2SO4} conc. (séchage).
  \item \textbf{Recueil :} Par déplacement d'air (gaz plus léger que l'air, $d \approx 0,9$) ou sur eau (peu soluble). \textbf{Interdit :} Recueil sur mercure.
  \item \textbf{Sécurité :} \ce{C2H2} forme des mélanges explosifs avec l'air (\SIrange{2,5}{82}{\%} vol.). \textbf{Ne jamais comprimer l'acétylène pur} (risque décomposition explosive). En bouteille industrielle : dissous dans l'acétone sur support poreux.
  \item \textbf{Test :} Flamme très vive, suieuse à l'air ; décolore l'eau de brome.
\end{itemize}
\end{experience}

\begin{popfigure}
\begin{tikzpicture}[scale=0.9, every node/.style={font=\small}]
% Generator
\draw[thick, popdark] (-2,0) rectangle (2,4);
\draw[thick, popdark] (-2,4) -- (-1,5) -- (1,5) -- (2,4);
\draw[thick, popdark] (-0.3,5) -- (-0.3,6) node[above] {Entonnoir à robinet};
\draw[fill=poporange!30] (-1.5,1) rectangle (-0.5,3) node[midway, align=center] {CaC$_2$};
\draw[fill=popblue!20] (0.5,0.5) rectangle (1.5,2.5) node[midway, align=center] {Eau};
\draw[->, thick, popblue] (0,0.5) -- (0,-0.5) node[below] {Écoulement eau};
% Wash bottle
\draw[thick, popdark] (4,0) circle (1.5 and 2);
\draw[thick, popdark] (4,2) ellipse (1.5 and 0.3);
\draw[thick, popdark] (2.5,2) -- (3.5,2) node[midway, above] {Gaz brut};
\draw[->, thick, popdark] (4.5,2) -- (5.5,2) node[midway, above] {Gaz lavé};
\draw[fill=popblue!20] (3,0) rectangle (5,1.8) node[midway, align=center] {Sol. CuSO$_4$/\ce{NH3}};
% Drying tube
\draw[thick, popdark] (7,0) -- (7,3) -- (9,3) -- (9,0);
\draw[thick, popdark] (7,0) -- (9,0);
\draw[fill=popgold!30] (7.2,0.2) rectangle (8.8,2.8) node[midway, align=center] {CaCl$_2$ anhydre};
\draw[->, thick, popdark] (9,1.5) -- (10,1.5) node[right] {Acétylène sec};
% Collection
\draw[thick, popdark] (11,0) -- (11,4) -- (13,4) -- (13,0);
\draw[thick, popdark] (11,0) -- (13,0);
\draw[fill=popmuted!30] (11.2,0.2) rectangle (12.8,3.5) node[midway, align=center] {Air déplacé};
\draw[->, thick, popdark] (10,1.5) -- (11,1.5);
\end{tikzpicture}
\legende{Schéma du dispositif de préparation de l'acétylène : générateur, lavage (purification), séchage, recueil par déplacement d'air.}
\end{popfigure}

\begin{savaistu}[L'acétylène et la soudure oxyacétylénique]
La combustion de l'acétylène dans l'oxygène pur produit une flamme atteignant \SI{3100}{\celsius} (la plus chaude des combustibles gazeux usuels). C'est la base de la \textbf{soudure oxyacétylénique} (chalumeau), très utilisée dans les garages et ateliers de métallurgie au Cameroun pour le brasage, le coupage et le soudage des aciers. Les bouteilles blanches (acétylène dissous dans l'acétone) et noires (oxygène) sont des standards de sécurité.
\end{savaistu}

\courssub{L'essentiel à retenir}

\begin{aretenir}[Synthèse des savoirs et savoir-faire]
\textbf{1. Structures et formules}
\begin{itemize}
  \item \cle{Alcène} : \ce{C=C}, carbone \textbf{trigonal} (sp$^2$), plan, \SI{120}{\degree}, formule générale \textbf{$C_nH_{2n}$}.
  \item \cle{Alcyne} : \ce{C#C}, carbone \textbf{digonal} (sp), linéaire, \SI{180}{\degree}, formule générale \textbf{$C_nH_{2n-2}$}.
  \item Représentations : formules développées, semi-développées (ex: \chemfig{CH_3-CH=CH-CH_3}), formules de Cram pour la stéréochimie.
\end{itemize}

\textbf{2. Nomenclature (IUPAC)}
\begin{itemize}
  \item Chaîne principale = plus longue chaîne \textbf{contenant} la liaison multiple.
  \item Numérotation = indice le plus bas pour la liaison multiple (suffixe \textbf{-ène} ou \textbf{-yne}).
  \item Substituants par ordre alphabétique.
\end{itemize}

\textbf{3. Isomérie}
\begin{itemize}
  \item Isomères de chaîne, de position (double/triple liaison), de fonction (alcène/alcyne/cycloalcane).
  \item \textbf{Stéréoisomérie Z/E :} Priorité CIP (numéro atomique) sur chaque C de \ce{C=C}. Même côté = \textbf{Z}, côtés opposés = \textbf{E}.
\end{itemize}

\textbf{4. Réactions chimiques types (à connaître par cœur)}
\begin{itemize}
  \item \textbf{Combustion complète :}
  \begin{align*}
    \ce{C_nH_{2n} + \frac{3n}{2} O2 &-> n CO2 + n H2O} \\
    \ce{C_nH_{2n-2} + \frac{3n-1}{2} O2 &-> n CO2 + (n-1) H2O}
  \end{align*}
  \item \textbf{Additions sur \ce{C=C} :}
  \begin{align*}
    \text{Hydrogénation :} &\quad \ce{CH2=CH2 + H2 ->[Ni] CH3-CH3} \\
    \text{Halogénation (test) :} &\quad \ce{CH2=CH2 + Br2 -> CH2Br-CH2Br} \quad \text{(décoloration eau de brome)} \\
    \text{Hydrohalogénation :} &\quad \ce{CH3-CH=CH2 + HBr -> CH3-CHBr-CH3} \quad \text{(Markovnikov)} \\
    \text{Hydratation :} &\quad \ce{CH2=CH2 + H2O ->[H^+] CH3-CH2OH}
  \end{align*}
  \item \textbf{Additions sur \ce{C#C} :} 2 équivalents possibles. Hydratation $\to$ cétone (via énol).
  \item \textbf{Polymérisation :} $\ce{n CH2=CH-R -> [-CH2-CH(R)-]_n}$ (PE, PP, PVC...).
  \item \textbf{Préparation acétylène :} $\ce{CaC2 + 2 H2O -> C2H2 + Ca(OH)2}$.
\end{itemize}

\textbf{5. Méthodes clés}
\begin{itemize}
  \item \textbf{Nommer} : Repérer liaison multiple $\to$ Chaîne principale $\to$ Numéroter (indice min) $\to$ Écrire (substituants + locant + suffixe).
  \item \textbf{Config Z/E} : Priorités CIP sur chaque C $\to$ Observer relative position.
  \item \textbf{Produit addition HX} : H sur C le plus hydrogéné (Markovnikov).
  \item \textbf{Test insaturation} : Eau de brome $\to$ décoloration rapide.
\end{itemize}
\end{aretenir}

\begin{aretenir}[Checklist avant le Probatoire]
\begin{itemize}
  \item $\square$ Je donne les formules générales : alcènes $C_nH_{2n}$, alcynes $C_nH_{2n-2}$.
  \item $\square$ Je décris la géométrie : trigonal plan \SI{120}{\degree} (alcène), digonal linéaire \SI{180}{\degree} (alcyne).
  \item $\square$ J'écris les formules (développées, semi-développées) à partir du nom et inversement.
  \item $\square$ Je distingue isomères de chaîne, de position, et stéréoisomères Z/E (règles CIP).
  \item $\square$ J'équilibre les équations de combustion complète.
  \item $\square$ J'écris les équations d'addition : \ce{H2}, \ce{Cl2}, \ce{Br2}, \ce{HCl}, \ce{HBr}, \ce{H2O} sur \ce{C=C}.
  \item $\square$ J'applique la règle de Markovnikov pour prévoir le produit majoritaire sur alcène dissymétrique.
  \item $\square$ Je sais que la décoloration de l'eau de brome est le test caractéristique des insaturations.
  \item $\square$ J'écris l'équation de polymérisation de l'éthylène, du propylène, du chloroéthène ; je cite des usages (PE, PP, PVC).
  \item $\square$ Je décris et schématise le dispositif de préparation de l'acétylène (\ce{CaC2} + \ce{H2O}), purification, séchage, recueil.
  \item $\square$ Je cite l'utilisation majeure de l'acétylène : soudure oxyacétylénique (flamme \SI{3100}{\celsius}).
\end{itemize}
\end{aretenir}

\courssub{Exercices d'entraînement}

\begin{exercice}[Nomenclature et isomérie]
\begin{enumerate}
  \item Donner la formule semi-développée et le nom IUPAC des isomères d'alcènes de formule brute $C_5H_{10}$.
  \item Parmi ces isomères, lesquels présentent une isomérie Z/E ? Représenter les deux formes pour chacun.
  \item Nommer : \chemfig{CH_3-CH_2-C(CH_3)=CH-CH_3} et \chemfig{HC#C-CH_2-CH(CH_3)_2}.
\end{enumerate}
\end{exercice}

\begin{exercice}[Réactions d'addition et règle de Markovnikov]
\begin{enumerate}
  \item Écrire l'équation de l'addition de \ce{HBr} sur : (a) le but-1-ène, (b) le 2-méthylbut-2-ène, (c) le pent-2-ène. Préciser le(s) produit(s) majeur(s).
  \item L'addition d'eau sur le propène en milieu acide donne un alcool. Lequel ? Écrire l'équation.
  \item L'acétylène réagit avec 2 molécules de \ce{Cl2}. Écrire les deux étapes et le produit final.
\end{enumerate}
\end{exercice}

\begin{exercice}[Dosage et réactivité — Type Probatoire]
On réalise le dosage d'une solution commerciale d'acétylène dissous dans l'acétone (bouteille de \SI{5}{L}) par réaction avec l'eau de brome.
On prélève \SI{10,0}{mL} de gaz acétylène (mesuré aux CNTP) que l'on fait bouillir dans \SI{50,0}{mL} d'eau de brome à \SI{0,100}{mol.L^{-1}}.
Après réaction, l'excès de brome est dosé par \SI{20,0}{mL} de solution de thiosulfate de sodium \SI{0,100}{mol.L^{-1}}.
\begin{enumerate}
  \item Rappeler l'équation de la réaction \ce{C2H2 + 2 Br2 -> C2H2Br4}.
  \item Calculer la quantité de matière de brome introduite, puis celle réagie avec l'acétylène.
  \item En déduire la quantité de matière d'acétylène dans les \SI{10,0}{mL} prélevés.
  \item Calculer le volume molaire de l'acétylène aux CNTP trouvé expérimentalement. Comparer à la valeur théorique (\SI{22,4}{L.mol^{-1}}).
\end{enumerate}
\end{exercice}

\begin{corrige}[Exercice 1]
\begin{enumerate}
  \item Isomères $C_5H_{10}$ (alcènes) :
  \begin{itemize}
    \item Pent-1-ène : \chemfig{CH_2=CH-CH_2-CH_2-CH_3} (pas de Z/E, C1 porte 2H)
    \item Pent-2-ène : \chemfig{CH_3-CH=CH-CH_2-CH_3} $\to$ \textbf{(Z)-pent-2-ène} et \textbf{(E)-pent-2-ène}
    \item 2-Méthylbut-1-ène : \chemfig{CH_2=C(CH_3)-CH_2-CH_3} (pas de Z/E, C1 porte 2H)
    \item 3-Méthylbut-1-ène : \chemfig{CH_2=CH-CH(CH_3)_2} (pas de Z/E)
    \item 2-Méthylbut-2-ène : \chemfig{CH_3-C(CH_3)=CH-CH_3} (pas de Z/E, C2 porte 2 groupes CH3)
  \end{itemize}
  \item Seuls le \textbf{pent-2-ène} présente une isomérie Z/E.
  \item Noms : \textbf{3-Méthylpent-2-ène} (attention : chaîne 5C, double liaison en 2, méthyle en 3) et \textbf{4-Méthylpent-1-yne}.
\end{enumerate}
\end{corrige}

\begin{corrige}[Exercice 2]
\begin{enumerate}
  \item (a) But-1-ène \ce{CH2=CH-CH2-CH3} + HBr $\to$ \textbf{2-bromobutane} (Markovnikov : H sur C1, Br sur C2).
  (b) 2-Méthylbut-2-ène \ce{CH3-C(CH3)=CH-CH3} + HBr $\to$ \textbf{2-bromo-2-méthylbutane} (H sur C3, Br sur C2 tertiaire).
  (c) Pent-2-ène \ce{CH3-CH=CH-CH2-CH3} (symétrique) + HBr $\to$ \textbf{2-bromopentane} et \textbf{3-bromopentane} (mélange 50/50, même carbocation secondaire).
  \item Hydratation propène : \ce{CH3-CH=CH2 + H2O ->[H^+] CH3-CH(OH)-CH3} $\to$ \textbf{Propan-2-ol} (isopropanol).
  \item Acétylène + \ce{Cl2} :
  \ce{HC#CH + Cl2 -> CHCl=CHCl} (dichloroéthène, isomères Z/E possibles) \\
  \ce{CHCl=CHCl + Cl2 -> CHCl2-CHCl2} (1,1,2,2-tétrachloroéthane).
\end{enumerate}
\end{corrige}

\begin{corrige}[Exercice 3]
\begin{enumerate}
  \item $\ce{C2H2(g) + 2 Br2(aq) -> C2H2Br4(aq)}$ (tétrabromoéthane, incolore).
  \item $n_{\text{Br}_2,\text{init}} = C \times V = \SI{0,100}{mol.L^{-1}} \times \SI{50,0E-3}{L} = \SI{5,00E-3}{mol}$.
  Dosage excès : $\ce{Br2 + 2 S2O3^2- -> 2 Br- + S4O6^2-}$.
  $n_{\ce{S2O3^2-}} = \SI{0,100}{mol.L^{-1}} \times \SI{20,0E-3}{L} = \SI{2,00E-3}{mol}$.
  $n_{\text{Br}_2,\text{excès}} = \frac{n_{\ce{S2O3^2-}}}{2} = \SI{1,00E-3}{mol}$.
  $n_{\text{Br}_2,\text{réagi}} = \SI{5,00E-3}{mol} - \SI{1,00E-3}{mol} = \SI{4,00E-3}{mol}$.
  \item Stœchiométrie : 1 mol \ce{C2H2} $\leftrightarrow$ 2 mol \ce{Br2}.
  $n_{\ce{C2H2}} = \frac{n_{\text{Br}_2,\text{réagi}}}{2} = \SI{2,00E-3}{mol}$.
  \item Volume molaire $V_m = \frac{V}{n} = \frac{\SI{10,0E-3}{L}}{\SI{2,00E-3}{mol}} = \SI{5,00}{L.mol^{-1}}$.
  \textbf{Attention :} Ce résultat (\SI{5,0}{L.mol^{-1}}) est très éloigné de \SI{22,4}{L.mol^{-1}}. Cela signale une erreur expérimentale majeure (fuite, gaz non pur, température/pression non CNTP, réaction incomplète, dosage erroné). Aux CNTP (0°C, 1 atm), $V_m = \SI{22,4}{L.mol^{-1}}$. À température ambiante (25°C, 1 bar), $V_m \approx \SI{24,8}{L.mol^{-1}}$.
\end{enumerate}
\end{corrige}

\findecours

\end{document}
